\documentclass[10pt]{article}
\usepackage[utf8]{inputenc}
\usepackage[T1]{fontenc}
\usepackage{amsmath}
\usepackage{amsfonts}
\usepackage{amssymb}
\usepackage[version=4]{mhchem}
\usepackage{stmaryrd}
\usepackage{bbold}
\usepackage{graphicx}
\usepackage[export]{adjustbox}
\graphicspath{ {./images/} }
\usepackage{fvextra, csquotes}
\usepackage{caption}
\usepackage{MnSymbol}

%New command to display footnote whose markers will always be hidden
\let\svthefootnote\thefootnote
\newcommand\blfootnotetext[1]{%
  \let\thefootnote\relax\footnote{#1}%
  \addtocounter{footnote}{-1}%
  \let\thefootnote\svthefootnote%
}

%Overriding the \footnotetext command to hide the marker if its value is `0`
\let\svfootnotetext\footnotetext
\renewcommand\footnotetext[2][?]{%
  \if\relax#1\relax%
    \ifnum\value{footnote}=0\blfootnotetext{#2}\else\svfootnotetext{#2}\fi%
  \else%
    \if?#1\ifnum\value{footnote}=0\blfootnotetext{#2}\else\svfootnotetext{#2}\fi%
    \else\svfootnotetext[#1]{#2}\fi%
  \fi
}

\DeclareUnicodeCharacter{27FC}{\ifmmode\longmapsto\else{$\longmapsto$}\fi}
\DeclareUnicodeCharacter{2218}{\ifmmode\circ\else{$\circ$}\fi}
\DeclareUnicodeCharacter{25CB}{\ifmmode\medcircle\else{$\medcircle$}\fi}
\DeclareUnicodeCharacter{25C7}{\ifmmode\diamond\else{$\diamond$}\fi}
\DeclareUnicodeCharacter{27F9}{\ifmmode\Longrightarrow\else{$\Longrightarrow$}\fi}
\DeclareUnicodeCharacter{21D2}{\ifmmode\Rightarrow\else{$\Rightarrow$}\fi}
\DeclareUnicodeCharacter{00D7}{\ifmmode\times\else{$\times$}\fi}
\DeclareUnicodeCharacter{25A1}{\ifmmode\square\else{$\square$}\fi}

\begin{document}
\captionsetup{singlelinecheck=false}
\section*{Chapter II}
\section*{The Cohomology of Groups}
We take a respite from number theory and do some homological algebra. In an appendix to the chapter, we review the general theory of derived functors.

\section*{1 Cohomology}
The category of $G$-modules\\
Let $G$ be a group. A $G$-module is an abelian group $M$ together with a map

$$
(g, m) \mapsto g m: G \times M \rightarrow M
$$

such that, for all $g, g^{\prime} \in G, m, m^{\prime} \in M$,

\begin{itemize}
  \item[(a)] $g\left(m+m^{\prime}\right)=g m+g m^{\prime}$;
  \item[(b)] $\left(g g^{\prime}\right)(m)=g\left(g^{\prime} m\right), 1 m=m$.
\end{itemize}

Equivalently, a $G$-module is an abelian group $M$ together with a homomorphism of groups $G \rightarrow \operatorname{Aut}(M)$. We also say that $G$ acts on $M$, and we say that the action is trivial if $g m=m$ for all $g \in G$ and $m \in M$.

A homomorphism of $G$-modules (or a $G$-homomorphism) is a map $\alpha: M \rightarrow N$ such that

\begin{itemize}
  \item[(a)] $\alpha\left(m+m^{\prime}\right)=\alpha(m)+\alpha\left(m^{\prime}\right) \quad$ (i.e., $\alpha$ is a homomorphism of abelian groups);
  \item[(b)] $\alpha(g m)=g(\alpha(m))$ for all $g \in G, m \in M$.
\end{itemize}

We write $\operatorname{Hom}_{G}(M, N)$ for the set of $G$-homomorphisms $M \rightarrow N$.\\
REMARK 1.1 The group algebra $\mathbb{Z}[G]$ of $G$ is the free abelian group with basis the elements of $G$ and with the multiplication provided by the group law on $G$. Thus the elements of $\mathbb{Z}[G]$ are the finite sums

$$
\sum_{i} n_{i} g_{i}, \quad n_{i} \in \mathbb{Z}, \quad g_{i} \in G,
$$

and

$$
\left(\sum_{i} n_{i} g_{i}\right)\left(\sum_{j} n_{j}^{\prime} g_{j}^{\prime}\right)=\sum_{i, j} n_{i} n_{j}^{\prime}\left(g_{i} g_{j}^{\prime}\right) .
$$

A $G$-module structure on an abelian group $M$ extends uniquely to a $\mathbb{Z}[G]$-module structure, and a homomorphism of abelian groups is a homomorphism of $G$-modules if and only if\\
it is a homomorphism of $\mathbb{Z}[G]$-modules. Thus the category $\operatorname{Mod}_{G}$ of $G$-modules can be identified with the category of modules over the ring $\mathbb{Z}[G]$. In particular, $\operatorname{Mod}_{G}$ is an abelian category.

If $M$ and $N$ are $G$-modules, then the set $\operatorname{Hom}(M, N)$ of homomorphisms $\varphi: M \rightarrow N$ ( $M$ and $N$ regarded only as abelian groups) becomes a $G$-module with the structures


\begin{align*}
\left(\varphi+\varphi^{\prime}\right)(m) & =\varphi(m)+\varphi^{\prime}(m) \\
(g \varphi)(m) & =g\left(\varphi\left(g^{-1} m\right)\right) . \tag{11}
\end{align*}


\section*{Induced modules}
For an $H$-module $M$, we define $\operatorname{Ind}_{H}^{G}(M)$ to be the set of maps (not necessarily homomorphisms) $\varphi: G \rightarrow M$ such that $\varphi(h g)=h \varphi(g)$ for all $h \in H$. Then $\operatorname{Ind}_{H}^{G}(M)$ becomes a $G$-module with the operations


\begin{align*}
\left(\varphi+\varphi^{\prime}\right)(x) & =\varphi(x)+\varphi^{\prime}(x) \\
(g \varphi)(x) & =\varphi(x g) . \tag{12}
\end{align*}


A homomorphism $\alpha: M \rightarrow M^{\prime}$ of $H$-modules defines a homomorphism

$$
\varphi \mapsto \alpha \circ \varphi: \operatorname{Ind}_{H}^{G}(M) \rightarrow \operatorname{Ind}_{H}^{G}\left(M^{\prime}\right)
$$

of $G$-modules.\\
Lemma 1.2 (a) For every $G$-module $M$ and $H$-module $N$,

$$
\operatorname{Hom}_{G}\left(M, \operatorname{Ind}_{H}^{G}(N)\right) \simeq \operatorname{Hom}_{H}(M, N) .
$$

(b) The functor

$$
\operatorname{Ind}_{H}^{G}: \operatorname{Mod}_{H} \rightarrow \operatorname{Mod}_{G}
$$

is exact.\\
Proof. (a) Given a $G$-homomorphism $\alpha: M \rightarrow \operatorname{Ind}_{H}^{G}(N)$, we define a map $\beta: M \rightarrow N$ by the rule

$$
\beta(m)=\alpha(m)\left(1_{G}\right),
$$

where $1_{G}$ is the identity element in $G$. For every $g \in G$,

$$
\beta(g m) \stackrel{\text { def }}{=}(\alpha(g m))\left(1_{G}\right)=(g(\alpha(m)))\left(1_{G}\right) \stackrel{(12)}{=} \alpha(m)(g) .
$$

Because $\alpha(m) \in \operatorname{Ind}_{H}^{G}(N)$, when $g \in H, \alpha(m)(g)=g\left(\alpha(m)\left(1_{G}\right)\right)=g(\beta(m))$. Therefore, $\beta$ is an $H$-homomorphism $M \rightarrow N$.

Conversely, given an $H$-homomorphism $\beta: M \rightarrow N$, we define $\alpha$ to be the map $M \rightarrow$ $\operatorname{Ind}_{H}^{G}(N)$ such that $\alpha(m)(g)=\beta(g m)$. Then $\alpha$ is a $G$-homomorphism.

Since the maps $\alpha \mapsto \beta$ and $\beta \mapsto \alpha$ are inverse, both are isomorphisms.\\
(b) Given an exact sequence

$$
0 \rightarrow M \rightarrow N \rightarrow P \rightarrow 0,
$$

we have to prove that

$$
0 \rightarrow \operatorname{Ind}_{H}^{G} M \rightarrow \operatorname{Ind}_{H}^{G} N \rightarrow \operatorname{Ind}_{H}^{G} P \rightarrow 0
$$

is exact. This is obvious except at the last position. Let $S$ be a set of right coset representatives for $H$ in $G$, so that $G=\bigcup_{s \in S} H s$, and let $\varphi \in \operatorname{Ind}_{H}^{G}(P)$. For each $s \in S$, choose an $n(s) \in N$ mapping to $\varphi(s)$ in $P$, and define $\tilde{\varphi}(h s)=h \cdot n(s)$. Then $\tilde{\varphi} \in \operatorname{Ind}_{H}^{G}(N)$ and maps to $\varphi$. \(\square\)

Let $\phi$ be the map $\operatorname{Ind}_{H}^{G}(N) \rightarrow N, \varphi \mapsto \varphi\left(1_{G}\right)$; this is an $H$-homomorphism, and the lemma shows that $\left(\operatorname{Ind}_{H}^{G}(N), \phi\right)$ has the following universal property:\\
\includegraphics[max width=\textwidth, alt={}, center]{e7cfc28e-5737-4710-b4d2-10a47512e540-03_256_350_849_1246}\\
for any $H$-homomorphism $\beta: M \rightarrow N$ from a $G$-module $M$ to $N$, there exists a unique $G$ -homomorphism $\alpha: M \rightarrow \operatorname{Ind}_{H}^{G}(N)$ such that $\phi \circ \alpha=\beta$.

When $H=\{1\}$, an $H$-module is just an abelian group. In this case, we drop the $H$ from the notation $\operatorname{Ind}_{H}^{G}$. Thus

$$
\begin{aligned}
\operatorname{Ind}^{G}\left(M_{0}\right) & =\left\{\varphi: G \rightarrow M_{0}\right\} & & \text { ( maps, not necessarily homomorphisms) } \\
& =\operatorname{Hom}\left(\mathbb{Z}[G], M_{0}\right) & & \text { ( homomorphisms of abelian groups) } .
\end{aligned}
$$

A $G$-module is said to be induced if it isomorphic to $\operatorname{Ind}^{G}\left(M_{0}\right)$ for some abelian group $M_{0}$.

Remark 1.3 Let $G$ be a finite group.

\begin{itemize}
  \item[(a)] A $G$-module $M$ is induced if and only if there exists an abelian group $M_{0} \subset M$ such that
$$
M=\bigoplus_{g \in G} g M_{0} \quad \text { (direct sum of abelian groups) },
$$
in which case there is an isomorphism of $G$-modules
$$
\varphi \mapsto \sum_{g \in G} g \otimes \varphi\left(g^{-1}\right): \operatorname{Ind}^{G}\left(M_{0}\right) \rightarrow \mathbb{Z}[G] \otimes_{\mathbb{Z}} M_{0} .
$$
Here $\mathbb{Z}[G] \otimes_{\mathbb{Z}} M_{0}$ is endowed with the $G$-structure such that
$$
g(z \otimes m)=g z \otimes m .
$$
  \item[(b)] Let $H$ be a subgroup of $G$. An induced $G$-module is also induced when considered as an $H$-module: let $S$ be a set of right coset representatives for $H$ in $G$, so that $G=$ $\bigcup_{s \in S} H s$; if $M=\bigoplus_{g \in G} g M_{0}$, then $M=\bigoplus_{h \in H} h M_{1}$ with $M_{1}=\bigoplus_{s \in S} s M_{0}$.
  \item[(c)] Let $M$ be a $G$-module, and let $M_{0}$ be $M$ regarded as an abelian group. Then
$$
\pi: \operatorname{Ind}^{G}\left(M_{0}\right) \rightarrow M, \quad \varphi \mapsto \sum_{g \in G} g \varphi\left(g^{-1}\right)
$$
is a surjective homomorphism of $G$-modules. It corresponds to the map
$$
\mathbb{Z}[G] \otimes M_{0} \rightarrow M, \quad\left(\sum n_{g} g\right) \otimes m \mapsto \sum n_{g} g m .
$$
\end{itemize}

Remark 1.4 Let $M$ and $N$ be $G$-modules. Then the rule

$$
g(m \otimes n)=g m \otimes g n
$$

defines a $G$-module structure on $M \otimes_{\mathbb{Z}} N$. Let $M_{0}$ be $M$ regarded as an abelian group. Then $\mathbb{Z}[G] \otimes_{\mathbb{Z}} M=\mathbb{Z}[G] \otimes_{\mathbb{Z}} M_{0}$ as abelian groups, but their $G$-module structures do not correspond. However, one checks easily that the map $\mathbb{Z}[G] \otimes_{\mathbb{Z}} M_{0} \rightarrow \mathbb{Z}[G] \otimes_{\mathbb{Z}} M$ sending $g \otimes m$ to $g \otimes g m$ is an isomorphism of $G$-modules.

\section*{Injective $G$-modules}
A $G$-module $I$ is said to be injective if every $G$-homomorphism from a submodule of a $G$-module extends to the whole module, or, equivalently, if $\operatorname{Hom}_{G}(\cdot, I)$ is an exact functor.

Proposition 1.5 The category $\operatorname{Mod}_{G}$ has enough injectives, i.e., every $G$-module $M$ can be embedded into an injective $G$-module, $M \hookrightarrow I$.

Proof. When $G=\{1\}$, so that $\operatorname{Mod}_{G}$ is the category of abelian groups, this is proved in the Appendix (A.4). Now let $M$ be a $G$-module, and let $M_{0}$ be $M$ regarded as an abelian group. We can embed $M_{0}$ into an injective abelian group, say, $M_{0} \hookrightarrow I$. On applying the functor $\operatorname{Ind}^{G}$, we obtain an inclusion $\operatorname{Ind}^{G}\left(M_{0}\right) \hookrightarrow \operatorname{Ind}^{G}(I)$ of $G$-modules. There is an inclusion of $G$-modules $M \hookrightarrow \operatorname{Ind}^{G}\left(M_{0}\right)$ sending $m$ to the function $g \mapsto g m$. On composing these maps, we obtain an injective homomorphism $M \hookrightarrow \operatorname{Ind}^{G}(I)$, and so it remains to show that $\operatorname{Ind}^{G}(I)$ is an injective $G$-module, but the natural isomorphism

$$
\operatorname{Hom}_{\operatorname{Mod}_{G}}\left(M, \operatorname{Ind}^{G}(I)\right) \simeq \operatorname{Hom}_{\mathrm{Ab}}(M, I)
$$

shows that $\operatorname{Hom}_{\operatorname{Mod}_{G}}\left(-, \operatorname{Ind}^{G}(I)\right)$ is exact. \(\square\)

\section*{Definition of the cohomology groups}
For a $G$-module $M$, define

$$
M^{G}=\{m \in M \mid g m=m \text { all } g \in G\} .
$$

The functor

$$
M \mapsto M^{G}: \operatorname{Mod}_{G} \rightarrow \mathrm{Ab}
$$

is left exact, i.e., if

$$
0 \rightarrow M^{\prime} \rightarrow M \rightarrow M^{\prime \prime} \rightarrow 0
$$

is exact, then

$$
0 \rightarrow M^{\prime G} \rightarrow M^{G} \rightarrow M^{\prime \prime G}
$$

is exact. This can be checked directly, or by observing that $(-)^{G}$ is isomorphic to the left exact functor $\operatorname{Hom}_{G}(\mathbb{Z},-)$.

Since the category of $G$-modules has enough injectives, we can apply the theory of derived functors (see the appendix to this chapter) to this situation.

Let $M$ be a $G$-module, and choose an injective resolution

$$
0 \rightarrow M \rightarrow I^{0} \xrightarrow{d^{0}} I^{1} \xrightarrow{d^{1}} I^{2} \xrightarrow{d^{2}} \cdots
$$

of $M$. The complex

$$
0 \xrightarrow{d^{-1}}\left(I^{0}\right)^{G} \xrightarrow{d^{0}}\left(I^{1}\right)^{G} \rightarrow \cdots \xrightarrow{d^{r-1}}\left(I^{r}\right)^{G} \xrightarrow{d^{r}}\left(I^{r+1}\right)^{G} \rightarrow \cdots
$$

need no longer be exact, and we define the $r$ th cohomology group of $G$ with coefficients in $M$ to be

$$
H^{r}(G, M)=\frac{\operatorname{Ker}\left(d^{r}\right)}{\operatorname{Im}\left(d^{r-1}\right)}
$$

These groups have the following basic properties.\\
1.6 The zeroth group $H^{0}(G, M)=M^{G}$, because

$$
0 \rightarrow M^{G} \rightarrow I^{0 G} \xrightarrow{d^{0}} I^{1 G}
$$

is exact, and $H^{0}(G, M) \stackrel{\text { def }}{=} \frac{\operatorname{Ker}\left(d^{0}\right)}{\operatorname{Im}\left(d^{-1}\right)}=\operatorname{Ker}\left(d^{0}\right)$.\\
1.7 If $I$ is an injective $G$-module, then $H^{r}(G, I)=0$ for all $r>0$, because

$$
0 \rightarrow I \rightarrow I \rightarrow 0 \rightarrow \cdots
$$

is an injective resolution of $I$.\\
1.8 Let $M \rightarrow I^{\bullet}$ and $N \rightarrow J^{\bullet}$ be injective resolutions of the $G$-modules $M$ and $N$. Any homomorphism $\alpha: M \rightarrow N$ of $G$-modules extends to a map of complexes\\
\includegraphics[max width=\textwidth, alt={}, center]{e7cfc28e-5737-4710-b4d2-10a47512e540-05_173_243_1378_915}\\
and the homomorphisms

$$
H^{r}\left(\alpha^{\bullet}\right): H^{r}\left(I^{\bullet}\right) \rightarrow H^{r}\left(J^{\bullet}\right)
$$

are independent of the choice of $\alpha^{\bullet}$. On applying this statement to the identity map id: $M \rightarrow$ $M$, we find that the groups $H^{r}(G, M)$ are well-defined up to a well-defined isomorphism. The statement for a general $\alpha$ then shows that $M \mapsto H^{r}(G, M)$ is a functor from the category of $G$-modules to the category of abelian groups.\\
1.9 A short exact sequence

$$
0 \rightarrow M^{\prime} \rightarrow M \rightarrow M^{\prime \prime} \rightarrow 0
$$

of $G$-modules gives rise to a long exact sequence

$$
0 \rightarrow H^{0}\left(G, M^{\prime}\right) \rightarrow \cdots \rightarrow H^{r}(G, M) \rightarrow H^{r}\left(G, M^{\prime \prime}\right) \xrightarrow{\delta^{r}} H^{r+1}\left(G, M^{\prime}\right) \rightarrow \cdots
$$

Moreover, the association\\
short exact sequence ⟼ long exact sequence\\
is functorial, i.e., a morphism of short exact sequences induces a morphism of long exact sequences (see A.11).

REMARK 1.10 The family of functors $\left(H^{r}(G, \cdot)\right)_{r \geq 0}$ and coboundary maps $\delta^{r}$ are uniquely determined by the properties (1.6, 1.7, 1.9).

\section*{Shapiro's lemma}
Let $M$ be a $G$-module, and regard $\mathbb{Z}$ as a $G$-module with the trivial action: $g m=m$ for all $g \in G, m \in \mathbb{Z}$. A homomorphism $\alpha: \mathbb{Z} \rightarrow M$ is uniquely determined by $\alpha(1)$, and $m \in M$ is the image of 1 under a $G$-homomorphism $\mathbb{Z} \rightarrow M$ if and only if it is fixed by $G$. Therefore


\begin{equation*}
\operatorname{Hom}_{G}(\mathbb{Z}, M) \simeq M^{G} . \tag{13}
\end{equation*}


Proposition 1.11 (Shapiro's Lemma) Let $H$ be a subgroup of $G$. For every $H$ -module $N$, there is a canonical isomorphism

$$
H^{r}\left(G, \operatorname{Ind}_{H}^{G}(N)\right) \rightarrow H^{r}(H, N)
$$

for all $r \geq 0$.\\
Proof. For $r=0$, the isomorphism is the composite


\begin{equation*}
N^{H} \stackrel{(13)}{\simeq} \operatorname{Hom}_{H}(\mathbb{Z}, N) \stackrel{(1.2)}{\simeq} \operatorname{Hom}_{G}\left(\mathbb{Z}, \operatorname{Ind}_{H}^{G}(N)\right) \stackrel{(13)}{\simeq} \operatorname{Ind}_{H}^{G}(N)^{G} . \tag{14}
\end{equation*}


Now choose an injective resolution $N \rightarrow I^{\bullet}$ of $N$. On applying the functor $\operatorname{Ind}_{H}^{G}$, we obtain an injective resolution $\operatorname{Ind}_{H}^{G}(N) \rightarrow \operatorname{Ind}_{H}^{G}\left(I^{\bullet}\right)$ of the $G$-module $\operatorname{Ind}_{H}^{G}(N)$ because the functor $\operatorname{Ind}_{H}^{G}$ is exact (1.2) and preserves injectives (proof of 1.5). Hence

$$
H^{r}\left(G, \operatorname{Ind}_{H}^{G}(N)\right)=H^{r}\left(\left(\operatorname{Ind}_{H}^{G}\left(I^{\bullet}\right)\right)^{G}\right) \stackrel{(14)}{\simeq} H^{r}\left(I^{\bullet H}\right)=H^{r}(H, N) .
$$ \(\square\)

Corollary 1.12 If $M$ is an induced $G$-module, then $H^{r}(G, M)=0$ for $r>0$.\\
Proof. If $M=\operatorname{Ind}^{G}\left(M_{0}\right)$, then

$$
H^{r}(G, M) \simeq H^{r}\left(\{1\}, M_{0}\right)=0 \quad \text { for } r>0 .
$$ \(\square\)

Remark 1.13 Consider an exact sequence

$$
0 \rightarrow M \rightarrow J \rightarrow N \rightarrow 0
$$

of $G$-modules. If $H^{r}(G, J)=0$ for all $r>0$, then the cohomology sequence becomes the exact sequence

$$
0 \rightarrow M^{G} \rightarrow J^{G} \rightarrow N^{G} \rightarrow H^{1}(G, M) \rightarrow 0
$$

and the collection of isomorphisms

$$
H^{r}(G, N) \xrightarrow{\approx} H^{r+1}(G, M), \quad r \geq 1 .
$$

For example, let $M$ be a $G$-module, and let $M_{*}$ be the induced module $\operatorname{Ind}^{G}\left(M_{0}\right)$, where $M_{0}$ is $M$ regarded as an abelian group. As we noted in the proof of (1.5), $M$ can be identified with the $G$-submodule of $M_{*}$ consisting of maps of the form $g \mapsto g m, m \in M$. Let $M_{\dagger}=M_{*} / M$. On applying the above remark to the sequence


\begin{equation*}
0 \rightarrow M \rightarrow M_{*} \rightarrow M_{\dagger} \rightarrow 0 \tag{15}
\end{equation*}


we find that

$$
H^{r}\left(G, M_{\dagger}\right) \simeq H^{r+1}(G, M), \quad \text { all } r \geq 1 .
$$

More generally, an exact sequence

$$
0 \rightarrow M \rightarrow J^{1} \rightarrow \cdots \rightarrow J^{s} \rightarrow N \rightarrow 0
$$

such that $H^{r}\left(G, J^{i}\right)=0$ for all $r, i>0$, defines isomorphisms

$$
H^{r}(G, N) \xrightarrow{\simeq} H^{r+s}(G, M), \quad \text { all } r \geq 1 .
$$

To prove this, break the sequence up into short exact sequences

$$
\begin{aligned}
& 0 \rightarrow M \rightarrow J^{1} \rightarrow N^{1} \rightarrow 0 \\
& 0 \rightarrow N^{1} \rightarrow J^{2} \rightarrow N^{2} \rightarrow 0 \\
& \cdots \\
& 0 \rightarrow N^{s-1} \rightarrow J^{s} \rightarrow N \rightarrow 0
\end{aligned}
$$

and note that we have isomorphisms

$$
H^{r}(G, N) \simeq H^{r+1}\left(G, N^{s-1}\right) \simeq H^{r+2}\left(G, N^{s-2}\right) \simeq \cdots .
$$

Remark 1.14 Let

$$
0 \rightarrow M \xrightarrow{\varepsilon} J^{0} \xrightarrow{d^{0}} J^{1} \xrightarrow{d^{1}} J^{2} \rightarrow \cdots
$$

be an exact sequence such that $H^{s}\left(G, J^{r}\right)=0$ for all $s>0$ and all $r$. Then

$$
H^{r}(G, M)=H^{r}\left(J^{\bullet G}\right) .
$$

This remark applies to every resolution of $M$ by induced modules.

\section*{Description of the cohomology groups by means of cochains}
Let $P_{r}, r \geq 0$, be the free $\mathbb{Z}$-module with basis the $(r+1)$-tuples $\left(g_{0}, \ldots, g_{r}\right)$ of elements of $G$, endowed the action of $G$ such that

$$
g\left(g_{0}, \ldots, g_{r}\right)=\left(g g_{0}, \ldots, g g_{r}\right) .
$$

Note that $P_{r}$ is also free as a $\mathbb{Z}[G]$-module, with basis $\left\{\left(1, g_{1}, \ldots, g_{r}\right) \mid g_{i} \in G\right\}$. Define a homomorphism $d_{r}: P_{r} \rightarrow P_{r-1}$ by the rule

$$
d_{r}\left(g_{0}, \ldots, g_{r}\right)=\sum_{i=0}^{r}(-1)^{i}\left(g_{0}, \ldots, \hat{g}_{i}, \ldots, g_{r}\right),
$$

where the symbol \^{} means that v . is omitted. Let $P_{\bullet}$ be

$$
\cdots \rightarrow P_{r} \xrightarrow{d_{r}} P_{r-1} \rightarrow \cdots \rightarrow P_{0}
$$

One checks easily that $d_{r-1} \circ d_{r}=0$, and so this is a complex. Let $\varepsilon$ be the map $P_{0} \rightarrow \mathbb{Z}$ sending each basis element to 1 .

Lemma 1.15 The complex $P_{\bullet} \stackrel{\varepsilon}{\rightarrow} \mathbb{Z} \rightarrow 0$ is exact.

Proof. Choose an element $o \in G$, and define $k_{r}: P_{r} \rightarrow P_{r+1}$ by

$$
k_{r}\left(g_{0}, \ldots, g_{r}\right)=\left(o, g_{0}, \ldots, g_{r}\right) .
$$

One checks easily that $d_{r+1} \circ k_{r}+k_{r-1} \circ d_{r}=1$. Hence, if $d_{r}(x)=0$, then $x=$ $d_{r+1}\left(k_{r}(x)\right)$. \(\square\)

Proposition 1.16 For every $G$-module $M$,

$$
H^{r}(G, M) \simeq H^{r}\left(\operatorname{Hom}_{G}\left(P_{\bullet}, M\right)\right) .
$$

Proof. This follows from the fact that $P_{\bullet} \rightarrow \mathbb{Z}$ is a projective resolution of $\mathbb{Z}$ - see Example A.14. \(\square\)

An element of $\operatorname{Hom}\left(P_{r}, M\right)$ can be identified with a function $\varphi: G^{r+1} \rightarrow M$, and $\varphi$ is fixed by $G$ if and only if

$$
\varphi\left(g g_{0}, \ldots, g g_{r}\right)=g\left(\varphi\left(g_{0}, \ldots, g_{r}\right)\right) \text { all } g, g_{0}, \ldots, g_{r} \in G .
$$

Thus $\operatorname{Hom}_{G}\left(P_{r}, M\right)$ can be identified with the set $\tilde{C}^{r}(G, M)$ of $\varphi$ 's satisfying this condition. Such $\varphi$ are called homogeneous $r$-cochains of $G$ with values in $M$. The boundary map $\tilde{d}^{r}: \tilde{C}^{r}(G, M) \rightarrow \tilde{C}^{r+1}(G, M)$ induced by $d_{r+1}$ is

$$
\left(\tilde{d}^{r} \varphi\right)\left(g_{0}, \ldots, g_{r+1}\right)=\sum(-1)^{i} \varphi\left(g_{0}, \ldots, \hat{g}_{i}, \ldots, g_{r+1}\right) .
$$

Proposition 1.16 says that

$$
H^{r}(G, M) \simeq \frac{\operatorname{Ker}\left(\tilde{d}^{r}\right)}{\operatorname{Im}\left(\tilde{d}^{r-1}\right)} .
$$

A homogenous cochain $\varphi: G^{r+1} \rightarrow M$ is determined by its values on the elements $\left(1, g_{1}, g_{1} g_{2}, \ldots, g_{1} \ldots g_{r}\right)$. We are therefore led to introduce the group $C^{r}(G, M)$ of inhomogeneous $r$-cochains of $G$ with values in $M$ consisting of all maps $\varphi: G^{r} \rightarrow M$. We set $G^{0}=\{1\}$, so that $C^{0}(G, M)=M$. Define

$$
d^{r}: C^{r}(G, M) \rightarrow C^{r+1}(G, M),
$$

by $\left(d^{r} \varphi\right)\left(g_{1}, \cdots, g_{r+1}\right)=$

$$
g_{1} \varphi\left(g_{2}, \ldots, g_{r+1}\right)+\sum_{j=1}^{r}(-1)^{j} \varphi\left(g_{1}, \ldots, g_{j} g_{j+1}, \ldots, g_{r+1}\right)+(-1)^{r+1} \varphi\left(g_{1}, \ldots, g_{r}\right) .
$$

Define

$$
Z^{r}(G, M)=\operatorname{Ker}\left(d^{r}\right) \quad(\text { group of } r \text {-cocycles })
$$

and

$$
B^{r}(G, M)=\operatorname{Im}\left(d^{r-1}\right) \quad \text { (group of } r \text {-coboundaries). }
$$

Proposition 1.17 The sequence of maps

$$
C^{0}(G, M) \xrightarrow{d^{0}} C^{1}(G, M) \xrightarrow{d^{1}} \cdots \xrightarrow{d^{r-1}} C^{r}(G, M) \xrightarrow{d^{r}} C^{r+1}(G, M) \rightarrow \cdots
$$

is a complex, i.e., $d^{r} \circ d^{r-1}=0$, and there is a canonical isomorphism

$$
H^{r}(G, M) \simeq \frac{Z^{r}(G, M)}{B^{r}(G, M)} .
$$

Proof. For $\varphi \in \tilde{C}^{r}(G, M)$, define

$$
\varphi^{\prime}\left(g_{1}, \ldots, g_{r}\right)=\varphi\left(1, g_{1}, g_{1} g_{2}, \ldots, g_{1} \cdots g_{r}\right) .
$$

Then $\varphi \mapsto \varphi^{\prime}$ is a bijection $\tilde{C}^{r}(G, M) \rightarrow C^{r}(G, M)$ transforming the boundary maps in $\tilde{C}^{\bullet}(G, M)$ into the boundary maps in $C^{\bullet}(G, M)$. \(\square\)

EXAMPLE 1.18 (a) A map $\varphi$ : $G \rightarrow M$ is a crossed homomorphism if

$$
\varphi(\sigma \tau)=\sigma \varphi(\tau)+\varphi(\sigma), \text { all } \sigma, \tau \in G .
$$

Notice that the condition implies that, for $e$ the identity element of $G$,

$$
\varphi(e)=\varphi(e e)=e \varphi(e)+\varphi(e)=2 \varphi(e)
$$

and so $\varphi(e)=0$. For every $m \in M$, the map $\sigma \mapsto \sigma m-m$ is a crossed homomorphism—called a principal crossed homomorphism. According to the proposition

$$
H^{1}(G, M)=\frac{\{\text { crossed homomorphisms } G \rightarrow M\}}{\{\text { principal crossed homomorphisms }\}} .
$$

If the action of $G$ on $M$ is trivial, then a crossed homomorphism is a homomorphism, and the principal crossed homomorphisms are zero. Thus, in this case,


\begin{equation*}
H^{1}(G, M) \simeq \operatorname{Hom}(G, M) . \tag{16}
\end{equation*}


(b) Let $M$ be an abelian group (with the law of composition written as multiplication). An extension of $G$ by $M$ is an exact sequence of groups

$$
1 \rightarrow M \rightarrow E \xrightarrow{\pi} G \rightarrow 1 .
$$

We set

$$
\sigma m=s(\sigma) \cdot m \cdot s(\sigma)^{-1}, \quad \sigma \in G, m \in M,
$$

where $s(\sigma)$ is any element of $E$ mapping to $\sigma$. Because $M$ is commutative, $\sigma m$ depends only on $\sigma$, and this defines an action of $G$ on $M$. Note that

$$
s(\sigma) \cdot m=\sigma m \cdot s(\sigma), \quad \text { all } \sigma \in G, \quad m \in M .
$$

Now choose a section $s$ to $\pi$, i.e., a map (not necessarily a homomorphism) $s: G \rightarrow E$ such that $\pi \circ s=\mathrm{id}$. Then $s(\sigma) s\left(\sigma^{\prime}\right)$ and $s\left(\sigma \sigma^{\prime}\right)$ both map to $\sigma \sigma^{\prime} \in G$, and so they differ by an element $\varphi\left(\sigma, \sigma^{\prime}\right) \in M$ :

$$
s(\sigma) s\left(\sigma^{\prime}\right)=\varphi\left(\sigma, \sigma^{\prime}\right) \cdot s\left(\sigma \sigma^{\prime}\right) .
$$

From

$$
s(\sigma)\left(s\left(\sigma^{\prime}\right) s\left(\sigma^{\prime \prime}\right)\right)=\left(s(\sigma) s\left(\sigma^{\prime}\right)\right) s\left(\sigma^{\prime \prime}\right)
$$

we deduce that

$$
\sigma \varphi\left(\sigma^{\prime}, \sigma^{\prime \prime}\right) \cdot \varphi\left(\sigma, \sigma^{\prime} \sigma^{\prime \prime}\right)=\varphi\left(\sigma, \sigma^{\prime}\right) \cdot \varphi\left(\sigma \sigma^{\prime}, \sigma^{\prime \prime}\right),
$$

i.e., that $\varphi \in Z^{2}(G, M)$. If $s$ is replaced by a different section, $\varphi$ is replaced by a cohomologous cocycle, and so the class of $\varphi$ in $H^{2}(G, M)$ is independent of the choice of $s$. Every such $\varphi$ arises from an extension. In this way, $H^{2}(G, M)$ classifies the isomorphism classes of extensions of $G$ by $M$ with a fixed action of $G$ on $M$.

Remark 1.19 Let $G$ be a group, and let $M$ be a $G$-module. For $m \in M$, let $\varphi_{m}: G \rightarrow M$ be the constant map $\sigma \mapsto m$. Then

$$
\left(d^{1} \varphi_{m}\right)(\sigma, \tau)=\sigma m-m+m=\sigma m .
$$

In particular, $\left(d^{1} \varphi_{m}\right)(1,1)=m$. Therefore, every class in $H^{2}(G, M)$ is represented by a 2-cocycle $\varphi$ with $\varphi(1,1)=0$. Such a 2-cocycle is said to be normalized.

EXAMPLE 1.20 Let $\varphi: G \rightarrow M$ be a crossed homomorphism. For every $\sigma \in G$,

$$
\begin{aligned}
\varphi\left(\sigma^{2}\right) & =\sigma \varphi(\sigma)+\varphi(\sigma) \\
\varphi\left(\sigma^{3}\right) & =\varphi\left(\sigma \cdot \sigma^{2}\right)=\sigma^{2} \varphi(\sigma)+\sigma \varphi(\sigma)+\varphi(\sigma) \\
& \cdots \\
\varphi\left(\sigma^{f}\right) & =\sigma^{f-1} \varphi(\sigma)+\cdots+\sigma \varphi(\sigma)+\varphi(\sigma) .
\end{aligned}
$$

Thus, if $G$ is a cyclic group of order $f$ generated by $\sigma$, then a crossed homomorphism $\varphi$ is determined by its value, $m$ say, on $\sigma$, and $m$ satisfies the equation


\begin{equation*}
\sigma^{f-1} m+\cdots+\sigma m+m=0 . \tag{17}
\end{equation*}


Conversely, if $m \in M$ satisfies (17), then the formulas $\varphi\left(\sigma^{i}\right)=\sigma^{i-1} m+\cdots+\sigma m+m$ define a crossed homomorphism $\varphi: G \rightarrow M$, and

$$
\varphi \text { is principal } \Longleftrightarrow m=\sigma x-x \text { for some } x \in M .
$$

Define maps

$$
\begin{array}{ll}
\operatorname{Nm}_{G}: M \rightarrow M, & m \mapsto \sum_{\sigma \in G} \sigma m \\
\sigma-1: M \rightarrow M, & m \mapsto \sigma m-m .
\end{array}
$$

Then, for a cyclic group $G$ of finite order, we have shown that the choice of a generator $\sigma$ for $G$ determines an isomorphism

$$
\varphi \mapsto \varphi(\sigma): H^{1}(G, M) \rightarrow \frac{\operatorname{Ker}\left(\operatorname{Nm}_{G}\right)}{(\sigma-1) M} .
$$

Remark 1.21 Let

$$
0 \rightarrow M \rightarrow N \rightarrow P \rightarrow 0
$$

be an exact sequence of $G$-modules. The boundary map

$$
\delta^{r}: H^{r}(G, P) \rightarrow H^{r+1}(G, M)
$$

has the following description: let $\gamma \in H^{r}(G, P)$ be represented by the $r$-cocycle $\varphi: G^{r} \rightarrow$ $P$; because $N$ maps onto $P$, there exists an $r$-cochain $\tilde{\varphi}: G^{r} \rightarrow N$ lifting $\varphi$; because $d \varphi=0, d \tilde{\varphi}$ takes values in $M$ - it is the cocycle representing $\delta^{r} \gamma$.

\section*{The cohomology of $L$ and $L^{\times}$}
Let $L$ be a finite Galois extension of the field $K$, and let $G=\operatorname{Gal}(L / K)$. Then both $L$ (regarded as a group under addition) and $L^{\times}$are $G$-modules.

Proposition 1.22 Let $L / K$ be a finite Galois extension with Galois group $G$. Then $H^{1}\left(G, L^{\times}\right)=0$.

Proof. Let $\varphi: G \rightarrow L^{\times}$be a crossed homomorphism. In multiplicative notation, this means that

$$
\varphi(\sigma \tau)=\sigma \varphi(\tau) \cdot \varphi(\sigma), \quad \sigma, \tau \in G,
$$

and we have to find a $c \in L^{\times}$such that $\varphi(\sigma)=\sigma c / c$. For $a \in L^{\times}$, let

$$
b=\sum_{\sigma \in G} \varphi(\sigma) \cdot \sigma a .
$$

Suppose $b \neq 0$. Then

$$
\tau b=\sum_{\sigma} \tau \varphi(\sigma) \cdot \tau \sigma a=\sum_{\sigma} \varphi(\tau)^{-1} \varphi(\tau \sigma) \tau \sigma a=\varphi(\tau)^{-1} b .
$$

Hence

$$
\varphi(\tau)=b / \tau b=\tau\left(b^{-1}\right) / b^{-1},
$$

which shows that $\varphi$ is principal.\\
It remains to show that there exists an $a$ for which $b \neq 0$. Recall (FT, 5.14, Dedekind's theorem on the independence of characters):

Let $L$ be a field and $H$ a group; then every finite set $\left\{f_{i}\right\}$ of distinct homomorphisms $H \rightarrow L^{\times}$is linearly independent over $L$, i.e.,

$$
\sum a_{i} f_{i}(\alpha)=0 \quad \text { all } \alpha \in H \Rightarrow a_{1}=a_{2}=\cdots=a_{n}=0 .
$$

When we apply this to the homomorphisms $\sigma: L^{\times} \rightarrow L^{\times}, \sigma \in G$, we find that $\sum_{\sigma} \varphi(\sigma) \sigma$ is not the zero map, and so there exists an $a$ for which $b \neq 0$. \(\square\)

Corollary 1.23 Let $L / K$ be a cyclic extension, and let $\sigma$ generate $\operatorname{Gal}(L / K)$. If $\operatorname{Nm}_{L / K} a=1$, then $a$ is of the form $\frac{\sigma b}{b}$.

Proof. We have $0=H^{1}\left(G, L^{\times}\right)=\operatorname{Ker}\left(\operatorname{Nm}_{G}\right) /(\sigma-1) L^{\times}$. \(\square\)

Corollary 1.23 occurs as Satz 90 of Hilbert's book, Die Theorie der algebraischen Zahlkörper, 1894/95, and Theorem 1.22 is Emmy Noether's generalization. Both are usually referred to as Hilbert's Theorem 90.

Proposition 1.24 Let $L / K$ be a finite Galois extension with Galois group $G$. Then $H^{r}(G, L)=0$ for all $r>0$.

Proof. Recall (FT, 5.18, Normal Basis Theorem):\\
There exists an $\alpha \in L$ such that $\{\sigma \alpha \mid \sigma \in G\}$ is a basis for $L$ as a $K$-vector space. (A basis of this form is said to be normal.)\\
A normal basis $(\sigma \alpha)_{\sigma \in G}$ defines an isomorphism of $G$-modules

$$
\sum_{\sigma \in G} a_{\sigma} \sigma \mapsto \sum_{\sigma \in G} a_{\sigma} \sigma \alpha: K[G] \rightarrow L .
$$

But $K[G]=\operatorname{Ind}_{\{1\}}^{G} K$, and so $H^{r}(G, L) \simeq H^{r}(\{1\}, K)=0$ for $r>0$ (Shapiro's lemma, 1.11). \(\square\)

The cohomology of products\\
A product $M=\prod_{i} M_{i}$ of $G$-modules becomes a $G$-module under the diagonal action:

$$
\sigma\left(\ldots, m_{i}, \ldots\right)=\left(\ldots, \sigma m_{i}, \ldots\right) .
$$

Proposition 1.25 For all families of $G$-modules $\left(M_{i}\right)_{i}$,

$$
H^{r}\left(G, \prod_{i} M_{i}\right) \simeq \prod_{i} H^{r}\left(G, M_{i}\right) .
$$

Proof. A product of exact sequences of abelian groups is again exact. From this it follows that a product $I=\prod I_{i}$ of injective $G$-modules is again injective, because

$$
\operatorname{Hom}_{G}(\cdot, I) \simeq \prod_{i} \operatorname{Hom}_{G}\left(\cdot, I_{i}\right)
$$

is exact. Let $M_{i} \rightarrow I_{i}^{\bullet}$ be an injective resolution of $M_{i}$. Then $\prod M_{i} \rightarrow \prod I_{i}^{\bullet}$ is an injective resolution of $\prod M_{i}$, and

$$
H^{r}\left(G, \prod_{i} M_{i}\right) \simeq H^{r}\left(\left(\prod_{i} I_{i}^{\bullet}\right)^{G}\right) \simeq H^{r}\left(\prod_{i}\left(I_{i}^{\bullet G}\right)\right) \simeq \prod_{i} H^{r}\left(I_{i}^{\bullet G}\right) \simeq \prod_{i} H^{r}\left(G, M_{i}\right) .
$$ \(\square\)

In particular, for any $G$-modules $M, N$,

$$
H^{r}(G, M \oplus N)=H^{r}(G, M) \oplus H^{r}(G, N) .
$$

This can be proved more directly by noting that to say that a module $P$ is a direct sum of modules $M$ and $N$ means that certain maps, and relations between the maps, exist (see p. 90). These maps and relations persist when we apply the additive functor $H^{r}(G, \cdot)$.

Remark 1.26 The formation of inverse limits of arbitrary abelian groups is not exact. Therefore, in general, one can not expect cohomology to commute with inverse limits.

Functorial properties of the cohomology groups\\
Let $M$ and $M^{\prime}$ respectively be $G$ and $G^{\prime}$ modules. Homomorphisms

$$
\alpha: G^{\prime} \rightarrow G, \quad \beta: M \rightarrow M^{\prime}
$$

are said to be compatible if

$$
\beta(\alpha(g) m)=g(\beta(m)) .
$$

Then $(\alpha, \beta)$ defines a homomorphism of complexes

$$
C^{\bullet}(G, M) \rightarrow C^{\bullet}\left(G^{\prime}, M^{\prime}\right), \quad \varphi \mapsto \beta \circ \varphi \circ \alpha^{r},
$$

and hence homomorphisms

$$
H^{r}(G, M) \rightarrow H^{r}\left(G^{\prime}, M^{\prime}\right) .
$$

Example 1.27 (a) Let $H$ be a subgroup of $G$. For every $H$-module $M$, the map

$$
\varphi \mapsto \varphi\left(1_{G}\right): \operatorname{Ind}_{H}^{G}(M) \rightarrow M
$$

is compatible with the inclusion $H \hookrightarrow G$, and the induced homomorphism

$$
H^{r}\left(G, \operatorname{Ind}_{H}^{G}(M)\right) \rightarrow H^{r}(H, M)
$$

is the isomorphism in Shapiro's Lemma 1.11. To show that the two isomorphisms agree, first check it for $r=0$, and then use dimension-shifting (II, 1.28) to get it for all $r$.\\
(b) Let $H$ be a subgroup of $G$. Let $\alpha$ be the inclusion $H \hookrightarrow G$, and let $\beta$ be the identity map on a $G$-module $M$. In this case, we obtain restriction homomorphisms

$$
\operatorname{Res}: H^{r}(G, M) \rightarrow H^{r}(H, M) .
$$

They can also be constructed as follows: let $M \rightarrow \operatorname{Ind}_{H}^{G}(M)$ be the homomorphism sending $m$ to the map $g \mapsto g m$; the composite of the homomorphism this defines on cohomology with the isomorphism in Shapiro's Lemma,

$$
H^{r}(G, M) \rightarrow H^{r}\left(G, \operatorname{Ind}_{H}^{G}(M)\right) \xrightarrow{\simeq} H^{r}(H, M)
$$

is the restriction map.\\
(c) Let $H$ be a normal subgroup of $G$, let $\alpha$ be the quotient map $G \rightarrow G / H$, and let $\beta$ be the inclusion $M^{H} \hookrightarrow M$. In this case, we obtain the inflation homomorphisms

$$
\operatorname{Inf}: H^{r}\left(G / H, M^{H}\right) \rightarrow H^{r}(G, M) .
$$

(d) For any $g_{0} \in G$, the homomorphisms $\alpha: G \rightarrow G, \sigma \mapsto g_{0} \sigma g_{0}^{-1}$, and $\beta: M \rightarrow M$, $m \mapsto g_{0}^{-1} m$, are compatible. I claim that the homomorphisms

$$
H^{r}(G, M) \rightarrow H^{r}(G, M)
$$

they define are each the identity map. For $r=0$, the homomorphism is

$$
m \mapsto g_{0}^{-1} m: M^{G} \rightarrow M^{G},
$$

which is obviously the identity. Let $r>0$, and suppose the statement is known for $r-1$. From the exact sequence (15, p. 62)

$$
0 \rightarrow M \rightarrow M_{*} \rightarrow M_{\dagger} \rightarrow 0,
$$

$M_{*}=\operatorname{Ind}^{G}\left(M_{0}\right)$, we obtain a diagram with exact rows\\
\includegraphics[max width=\textwidth, alt={}, center]{e7cfc28e-5737-4710-b4d2-10a47512e540-13_198_1026_1998_520}\\
The 0's at right result from the fact that $M_{*}$ is an induced module. The vertical maps are those defined by the pair $(\alpha, \beta)$ (for the different modules). One checks easily that the diagram commutes. By induction, the middle vertical map is the identity, which implies that the third is also.

Remark 1.28 (a) The method of proof in (d) is called dimension shifting.\\
(b) Let $H$ be a normal subgroup of $G$. For every $G$-module $M$, the recipe in (d) gives an action of $G$ on $H^{r}(H, M)$, which the above result shows to factor through $G / H$.

Example 1.29 We shall need one more functorial map of cohomology groups. Let $H$ be a subgroup of finite index of $G$, and let $S$ be a set of left coset representatives for $H$ in $G$, $G=\bigcup_{s \in S} s H$. Let $M$ be a $G$-module. For every $m \in M^{H}$,

$$
\operatorname{Nm}_{G / H} m \stackrel{\text { def }}{=} \sum_{s \in S} s m
$$

is independent of the choice of $S$, and is fixed by $G$. Thus $\operatorname{Nm}_{G / H}$ is a homomorphism

$$
M^{H} \rightarrow M^{G} .
$$

This can be extended to a corestriction homomorphism ${ }^{1}$

$$
\text { Cor: } H^{r}(H, M) \rightarrow H^{r}(G, M)
$$

for all $r$ as follows: for every $G$-module $M$, there is a canonical homomorphism of $G$ -modules

$$
\varphi \mapsto \sum_{s \in S} s \varphi\left(s^{-1}\right): \operatorname{Ind}_{H}^{G} M \rightarrow M ;
$$

the map on cohomology which it defines, when composed with the isomorphism in Shapiro's lemma, gives Cor,

$$
H^{r}(H, M) \xrightarrow{\simeq} H^{r}\left(G, \operatorname{Ind}_{H}^{G} M\right) \rightarrow H^{r}(G, M) .
$$

Proposition 1.30 Let $H$ be a subgroup of $G$ of finite index. The composite

$$
\text { Cor ∘ Res: } H^{r}(G, M) \rightarrow H^{r}(G, M)
$$

is multiplication by $(G: H)$.\\
Proof. For $m \in M$, let $\varphi_{m}$ be the map $g \mapsto g m: G \rightarrow M$. Then Cor ∘ Res is the map on cohomology defined by the composite of

$$
M \rightarrow \operatorname{Ind}_{H}^{G}(M) \rightarrow M, \quad m \mapsto \varphi_{m} \mapsto \sum_{s} s \varphi_{m}\left(s^{-1}\right)=\sum_{s} m=(G: H) m .
$$ \(\square\)

Corollary 1.31 If $(G: 1)=m$, then $m H^{r}(G, M)=0$ for $r>0$.\\
Proof. According to the proposition, multiplication by $m$ factors through $H^{r}(\{1\}, M)=$ 0,

$$
H^{r}(G, M) \xrightarrow{\text { Res }} H^{r}(\{1\}, M) \xrightarrow{\text { Cor }} H^{r}(G, M) .
$$ \(\square\)

Corollary 1.32 If $G$ is finite and $M$ is finitely generated as an abelian group, then $H^{r}(G, M)$ is finite.

Proof. From the description of $H^{r}(G, M)$ as the group of cocycles modulo coboundaries, it is clear that $H^{r}(G, M)$ is finitely generated as an abelian group, and we have just seen that it is killed by $(G: 1)$. Therefore it is finite. \(\square\)

\footnotetext{${ }^{1}$ The corestriction homomorphism is also called the transfer homomorphism.
}For an abelian group $A$ and prime number $p$, the $p$-primary component $A(p)$ of $A$ is the subgroup consisting of all elements killed by a power of $p$.

Corollary 1.33 Let $G$ be a finite group, and let $G_{p}$ be its Sylow $p$-subgroup. For every $G$-module $M$, the restriction map

$$
\operatorname{Res}: H^{r}(G, M) \rightarrow H^{r}\left(G_{p}, M\right)
$$

is injective on the $p$-primary component of $H^{r}(G, M)$.\\
Proof. The composite

$$
\text { Cor ∘ Res : } H^{r}(G, M) \rightarrow H^{r}\left(G_{p}, M\right) \rightarrow H^{r}(G, M)
$$

is multiplication by $\left(G: G_{p}\right)$, which is not divisible by $p$, and so is injective on the $p$ -primary component of $H^{r}(G, M)$. \(\square\)

The inflation-restriction exact sequence\\
Proposition 1.34 Let $H$ be a normal subgroup of $G$, and let $M$ be a $G$-module. Let $r$ be an integer $>0$. If $H^{j}(H, M)=0$ for all $j$ with $0<j<r$, then the sequence

$$
0 \rightarrow H^{r}\left(G / H, M^{H}\right) \xrightarrow{\operatorname{Inf}} H^{r}(G, M) \xrightarrow{\operatorname{Res}} H^{r}(H, M)
$$

is exact.\\
Proof. We first prove this for $r=1$. In this case, the hypothesis on $H^{j}(H, M)$ is vacuous. Let $\varphi: G \rightarrow M$ be a crossed homomorphism whose restriction to $H$ is principal, say, $\varphi(h)=h m_{0}-m_{0}$. Define $\varphi^{\prime}(g)=\varphi(g)-\left(g m_{0}-m_{0}\right)$. Then $\varphi^{\prime}$ represents the same class in $H^{1}(G, M)$, but now $\varphi^{\prime}(h)=0$ for all $h \in H$, and so $\varphi^{\prime}$ comes by "inflation" from a crossed homomorphism $G / H \rightarrow M$.

We show that $\varphi^{\prime}$ takes values in $M^{H}$. As $\varphi^{\prime}$ is a crossed homomorphism and $\varphi^{\prime}(h)=0$ for $h \in H$, we have

$$
\begin{aligned}
& \varphi^{\prime}(h g)=h \varphi^{\prime}(g)+\varphi^{\prime}(h)=h \varphi^{\prime}(g) \\
& \varphi^{\prime}(g h)=g \varphi^{\prime}(h)+\varphi^{\prime}(g)=\varphi^{\prime}(g) .
\end{aligned}
$$

The second equation says that $\varphi^{\prime}$ is constant on left cosets of $H$. As $H$ is normal, left cosets and right cosets are the same, and so $\varphi^{\prime}$ is constant on right cosets. Since $h g$ and $g$ are in the same right coset $H g$, we have $\varphi^{\prime}(h g)=\varphi^{\prime}(g)$. Combining this with the first equation gives $h \varphi^{\prime}(g)=\varphi^{\prime}(g)$, i.e., $\varphi^{\prime}(g) \in M^{H}$.

We have shown that $\varphi^{\prime}$ comes by "inflation" from a crossed homomorphism $G / H \rightarrow$ $M^{H}$, which proves the exactness at $H^{1}(G, M)$. The proof of the exactness at $H^{1}\left(G / H, M^{H}\right)$ is easy.

Now assume that $r>1$ and that the statement is true for $r-1$. Consider the exact sequence (15), p. 62,

$$
0 \rightarrow M \rightarrow M_{*} \rightarrow M_{\dagger} \rightarrow 0 .
$$

Then

$$
H^{j}\left(H, M_{\dagger}\right) \simeq H^{j+1}(H, M), \quad j>0,
$$

and so $H^{j}\left(H, M_{\dagger}\right)=0$ for $0<j<r-1$. By induction, the sequence

$$
0 \rightarrow H^{r-1}\left(G / H, M_{\dagger}^{H}\right) \xrightarrow{\operatorname{Inf}} H^{r-1}\left(G, M_{\dagger}\right) \xrightarrow{\operatorname{Res}} H^{r-1}\left(H, M_{\dagger}\right)
$$

is exact, and this is isomorphic to the sequence

$$
0 \rightarrow H^{r}\left(G / H, M^{H}\right) \xrightarrow{\operatorname{Inf}} H^{r}(G, M) \xrightarrow{\operatorname{Res}} H^{r}(H, M) .
$$ \(\square\)

Remark 1.35 (a) With a little more effort, one can extend the sequence to

$$
\begin{aligned}
0 \rightarrow & H^{1}\left(G / H, M^{H}\right) \xrightarrow{\operatorname{Inf}} H^{1}(G, M) \xrightarrow{\operatorname{Res}} H^{1}(H, M)^{G / H} \rightarrow \\
& H^{2}\left(G / H, M^{H}\right) \xrightarrow{\operatorname{Inf}} H^{2}(G, M)^{*} \rightarrow H^{1}\left(G / H, H^{1}(H, M)\right) \rightarrow H^{3}\left(G / H, M^{H}\right),
\end{aligned}
$$

where $H^{2}(G, M)^{*}=\operatorname{Ker}\left(\operatorname{Res}: H^{2}(G, M) \rightarrow H^{2}(H, M)\right)$. See Dekimpe, Hartl, and Wauters, J. Algebra 369 (2012), 70-95.\\
(b) For the experts, the Hochschild-Serre (or Lyndon) spectral sequence takes the form

$$
H^{i}\left(G / H, H^{j}(H, M)\right) \Rightarrow H^{i+j}(G, M) ;
$$

here $G / H$ acts on $H^{j}(H, M)$ as described in (1.28b). It is possible to read off from this sequence, an exact sequence

$$
0 \rightarrow H^{1}\left(G / H, M^{H}\right) \xrightarrow{\operatorname{Inf}} H^{1}(G, M) \xrightarrow{\operatorname{Res}} H^{1}(H, M) .
$$

The hypotheses that $H^{i}(H, M)=0$ for $0<i<r$ forces $r$ rows of the spectral sequence to vanish, and allows one to read off the same sequence with 1 replaced with $r$.

Example 1.36 Let $\Omega \supset L$ be Galois extensions of $K$. Then $H \stackrel{\text { def }}{=} \operatorname{Gal}(\Omega / L)$ is a normal subgroup of $G \stackrel{\text { def }}{=} \operatorname{Gal}(\Omega / K)$. According to Proposition 1.22, $H^{1}\left(H, \Omega^{\times}\right)=0$, and so there is an exact sequence

$$
0 \rightarrow H^{2}\left(G / H, L^{\times}\right) \rightarrow H^{2}\left(G, \Omega^{\times}\right) \rightarrow H^{2}\left(H, \Omega^{\times}\right) .
$$

A direct proof (not involving dimension shifting) that this sequence is exact can be found in Artin 1951, 6.4.

\section*{Cup-products}
Lemma 1.37 The exact sequence (15), p. 62,

$$
0 \rightarrow M \rightarrow M_{*} \rightarrow M_{\dagger} \rightarrow 0,
$$

is split as a sequence of abelian groups, i.e., $M_{*} \approx M \oplus M_{\dagger}$ (as abelian groups).\\
Proof. Recall that $M_{*}$ consists of the maps $\varphi: G \rightarrow M$ and that the first homomorphism in the sequence sends an element $m$ of $M$ to $\varphi_{m}$, where $\varphi_{m}(g)=g m$. For $\varphi \in M_{*}$, let $s(\varphi)=\varphi\left(1_{G}\right)$. Then $s$ is a homomorphism of abelian groups $M_{*} \rightarrow M$ such that $s\left(\varphi_{m}\right)=m$, and so

$$
M_{*} \simeq \operatorname{Ker}(1-s) \oplus \operatorname{Ker}(s) \simeq M \oplus M_{\dagger}
$$

as abelian groups. \(\square\)

Let $G$ be a group. For $G$-modules $M$ and $N$, we write $M \otimes N$ for $M \otimes_{\mathbb{Z}} N$ regarded as a $G$-module with

$$
g(m \otimes n)=g m \otimes g n, \quad g \in G, \quad m \in M, \quad n \in N .
$$

Proposition 1.38 There exists one and only one family of bi-additive pairings

$$
(m, n) \mapsto m \cup n: H^{r}(G, M) \times H^{s}(G, N) \rightarrow H^{r+s}(G, M \otimes N),
$$

defined for all $G$-modules $M, N$ and all integers $r, s \geq 0$, satisfying the following conditions:

\begin{itemize}
  \item[(a)] these maps become morphisms of functors when the two sides are regarded as covariant bifunctors on $(M, N)$;
  \item[(b)] for $r=s=0$, the pairing is
$$
(m, n) \mapsto m \otimes n: M^{G} \otimes N^{G} \rightarrow(M \otimes N)^{G} ;
$$
  \item[(c)] if $0 \rightarrow M^{\prime} \rightarrow M \rightarrow M^{\prime \prime} \rightarrow 0$ is an exact sequence of $G$-modules such that
$$
0 \rightarrow M^{\prime} \otimes N \rightarrow M \otimes N \rightarrow M^{\prime \prime} \otimes N \rightarrow 0
$$
is exact, then
$$
\left(\delta m^{\prime \prime}\right) \cup n=\delta\left(m^{\prime \prime} \cup n\right), \quad m^{\prime \prime} \in H^{r}\left(G, M^{\prime \prime}\right), \quad n \in H^{s}(G, N) .
$$
Here $\delta$ denotes the connecting homomorphism $H^{r}\left(G, M^{\prime \prime}\right) \rightarrow H^{r+1}\left(G, M^{\prime}\right)$ or $H^{r+s}\left(G, M^{\prime \prime} \otimes N\right) \rightarrow H^{r+s+1}\left(G, M^{\prime} \otimes N\right)$.
  \item[(d)] if $0 \rightarrow N^{\prime} \rightarrow N \rightarrow N^{\prime \prime} \rightarrow 0$ is an exact sequence of $G$-modules such that
$$
0 \rightarrow M \otimes N^{\prime} \rightarrow M \otimes N \rightarrow M \otimes N^{\prime \prime} \rightarrow 0
$$
is exact, then
$$
m \cup \delta n^{\prime \prime}=(-1)^{r} \delta\left(m \cup n^{\prime \prime}\right), \quad m \in H^{r}(G, M), \quad n^{\prime \prime} \in H^{s}\left(G, N^{\prime \prime}\right) .
$$
\end{itemize}

Proof. The uniqueness is proved using dimension shifting, a key point being that, because the sequence (15) is split-exact as a sequence of abelian groups (see 1.37), it remains exact when tensored over $\mathbb{Z}$ with a $G$-module. For the existence, one proves that the pairing defined as follows has the required properties: let $m \in H^{r}(G, M)$ and $n \in H^{s}(G, N)$ be represented by the (inhomogeneous) cocycles $\varphi$ and $\psi$; then $m \cup n$ is represented by the cocycle

$$
\left(g_{1}, \ldots, g_{r+s}\right) \mapsto \varphi\left(g_{1}, \ldots, g_{r}\right) \otimes g_{1} \cdots g_{r} \psi\left(g_{r+1}, \cdots, g_{r+s}\right) .
$$

It takes several pages to write out the details, but the proof is not difficult. \(\square\)

Proposition 1.39 The following formulas hold:

\begin{itemize}
  \item[(a)] $(x \cup y) \cup z=x \cup(y \cup z)$ (in $H^{r+s+t}(G, M \otimes N \otimes P)$ );
  \item[(b)] $x \cup y=(-1)^{r s} y \cup x$ if $x \in H^{r}(G, M), y \in H^{s}(G, N)$ (we identify $M \otimes N$ and $N \otimes M$ );
  \item[(c)] $\operatorname{Res}(x \cup y)=\operatorname{Res}(x) \cup \operatorname{Res}(y)$;
\end{itemize}

\begin{itemize}
  \item[(d)] $\operatorname{Cor}(x \cup \operatorname{Res} y)=\operatorname{Cor}(x) \cup y$.
  \item[(e)] $\operatorname{Inf}(x \cup y)=\operatorname{Inf}(x) \cup \operatorname{Inf}(y)$.
\end{itemize}

Proof. In each case, one verifies the formula when $x, y, \ldots$ have degree 0 , and then uses dimension shifting to deduce the general case. For example, the proof of (d) is written out in detail in Atiyah and Wall 1967, Proposition 9. For (e), see Weiss 1969, 4-3-9. \(\square\)

A pairing


\begin{equation*}
x, y \mapsto(x, y): M \times N \rightarrow P \tag{18}
\end{equation*}


such that

$$
g(x, y)=(g x, g y), \quad \text { all } g \in G,
$$

defines a homomorphism of $G$-modules

$$
M \otimes N \rightarrow P,
$$

and hence maps

$$
H^{r}(G, M \otimes N) \rightarrow H^{r}(G, P) .
$$

The composites of the cup-product pairings with these maps, namely, the pairings

$$
H^{r}(G, M) \times H^{s}(G, N) \rightarrow H^{r+s}(G, P),
$$

will be referred to as the cup-product pairings defined by (18).

\section*{2 Homology}
Definition of the homology groups\\
For a $G$-module $M$, let $M_{G}$ be the largest quotient of $M$ on which $G$ acts trivially. Thus $M_{G}$ is the quotient of $M$ by the subgroup generated by

$$
\{g m-m \mid g \in G, \quad m \in M\} .
$$

Note that this is the dual notion to $M^{G}$, which is the largest subobject of $M$ on which $G$ acts trivially. The definition of the cohomology groups dualizes to give us homology groups.

In detail, the functor

$$
M \mapsto M_{G}: \operatorname{Mod}_{G} \rightarrow \mathbf{A b}
$$

is right exact, i.e., if

$$
0 \rightarrow M^{\prime} \rightarrow M \rightarrow M^{\prime \prime} \rightarrow 0
$$

is exact, then

$$
M_{G}^{\prime} \rightarrow M_{G} \rightarrow M_{G}^{\prime \prime} \rightarrow 0
$$

is exact. This can be checked directly, or by observing that passing from $M$ to $M_{G}$ amounts to tensoring $M$ with $\mathbb{Z}[G] / I_{G}$, where $I_{G}$ is the augmentation ideal

$$
\operatorname{Ker}\left(\sum n_{g} g \mapsto \sum n_{g}: \mathbb{Z}[G] \rightarrow \mathbb{Z}\right) .
$$

An object $P$ of an abelian category is projective if for any object $N$ and quotient object $M$, every morphism $P \rightarrow M$ lifts to a morphism $P \rightarrow N$. For example, every free $\mathbb{Z}[G]$-module is projective as a $\mathbb{Z}[G]$ - (equivalently, $G$-) module.

Let $M$ be a $G$-module. Let $\left(m_{i}\right)_{i \in I}$ be a family of generators for $M$ as a $\mathbb{Z}[G]$-module, and let $\mathbb{Z}[G]^{(I)}$ be a direct sum of copies of $\mathbb{Z}[G]$ indexed by $I$. The map $\sum_{i \in I} \gamma_{i} \mapsto$ $\sum \gamma_{i} m_{i}$ is a surjective $G$-homomorphism $\mathbb{Z}[G]^{(I)} \rightarrow M$. This shows that every $G$-module is a quotient of a free $G$-module, and hence that the category $\operatorname{Mod}_{G}$ has enough projectives.

Let $M$ be a $G$-module, and choose a projective resolution

$$
\cdots \longrightarrow P_{2} \xrightarrow{d_{2}} P_{1} \xrightarrow{d_{1}} P_{0} \xrightarrow{d_{0}} M \longrightarrow 0
$$

of $M$. The complex

$$
\cdots \longrightarrow\left(P_{2}\right)_{G} \xrightarrow{d_{2}}\left(P_{1}\right)_{G} \xrightarrow{d_{1}}\left(P_{0}\right)_{G} \longrightarrow 0
$$

need no longer be exact, and we set

$$
H_{r}(G, M)=\frac{\operatorname{Ker}\left(d_{r}\right)}{\operatorname{Im}\left(d_{r+1}\right)} .
$$

These groups have the following basic properties.\\
2.1 The zeroth group $H_{0}(G, M)=M_{G}$, because

$$
P_{1 G} \rightarrow P_{0 G} \rightarrow M_{G} \rightarrow 0
$$

is exact.\\
2.2 If $P$ is a projective $G$-module, then $H_{r}(G, P)=0$ for all $r>0$, because

$$
\cdots \rightarrow 0 \rightarrow P \rightarrow P \rightarrow 0
$$

is a projective resolution of $P$.\\
2.3 Let $P_{\bullet} \rightarrow M$ and $Q_{\bullet} \rightarrow N$ be projective resolutions of the $G$-modules $M$ and $N$. Any homomorphism $\alpha: M \rightarrow N$ of $G$-modules extends to a morphism of complexes\\
\includegraphics[max width=\textwidth, alt={}, center]{e7cfc28e-5737-4710-b4d2-10a47512e540-19_176_229_1752_923}\\
and the homomorphisms

$$
H_{r}\left(\alpha_{\bullet}\right): H_{r}\left(P_{\bullet}\right) \rightarrow H_{r}\left(Q_{\bullet}\right)
$$

are independent of the choice of $\alpha_{\bullet}$. When we apply this statement to the identity map id: $M \rightarrow M$, we find that the groups $H_{r}(G, M)$ are well-defined up to a well-defined isomorphism. The statement for a general $\alpha$ then shows that $M \mapsto H_{r}(G, M)$ is a functor from the category of $G$-modules to the category of abelian groups.\\
2.4 A short exact sequence

$$
0 \rightarrow M^{\prime} \rightarrow M \rightarrow M^{\prime \prime} \rightarrow 0
$$

of $G$-modules gives rise to a long exact sequence

$$
\cdots \rightarrow H_{r}(G, M) \rightarrow H_{r}\left(G, M^{\prime \prime}\right) \xrightarrow{\delta_{r}} H_{r-1}\left(G, M^{\prime}\right) \rightarrow \cdots \rightarrow H_{0}\left(G, M^{\prime \prime}\right) \rightarrow 0 .
$$

Moreover, the association\\
short exact sequence ⟼ long exact sequence\\
is functorial, i.e., a morphism of short exact sequences induces a morphism of long exact sequences.

Remark 2.5 The family of functors $\left(H_{r}(G, \cdot)\right)_{r \geq 0}$ and the boundary maps $\delta_{r}$ are uniquely determined by the properties (2.1, 2.2, 2.4).

Just as in the case of cohomology, it is possible to give an explicit description of $H_{r}(G, M)$ as the quotient of a group of $r$-cycles by a subgroup of $r$-boundaries-see the references later in this chapter.

The group $H_{1}(G, \mathbb{Z})$\\
Using only the properties of the homology groups listed above, we shall compute $H_{1}(G, \mathbb{Z})$.\\
The augmentation map is

$$
\mathbb{Z}[G] \rightarrow \mathbb{Z}, \quad \sum n_{g} g \mapsto \sum n_{g} .
$$

Its kernel is called the augmentation ideal $I_{G}$. Clearly $I_{G}$ is a free $\mathbb{Z}$-submodule of $\mathbb{Z}[G]$ with basis $\{g-1 \mid g \in G, g \neq 1\}$, and so

$$
M / I_{G} M=M_{G} \stackrel{(2.1)}{\simeq} H_{0}(G, M) .
$$

Consider the exact (augmentation) sequence:


\begin{equation*}
0 \rightarrow I_{G} \rightarrow \mathbb{Z}[G] \rightarrow \mathbb{Z} \rightarrow 0 . \tag{19}
\end{equation*}


The $G$-module $\mathbb{Z}[G]$ is projective (because it is a free $\mathbb{Z}[G]$-module), and so $H_{1}(G, \mathbb{Z}[G])=$ 0 . Therefore we obtain an exact sequence of homology groups

$$
0 \rightarrow H_{1}(G, \mathbb{Z}) \rightarrow I_{G} / I_{G}^{2} \rightarrow \mathbb{Z}[G] / I_{G} \rightarrow \mathbb{Z} \rightarrow 0 .
$$

The middle map is induced by the inclusion $I_{G} \hookrightarrow \mathbb{Z}[G]$, and so is zero. Therefore the sequence shows that


\begin{equation*}
H_{1}(G, \mathbb{Z}) \xrightarrow{\simeq} I_{G} / I_{G}^{2} \tag{20}
\end{equation*}


and

$$
\mathbb{Z}[G]_{G} \xrightarrow{\simeq} \mathbb{Z}
$$

i.e., $\mathbb{Z}$ is the largest quotient of $\mathbb{Z}[G]$ on which $G$ acts trivially. Note that $I_{G}^{2} \stackrel{\text { def }}{=} I_{G} \cdot I_{G}$ is the $\mathbb{Z}$-submodule of $\mathbb{Z}[G]$ generated by elements of the form

$$
(g-1)\left(g^{\prime}-1\right), \quad g, g^{\prime} \in G .
$$

Lemma 2.6 Let $G^{c}$ be the commutator subgroup of $G$, so that $G / G^{c}$ is the largest abelian quotient $G^{\mathrm{ab}}$ of $G$. Then the map $g \mapsto(g-1)+I_{G}^{2}$ induces an isomorphism


\begin{equation*}
\frac{G}{G^{c}} \rightarrow \frac{I_{G}}{I_{G}^{2}} . \tag{21}
\end{equation*}


Proof. Consider the map

$$
g \mapsto(g-1)+I_{G}^{2}: G \rightarrow I_{G} / I_{G}^{2} .
$$

This is a homomorphism because

$$
\begin{aligned}
g g^{\prime}-1 & =(g-1)\left(g^{\prime}-1\right)+(g-1)+\left(g^{\prime}-1\right) \\
& \equiv(g-1)+\left(g^{\prime}-1\right) \bmod I_{G}^{2} .
\end{aligned}
$$

Since $I_{G} / I_{G}^{2}$ is commutative, the map factors through $G^{\mathrm{ab}}$. To prove that it is an isomorphism, we construct an inverse. Recall that $I_{G}$ is freely generated as a $\mathbb{Z}$-module by the elements $g-1$. Consider the homomorphism $I_{G} \rightarrow G^{\mathrm{ab}}$ mapping $g-1$ to the class of $g$. From the identity

$$
(g-1)\left(g^{\prime}-1\right)=\left(g g^{\prime}-1\right)-(g-1)-\left(g^{\prime}-1\right)
$$

we see that $(g-1)\left(g^{\prime}-1\right)$ maps to 1 . Since $I_{G}^{2}$ is generated by elements of this form, this shows that the map factors through $I_{G} / I_{G}^{2}$. The two maps we have constructed are inverse. \(\square\)

Proposition 2.7 There is a canonical isomorphism

$$
H_{1}(G, \mathbb{Z}) \simeq G^{\mathrm{ab}} .
$$

Proof. Combine the isomorphism (20) with the inverse of (21). \(\square\)

Aside 2.8 For any group $G$, there exists a topological space $B G$, called the classifying space of $G$, such that $G=\pi_{1}(B G)$ and $H_{r}(B G, \mathbb{Z})=H_{r}(G, \mathbb{Z})$ for all $r$ (Rosenberg 1994, 5.1.16, 5.1.27). In terms of $B G$, the proposition simply states that $H_{1}(B G, \mathbb{Z}) \simeq \pi_{1}(B G)^{\mathrm{ab}}$.

\section*{3 The Tate groups}
Throughout this section, $G$ is a finite group.\\
For a $G$-module $M$, the norm map $\operatorname{Nm}_{G}: M \rightarrow M$ is defined to be

$$
m \mapsto \sum_{g \in G} g m .
$$

Let $g^{\prime} \in G$. As $g$ runs through the elements of $G$, so also do $g^{\prime} g$ and $g g^{\prime}$, and so $g^{\prime}\left(\operatorname{Nm}_{G}(m)\right)=\operatorname{Nm}_{G}(m)=\operatorname{Nm}_{G}\left(g^{\prime} m\right)$. Hence

$$
\operatorname{Im}\left(\operatorname{Nm}_{G}\right) \subset M^{G}, \quad I_{G} M \subset \operatorname{Ker} \operatorname{Nm}_{G} .
$$

Therefore, we get an exact commutative diagram,\\
\includegraphics[max width=\textwidth, alt={}, center]{e7cfc28e-5737-4710-b4d2-10a47512e540-21_195_1424_2219_356}\\
As $H_{0}(G, M)=M / I_{G} M$ and $H^{0}(G, M)=M^{G}$, the bottom row of this can be rewritten

$$
0 \rightarrow \operatorname{Ker}\left(\operatorname{Nm}_{G}\right) / I_{G} M \rightarrow H_{0}(G, M) \xrightarrow{\operatorname{Nm}_{G}} H^{0}(G, M) \rightarrow M^{G} / \operatorname{Nm}_{G}(M) \rightarrow 0 .
$$

Tate defined

$$
H_{T}^{r}(G, M)= \begin{cases}H^{r}(G, M) & r>0 \\ M^{G} / \mathrm{Nm}_{G}(M) & r=0 \\ \operatorname{Ker}\left(\operatorname{Nm}_{G}\right) / I_{G} M & r=-1 \\ H_{-r-1}(G, M) & r<-1\end{cases}
$$

Thus, the exact sequence now becomes

$$
0 \rightarrow H_{T}^{-1}(G, M) \rightarrow H_{0}(G, M) \xrightarrow{\mathrm{Nm}_{G}} H^{0}(G, M) \rightarrow H_{T}^{0}(G, M) \rightarrow 0 .
$$

The groups $H_{T}^{r}(G, M)$ are known as the Tate cohomology groups. Often $H_{T}^{r}(G, M)$ is denoted by $\hat{H}^{r}(G, M)$. Since it causes no ambiguity, we sometimes omit the subscript $T$ when $r>0$.

For any short exact sequence of $G$-modules

$$
0 \rightarrow M^{\prime} \rightarrow M \rightarrow M^{\prime \prime} \rightarrow 0
$$

we get a diagram\\
\includegraphics[max width=\textwidth, alt={}, center]{e7cfc28e-5737-4710-b4d2-10a47512e540-22_453_1439_1242_365}

On applying the extended snake lemma (A.1) to the middle part of the diagram, we get a (very) long exact sequence


\begin{equation*}
\cdots \rightarrow H_{T}^{r}\left(G, M^{\prime}\right) \rightarrow H_{T}^{r}(G, M) \rightarrow H_{T}^{r}\left(G, M^{\prime \prime}\right) \xrightarrow{\delta} H_{T}^{r+1}(G, M) \rightarrow \cdots \tag{22}
\end{equation*}


( $-\infty<r<\infty$ ).\\
Proposition 3.1 If $M$ is induced, then $H_{T}^{r}(G, M)=0$ for all $r \in \mathbb{Z}$.\\
Proof. For $r>0$, this was proved in (1.12).\\
Case $r=0$. Recall (1.3) that $M \simeq \mathbb{Z}[G] \otimes_{\mathbb{Z}} X$ for some abelian group $X$, and so an element $m$ of $M$ can be written uniquely in the form $m=\sum_{g}\left(g \otimes x_{g}\right)$ with $x_{g} \in X$. If $m$ is fixed by $g^{\prime} \in G$, then $\sum_{g}\left(g^{\prime} g \otimes x_{g}\right)=\sum_{g}\left(g \otimes x_{g}\right)$; in particular, ${ }^{2} g^{\prime} \otimes x_{e}=$ $g^{\prime} \otimes x_{g^{\prime}}$, i.e., $x_{e}=x_{g^{\prime}}$. Therefore, if $m \in M^{G}$, then $x_{g}=x_{e}$ for all $g \in G$, and $m=\left(\sum_{g} g\right) \otimes x_{e}=\operatorname{Nm}_{G}\left(e \otimes x_{e}\right)$.

\footnotetext{${ }^{2}$ Here, $e$ is the neutral element in $G$.
}Case $r=-1$. If $\operatorname{Nm}_{G}\left(\sum_{g} g \otimes x_{g}\right)=0$, then $\sum_{g} x_{g}=0$, and so

$$
\sum_{g}\left(g \otimes x_{g}\right)=\sum_{g}\left((g-1) \otimes x_{g}\right) \in I_{G} M .
$$

Case $r<-1$. Let $M=\mathbb{Z}[G] \otimes_{\mathbb{Z}} X$. Write $X$ as a quotient $X_{0} \rightarrow X$ of a free abelian group $X_{0}$. The kernel of this map will also be a free abelian group, ${ }^{3}$ and so we have an exact sequence

$$
0 \rightarrow X_{1} \rightarrow X_{0} \rightarrow X \rightarrow 0
$$

with $X_{0}$ and $X_{1}$ free abelian groups. On tensoring this with $\mathbb{Z}[G]$, we obtain an exact sequence of $\mathbb{Z}[G]$-modules


\begin{equation*}
0 \rightarrow M_{1} \rightarrow M_{0} \rightarrow M \rightarrow 0 . \tag{23}
\end{equation*}


Now

$$
H_{T}^{r}\left(G, M_{0}\right)=0=H_{T}^{r}\left(G, M_{1}\right)
$$

for $r<-1$ because $M_{0}$ and $M_{1}$ are free $\mathbb{Z}[G]$-modules and for $r \geq-1$ because $M_{0}$ and $M_{1}$ are induced. The (very) long exact sequence of (23) now shows that $H_{T}^{r}(G, M)=0$ for all $r$. \(\square\)

We know (15), p. 62, that every $G$-module $M$ is a submodule of an induced module. It is also a quotient of an induced $G$-module because, on tensoring the augmentation sequence (19) with $M$, we obtain an exact sequence of $G$-modules

$$
0 \rightarrow M^{\prime} \rightarrow \mathbb{Z}[G] \otimes_{\mathbb{Z}} M \rightarrow M \rightarrow 0 .
$$

This allows us to prove results by dimension shifting in both directions. Using this, it is possible to extend to all $r$, most of the results we proved for the groups $H^{r}(G, M)$ with $r \geq 0$. For example, Shapiro's lemma and its consequences are true, and the restriction and corestriction maps we defined in (1.27) and (1.29) extend to all the Tate cohomology groups. Specifically, there are canonical homomorphisms: ${ }^{4}$\\
Res: $H_{T}^{r}(G, M) \rightarrow H_{T}^{r}(H, M) \quad(H$ a subgroup of $G) ;$\\
Cor: $H_{T}^{r}(H, M) \rightarrow H_{T}^{r}(G, M) \quad(H$ a subgroup of $G)$.\\
Moreover, there is a natural action of $G / H$ on $H_{T}^{r}(H, M)$ for all $r \in \mathbb{Z}$.\\
The composite Res ○ Cor is still multiplication by ( $G: H$ ). As $H_{T}^{r}(\{1\}, M)=0$ for all $r$, the argument in the proof of (1.30) shows that $H_{T}^{r}(G, M)$ is killed by $|G|$ for all $r$.

Note that

$$
H_{T}^{-2}(G, \mathbb{Z})=H_{1}(G, \mathbb{Z}) \simeq G / G^{c} .
$$

Proposition 3.2 Let $H$ be a subgroup of $G$.

\begin{itemize}
  \item[(a)] The map Cor: $H_{T}^{-2}(H, \mathbb{Z}) \rightarrow H_{T}^{-2}(G, \mathbb{Z})$ corresponds to the map $H / H^{c} \rightarrow G / G^{c}$ induced by the inclusion $H \hookrightarrow G$.
  \item[(b)] The map Res: $H_{T}^{-2}(G, \mathbb{Z}) \rightarrow H_{T}^{-2}(H, \mathbb{Z})$ corresponds to the Verlagerung map $G / G^{c} \rightarrow$ $H / H^{c}(V, 3.18)$.
\end{itemize}

\footnotetext{${ }^{3}$ We are using that every subgroup of a free abelian group is free abelian. Darij Grinberg points out that this fact is highly nontrivial (and nonconstructive), and that it is possible to avoid using it.\\
${ }^{4}$ However, the inflation map $\operatorname{Inf}: H_{T}^{r}\left(G / H, M^{H}\right) \rightarrow H_{T}^{r}(G, M), H$ a normal subgroup of $G$, is defined only for $r \geq 1$.
}Proof. See Weiss 1969, 3-5-5, 3-5-7. I'll include the proof in a future version of the notes. \(\square\)

There is a uniform explicit description of all the groups $H_{T}^{r}(G, M)$. In fact, there is an explicit complex $L_{\bullet}$ of $G$-modules (infinite in both directions) such that

$$
H_{T}^{r}(G, M)=H^{r}\left(\operatorname{Hom}_{G}\left(L_{\bullet}, M\right)\right) .
$$

For $r>0$ this leads directly to the description we gave above of $H^{r}(G, M)$ in terms of inhomogeneous cocycles and coboundaries. See Weiss 1969, Chapter I.

\section*{Cup-products}
When the group $G$ is finite, the cup-products extend in a unique way to all the cohomology groups, and have the same list of properties. See Weiss 1969, Chapter IV. A future version of the notes will contain a complete list of the properties we need (but not the proofs, which are lengthy but not difficult).

\section*{The cohomology of finite cyclic groups}
We first compute the cohomology of $\mathbb{Q}, \mathbb{Z}$, and $\mathbb{Q} / \mathbb{Z}$ regarded as $G$-modules with the trivial action.

Lemma 3.3 For every finite group $G$,

\begin{itemize}
  \item[(a)] $H_{T}^{r}(G, \mathbb{Q})=0$ all $r \in \mathbb{Z}$;
  \item[(b)] $H_{T}^{0}(G, \mathbb{Z})=\mathbb{Z} /(G: 1) \mathbb{Z}$ and $H^{1}(G, \mathbb{Z})=0$;
  \item[(c)] there is a canonical isomorphism
$$
\operatorname{Hom}(G, \mathbb{Q} / \mathbb{Z}) \rightarrow H^{2}(G, \mathbb{Z}) .
$$
\end{itemize}

Proof. (a) The group $\mathbb{Q}$ is uniquely divisible, i.e., for all nonzero integers $m$, multiplication by $m: \mathbb{Q} \rightarrow \mathbb{Q}$ is an isomorphism. Therefore the map $H_{T}^{r}(m): H_{T}^{r}(G, \mathbb{Q}) \rightarrow$ $H_{T}^{r}(G, \mathbb{Q})$, which is also multiplication by $m$, is an isomorphism. When we take $m=$ $(G: 1)$, we find that multiplication by $m$ on $H_{T}^{r}(G, \mathbb{Q})$ is both zero (see 1.30) and an isomorphism, which is possible only if $H_{T}^{r}(G, \mathbb{Q})=0$.\\
(b) Because $G$ acts trivially on $\mathbb{Z}, \mathbb{Z}^{G}=\mathbb{Z}$ and the norm map is multiplication by $(G: 1)$. Hence $H_{T}^{0}(G, \mathbb{Z})=\mathbb{Z} /(G: 1) \mathbb{Z}$. Moreover, $H^{1}(G, \mathbb{Z})=\operatorname{Hom}(G, \mathbb{Z})$ (see 1.18a), but, because $\mathbb{Z}$ is torsion-free, all homomorphisms $G \rightarrow \mathbb{Z}$ are zero.\\
(c) The cohomology sequence of

$$
0 \rightarrow \mathbb{Z} \rightarrow \mathbb{Q} \rightarrow \mathbb{Q} / \mathbb{Z} \rightarrow 0
$$

is an exact sequence\\
\includegraphics[max width=\textwidth, alt={}]{e7cfc28e-5737-4710-b4d2-10a47512e540-24_152_1092_2350_478} \(\square\)

Let $G$ be a finite cyclic group of order $m$, with generator $\sigma$. Then, by definition,

$$
\begin{aligned}
H_{T}^{0}(G, M) & =\frac{\operatorname{Ker}(\sigma-1)}{\operatorname{Im}\left(\mathrm{Nm}_{G}\right)} \\
H_{T}^{-1}(G, M) & =\frac{\operatorname{Ker}\left(\mathrm{Nm}_{G}\right)}{\operatorname{Im}(\sigma-1)} .
\end{aligned}
$$

Recall that in (1.20) we showed that $H^{1}(G, M)$ is isomorphic to the same group as $H_{T}^{-1}(G, M)$. In fact, the cohomology groups of a finite cyclic group are periodic with period 2.

Proposition 3.4 Let $G$ be a cyclic group of finite order. The choice of a generator for $G$ determines isomorphisms

$$
H_{T}^{r}(G, M) \xrightarrow{\simeq} H_{T}^{r+2}(G, M)
$$

for all $G$-modules $M$ and all $r \in \mathbb{Z}$.\\
Proof. Let $\sigma$ generate $G$. Then the sequence

$$
0 \rightarrow \mathbb{Z} \xrightarrow{m \mapsto \sum_{g \in G} g m} \mathbb{Z}[G] \xrightarrow{\sigma-1} \mathbb{Z}[G] \xrightarrow{\sigma^{i} \mapsto 1} \mathbb{Z} \rightarrow 0
$$

is exact. Because the groups in the sequence and the kernel $I_{G}$ of $\mathbb{Z}[G] \rightarrow \mathbb{Z}$ are free $\mathbb{Z}$-modules, the sequence remains exact after it is tensored with $M$. Thus

$$
0 \rightarrow M \rightarrow \mathbb{Z}[G] \otimes_{\mathbb{Z}} M \rightarrow \mathbb{Z}[G] \otimes_{\mathbb{Z}} M \rightarrow M \rightarrow 0
$$

is an exact sequence of $G$-modules. Recall (1.4) that $\mathbb{Z}[G] \otimes_{\mathbb{Z}} M \approx \mathbb{Z}[G] \otimes_{\mathbb{Z}} M_{0}$, where $M_{0}$ is $M$ regarded as abelian group, and so $H_{T}^{r}\left(G, \mathbb{Z}[G] \otimes_{\mathbb{Z}} M\right)=0$ for all $r$ by (3.1). Therefore (see 1.13), the sequence defines isomorphisms

$$
H_{T}^{r}(G, M) \xrightarrow{\simeq} H_{T}^{r+2}(G, M)
$$

for all $r$. \(\square\)

REMARK 3.5 Let $\gamma$ be the element of $H^{2}(G, \mathbb{Z})$ corresponding under the isomorphism $H^{2}(G, \mathbb{Z}) \simeq \operatorname{Hom}(G, \mathbb{Q} / \mathbb{Z})$ to the map sending the chosen generator $\sigma$ of $G$ to $1 / m$. Then the map $H_{T}^{r}(G, M) \rightarrow H_{T}^{r+2}(G, M)$ is $x \mapsto x \cup \gamma$. In a future version, I'll prove this. Serre, Local Fields, VIII, Section 4 gives the following hint: "the period isomorphisms are given by cup product with $\gamma$ : this follows, for example, from the formulas for the cup product given in Cartan-Eilenberg (Homological Algebra) p. 252." Alternatively, see Weiss 1969, 4-5-10.

Let $G$ be a finite cyclic group, and let $M$ be a $G$-module. When the cohomology groups $H^{r}(G, M)$ are finite, we define the Herbrand quotient of $M$ to be

$$
h(M)=\frac{\left|H_{T}^{0}(G, M)\right|}{\left|H_{T}^{1}(G, M)\right|} .
$$

Proposition 3.6 Let $0 \rightarrow M^{\prime} \rightarrow M \rightarrow M^{\prime \prime} \rightarrow 0$ be an exact sequence of $G$-modules. If any two of the Herbrand quotients $h\left(M^{\prime}\right), h(M), h\left(M^{\prime \prime}\right)$ are defined, then so also is the third, and

$$
h(M)=h\left(M^{\prime}\right) h\left(M^{\prime \prime}\right) .
$$

Proof. We can truncate the (very) long exact cohomology sequence as follows,

$$
0 \rightarrow K \rightarrow H_{T}^{0}\left(M^{\prime}\right) \rightarrow H_{T}^{0}(M) \rightarrow H_{T}^{0}\left(M^{\prime \prime}\right) \rightarrow H^{1}\left(M^{\prime}\right) \rightarrow H^{1}(M) \rightarrow H^{1}\left(M^{\prime \prime}\right) \rightarrow K^{\prime} \rightarrow 0,
$$

where

$$
K=\operatorname{Coker}\left(H_{T}^{-1}(M) \rightarrow H_{T}^{-1}\left(M^{\prime \prime}\right)\right) \approx \operatorname{Coker}\left(H^{1}(M) \rightarrow H^{1}\left(M^{\prime \prime}\right)\right)=K^{\prime} .
$$

The first statement is now obvious, and the second follows from the next lemma. \(\square\)

Lemma 3.7 Let

$$
0 \rightarrow A_{0} \rightarrow A_{1} \rightarrow \cdots \rightarrow A_{r} \rightarrow 0
$$

be a exact sequence of finite groups. Then

$$
\frac{\left|A_{0}\right| \cdot\left|A_{2}\right| \cdots}{\left|A_{1}\right| \cdot\left|A_{3}\right| \cdots}=1 .
$$

Proof. For a short exact sequence, that is, $r=2$, this is obvious, but every exact sequence can be broken up into short exact sequences,

$$
\begin{gathered}
0 \rightarrow A_{0} \rightarrow A_{1} \rightarrow C_{1} \rightarrow 0 \\
0 \rightarrow C_{1} \rightarrow A_{2} \rightarrow C_{2} \rightarrow 0 \\
\ldots \\
0 \rightarrow C_{r-1} \rightarrow A_{r-1} \rightarrow A_{r} \rightarrow 0
\end{gathered}
$$

Here $C_{i}=\operatorname{Coker}\left(A_{i-1} \rightarrow A_{i}\right)=\operatorname{Ker}\left(A_{i+1} \rightarrow A_{i+2}\right)$. From these sequences we find that

$$
1=\frac{\left|A_{0}\right| \cdot\left|C_{1}\right|}{\left|A_{1}\right|}=\frac{\left|A_{0}\right| \cdot\left|A_{2}\right|}{\left|A_{1}\right| \cdot\left|C_{2}\right|}=\cdots .
$$ \(\square\)

Proposition 3.8 If $M$ is a finite module, then $h(M)=1$.\\
Proof. Consider the exact sequences

$$
0 \rightarrow M^{G} \rightarrow M \xrightarrow{g-1} M \rightarrow M_{G} \rightarrow 0
$$

and

$$
0 \rightarrow H_{T}^{-1}(M) \rightarrow M_{G} \xrightarrow{\mathrm{Nm}_{G}} M^{G} \rightarrow H_{T}^{0}(M) \rightarrow 0,
$$

where $g$ is any generator of $G$. From the first sequence we find that $M^{G}$ and $M_{G}$ have the same order, and then from the second that $H_{T}^{-1}(M)$ and $H_{T}^{0}(M)$ have the same order. \(\square\)

Corollary 3.9 Let $\alpha: M \rightarrow N$ be a homomorphism of $G$-modules with finite kernel and cokernel. If either $h(M)$ or $h(N)$ is defined, then so also is the other, and they are equal.

Proof. Suppose $h(N)$ is defined, and consider the exact sequences:

$$
\begin{aligned}
& 0 \rightarrow \alpha(M) \rightarrow N \rightarrow \operatorname{Coker}(\alpha) \rightarrow 0 \\
& 0 \rightarrow \operatorname{Ker}(\alpha) \rightarrow M \rightarrow \alpha(M) \rightarrow 0
\end{aligned}
$$

From the first sequence, we find that $h(\alpha M)$ is defined and equals $h(N)$, and from the second sequence we find that $h(M)$ is defined and equals $h(\alpha M)$. \(\square\)

\section*{Tate's Theorem}
For the remainder of this section, all cohomology groups will be the Tate groups, and so we drop the subscript $T$ except in the main statements.

Theorem 3.10 Let $G$ be a finite group, and let $M$ be a $G$-module. If

$$
H^{1}(H, M)=0=H^{2}(H, M)
$$

for all subgroups $H$ of $G$, then $H_{T}^{r}(G, M)=0$ for all $r \in \mathbb{Z}$.\\
Proof. If $G$ is cyclic, this follows from the periodicity of the cohomology.\\
Assume now that $G$ is solvable. We shall prove the theorem in this case by induction on the order of $G$.

Because $G$ is solvable, it contains a proper normal subgroup $H$ such that $G / H$ is cyclic. Because $H$ has order less than that of $G$, and the pair $(H, M)$ satisfies the hypotheses of the theorem, $H^{r}(H, M)=0$ for all $r$. Therefore (see 1.34), we have exact sequences


\begin{equation*}
0 \rightarrow H^{r}\left(G / H, M^{H}\right) \rightarrow H^{r}(G, M) \rightarrow H^{r}(H, M) \tag{*}
\end{equation*}


for all $r \geq 1$. Because $H^{1}(G, M)=0=H^{2}(G, M), H^{1}\left(G / H, M^{H}\right)=H^{2}\left(G / H, M^{H}\right)=$ 0 , and because $G / H$ is cyclic, this implies that $H^{r}\left(G / H, M^{H}\right)=0$ for all $r$. Therefore, the sequences $(*)$ show that $H^{r}(G, M)=0$ for all $r>0$. We next show that $H^{0}(G, M)=0$. Let $x \in M^{G}$. Because $H^{0}\left(G / H, M^{H}\right)=0$, there exists a $y \in M^{H}$ such that $\operatorname{Nm}_{G / H}(y)=x$, and because $H^{0}(H, M)=0$, there exists a $z \in M$ such that $\operatorname{Nm}_{H}(z)=y$. Now

$$
\operatorname{Nm}_{G}(z)=\left(\operatorname{Nm}_{G / H} \circ \operatorname{Nm}_{H}\right)(z)=x .
$$

Thus, we now know that $H^{r}(G, M)=0$ for all $r \geq 0$.\\
To proceed further, we use the exact sequence

$$
0 \rightarrow M^{\prime} \rightarrow \mathbb{Z}[G] \otimes_{\mathbb{Z}} M \rightarrow M \rightarrow 0
$$

obtained by tensoring the augmentation sequence (19) with $M$. Recall that

$$
\mathbb{Z}[G] \otimes_{\mathbb{Z}} M \stackrel{1.4}{\simeq} \mathbb{Z}[G] \otimes_{\mathbb{Z}} M_{0} \stackrel{1.3}{\simeq} \operatorname{Ind}^{G}\left(M_{0}\right)
$$

is induced, and so $H^{r}\left(H, \mathbb{Z}[G] \otimes_{\mathbb{Z}} M\right)=0$ for all $r$ and all subgroups $H$ of $G$. Therefore,

$$
H^{r}(H, M) \simeq H^{r+1}\left(H, M^{\prime}\right)
$$

for all $r$ and all $H$. In particular, $M^{\prime}$ satisfies the hypotheses of the theorem, and so (by what we have proved) $H^{r}\left(G, M^{\prime}\right)=0$ for $r \geq 0$. In particular,

$$
0=H^{0}\left(G, M^{\prime}\right)=H^{-1}(G, M) .
$$

The argument, when repeated, gives that $H^{-2}(G, M)=0$, etc.. This proves the theorem when $G$ is solvable.

Now consider the case of an arbitrary finite group $G$. If $G$ and $M$ satisfy the hypotheses of the theorem, so also do $G_{p}$ and $M$, where $G_{p}$ is a Sylow $p$-subgroup. Therefore $H^{r}\left(G_{p}, M\right)=0$ for all $r$ and $p$, and so (see 1.33), the $p$-primary component of $H^{r}(G, M)$ is zero for all $r$ and $p$. This implies that $H^{r}(G, M)=0$ for all $r$. \(\square\)

Theorem 3.11 (Tate's theorem) Let $G$ be a finite group and let $C$ be a $G$-module. Suppose that for all subgroups $H$ of $G$ (including $H=G$ ),

\begin{itemize}
  \item[(a)] $H^{1}(H, C)=0$, and
  \item[(b)] $H^{2}(H, C)$ is a cyclic group of order equal to $(H: 1)$.
\end{itemize}

Then, for all $r$, there is an isomorphism

$$
H_{T}^{r}(G, \mathbb{Z}) \rightarrow H_{T}^{r+2}(G, C)
$$

depending only on the choice of a generator for $H^{2}(G, C)$.\\
Proof. Choose a generator $\gamma$ for $H^{2}(G, C)$. Because $\operatorname{Cor} \circ \operatorname{Res}=(G: H), \operatorname{Res}(\gamma)$ generates $H^{2}(H, C)$ for any subgroup $H$ of $G$.

Let $\varphi$ be a cocycle representing $\gamma$. Define $C(\varphi)$ to be the direct sum of $C$ with the free abelian group having as basis symbols $x_{\sigma}$, one for each $\sigma \in G, \sigma \neq 1$, and extend the action of $G$ on $C$ to an action on $C(\varphi)$ by setting

$$
\sigma x_{\tau}=x_{\sigma \tau}-x_{\sigma}+\varphi(\sigma, \tau) .
$$

The symbol " $x_{1}$ " is to be interpreted as $\varphi(1,1)$. This does define an action of $G$ on $C(\varphi)$ because

$$
\rho \sigma x_{\tau}=x_{\rho \sigma \tau}-x_{\rho \sigma}+\varphi(\rho \sigma, \tau),
$$

whereas

$$
\begin{aligned}
\rho\left(\sigma x_{\tau}\right) & =\rho\left(x_{\sigma \tau}-x_{\sigma}+\varphi(\sigma, \tau)\right) \\
& =x_{\rho \sigma \tau}-x_{\rho}+\varphi(\rho, \sigma \tau)-\left(x_{\rho \sigma}-x_{\rho}+\varphi(\rho, \sigma)\right)+\rho \varphi(\sigma, \tau) .
\end{aligned}
$$

These agree because $\varphi$ satisfies the cocycle condition

$$
\rho \varphi(\sigma, \tau)+\varphi(\rho, \sigma \tau)=\varphi(\rho \sigma, \tau)+\varphi(\rho, \sigma) .
$$

Note that $\varphi$ is the coboundary of the 1 -cochain $\sigma \mapsto x_{\sigma}$, and so $\gamma$ maps to zero in $H^{2}(G, C(\varphi))$. For this reason, $C(\varphi)$ is called the splitting module for $\gamma$.

We shall first show that the hypotheses imply that

$$
H^{1}(H, C(\varphi))=0=H^{2}(H, C(\varphi))
$$

for all subgroups $H$ of $G$.\\
Recall that we have an exact sequence

$$
0 \rightarrow I_{G} \rightarrow \mathbb{Z}[G] \rightarrow \mathbb{Z} \rightarrow 0,
$$

where $I_{G}$ is the free abelian group with basis the elements $\sigma-1, \sigma \in G, \sigma \neq 1$ Because $\mathbb{Z}[G]$ is induced, $H^{r}(H, \mathbb{Z}[G])=0$ for all $r$, and so

$$
\begin{array}{ll}
H^{1}\left(H, I_{G}\right) & \simeq H^{0}(H, \mathbb{Z}) \simeq \mathbb{Z} /(H: 1) \mathbb{Z}, \\
H^{2}\left(H, I_{G}\right) & \simeq H^{1}(H, \mathbb{Z})=0
\end{array}
$$

Define $\alpha$ to be the additive map $C(\varphi) \rightarrow \mathbb{Z}[G]$ such that

$$
\begin{aligned}
\alpha(c) & =0 \text { for all } c \in C \\
\alpha\left(x_{\sigma}\right) & =\sigma-1 .
\end{aligned}
$$

Clearly,

$$
0 \rightarrow C \rightarrow C(\varphi) \xrightarrow{\alpha} I_{G} \rightarrow 0
$$

is an exact sequence of $G$-modules. Its cohomology sequence reads

$$
0 \rightarrow H^{1}(H, C(\varphi)) \rightarrow H^{1}\left(H, I_{G}\right) \rightarrow H^{2}(H, C) \xrightarrow{0} H^{2}(H, C(\varphi)) \rightarrow 0
$$

The zeros at the ends use that $H^{1}(H, C)=0$ and $H^{2}\left(H, I_{G}\right)=0$. The map $H^{2}(H, C) \rightarrow$ $H^{2}(H, C(\varphi))$ is zero because $H^{2}(H, C)$ is generated by $\operatorname{Res}(\gamma)$, and this maps to the restriction of the image of $\gamma$ in $H^{2}(G, C(\varphi))$, which is zero. Therefore, $H^{1}\left(H, I_{G}\right) \rightarrow$ $H^{2}(H, C)$ is onto, and hence is an isomorphism (because the two groups have the same order). Its kernel and cokernel, namely, $H^{1}(H, C(\varphi))$ and $H^{2}(H, C(\varphi))$, are therefore both zero

We deduce from the Theorem 3.10 that $H^{r}(H, C(\varphi))=0$ for all $r$. On splicing the two short exact sequences together, we obtain an exact sequence


\begin{equation*}
0 \rightarrow C \rightarrow C(\varphi) \rightarrow \mathbb{Z}[G] \rightarrow \mathbb{Z} \rightarrow 0 \tag{24}
\end{equation*}


with the property that $H^{r}(G, C(\varphi))=0=H^{r}(G, \mathbb{Z}[G])$ for all $r$. Therefore, the double boundary map is an isomorphism (1.13)

$$
H^{r}(G, \mathbb{Z}) \rightarrow H^{r+2}(G, C) .
$$ \(\square\)

REMARK 3.12 If $M$ is a $G$-module such that $\operatorname{Tor}_{1}^{\mathbb{Z}}(M, C)=0$, for example, if either $M$ or $C$ is torsion-free as a $\mathbb{Z}$-module, then one can tensor the above 4-term sequence with $M$ and obtain isomorphisms

$$
H_{T}^{r}(G, M) \rightarrow H_{T}^{r+2}(G, M \otimes C) .
$$

REMARK 3.13 The map $H_{T}^{r}(G, \mathbb{Z}) \rightarrow H_{T}^{r+2}(G, C)$ is cup-product with the chosen element $\gamma \in H^{2}(G, C)$.

Example 3.14 Let $K$ be a local field. We shall prove that for any finite Galois extension $L$ of $K$ with Galois group $G, H^{2}\left(G, L^{\times}\right)$is cyclic of order $[L: K]$ with a canonical generator $u_{L / K}$. Since we know that $H^{1}\left(G, L^{\times}\right)=0$ (Hilbert's Theorem 90), Tate's theorem shows that cup-product with $u_{L / K}$ is an isomorphism

$$
G^{\mathrm{ab}}=H_{T}^{-2}(G, \mathbb{Z}) \rightarrow H_{T}^{0}\left(G, L^{\times}\right)=K^{\times} / \operatorname{Nm} L^{\times} .
$$

The inverse isomorphism $K^{\times} / \operatorname{Nm} L^{\times} \xrightarrow{\simeq} G^{\text {ab }}$ is the local Artin map. The global Artin map can be obtained by a similar argument.

Notes The book Weiss 1969 gives a complete and very detailed account of the basic theory of Tate cohomology groups. For a recent paper on the Tate sequence (24), its generalizations, and associated boundary maps, see Buckingham, Acta Arith. 149 (2011), no. 4, 383-402.

\section*{4 The Cohomology of Profinite Groups}
Direct limits\\
A partially ordered set $(I, \leq)$ is said to be directed if for any two elements $i$ and $j$ of $I$, there exists a $k$ such that $i, j \leq k$. Suppose that for every element $i$ of a directed set $(I, \leq)$ we have a set $A_{i}$, and for every inequality $i \leq j$ we have a map $\alpha_{j i}: A_{i} \rightarrow A_{j}$. If

\begin{itemize}
  \item[(a)] $\alpha_{i i}=\mathrm{id}$ for all $i \in I$, and
  \item[(b)] $\alpha_{k j} \circ \alpha_{j i}=\alpha_{k i}$ whenever $i \leq j \leq k$,
\end{itemize}

then the family $\left(A_{i}, \alpha_{j i}\right)$ is called a direct system. On the disjoint union $\coprod A_{i}$ of the $A_{i}$, introduce the equivalence relation under which $a_{i} \in A_{i}$ is related to $a_{j} \in A_{j}$ if and only if $\alpha_{k i}\left(a_{i}\right)=\alpha_{k j}\left(a_{j}\right)$ for some $k \geq i, j$. The corresponding quotient set is called the direct limit of the $A_{i}$ (relative to the $\alpha_{j i}$ ):

$$
A=\underset{\longrightarrow}{\lim } A_{i} .
$$

There is for each $i$ a canonical map

$$
\alpha_{i}: A_{i} \rightarrow A,
$$

possessing the following properties:

\begin{itemize}
  \item[(a)] $\alpha_{i}=\alpha_{j} \circ \alpha_{j i}$ for $j \geq i$;
  \item[(b)] $\alpha_{i}\left(a_{i}\right)=\alpha_{j}\left(a_{j}\right) \Longleftrightarrow \alpha_{k i}\left(a_{i}\right)=\alpha_{k j}\left(a_{j}\right)$ for some $k \geq i, j$;
  \item[(c)] $A=\bigcup \alpha_{i}\left(A_{i}\right)$.
\end{itemize}

The system ( $A, \alpha_{i}$ ) has the following universal property: let $T$ be a set and let ( $\beta_{i}$ ), $\beta_{i}: A_{i} \rightarrow T$, be a family of maps such that $\beta_{i}=\beta_{j} \circ \alpha_{j i}$ for $i \leq j$; then there exists a unique map $\beta: A \rightarrow T$ such that $\beta_{i}=\beta \circ \alpha_{i}$ for all $i$.

If the $A_{i}$ are abelian groups and the $\alpha_{i j}$ are homomorphisms, then $A$ has a unique structure of an abelian group for which the $\alpha_{i}$ are homomorphisms.

Lemma 4.1 For any direct system of exact sequences

$$
A_{i} \rightarrow B_{i} \rightarrow C_{i},
$$

the sequence

$$
\xrightarrow{\lim } A_{i} \rightarrow \underset{\longrightarrow}{\lim } B_{i} \rightarrow \underset{\longrightarrow}{\lim } C_{i}
$$

is again exact. Therefore the formation of direct limits commutes with passage to cohomology in complexes.

Proof. Exercise for the reader. \(\square\)

Profinite groups\\
Let $G$ be a profinite group. This means that $G$ is a compact ${ }^{5}$ topological group for which the open normal subgroups form a fundamental system of neighbourhoods of 1. Note that every open subgroup is of finite index (because its cosets cover $G$ ). For example, a finite group with the discrete topology is profinite, and every discrete profinite group is finite. A Galois

\footnotetext{${ }^{5}$ Following Bourbaki, I require compact spaces to be Hausdorff.
}
group $G=\operatorname{Gal}(L / K)$ is a profinite group - the open subgroups are exactly those fixing a finite extension of $K$ contained in $L$ - and every profinite group occurs as a Galois group. For a profinite group, we use the topology to modify our notion of cohomology group.

First, we consider only those $G$-modules for which the map

$$
G \times M \rightarrow M
$$

is continuous when $M$ is endowed with the discrete topology, i.e., the topology in which every subset is open. Equivalent conditions:

\begin{itemize}
  \item[-] $M=\bigcup M^{H}, H$ runs through the open subgroups of $G$;
  \item[] ◇ the stabilizer in $G$ of any element of $M$ is open.
\end{itemize}

A module satisfying these conditions is called a discrete $G$-module.\\
The discrete $G$-modules form an abelian category with enough injectives, and so we can define cohomology groups $H_{\text {cts }}^{r}(G, M), r \geq 0$, by taking injective resolutions, just as before. Moreover, the groups can be calculated using continuous cocycles: for $r \geq 0$, let $C_{\text {cts }}^{r}(G, M)$ be the group of continuous maps $G^{r} \rightarrow M$, and define $d^{r}: C_{\text {cts }}^{r}(G, M) \rightarrow$ $C_{\text {cts }}^{r+1}(G, M)$ as before; then

$$
H_{\mathrm{cts}}^{r}(G, M)=\frac{Z_{\mathrm{cts}}^{r}(G, M)}{B_{\mathrm{cts}}^{r}(G, M)},
$$

where $Z_{\text {cts }}^{r}(G, M)=\operatorname{Ker}\left(d^{r}\right)$ and $B_{\text {cts }}^{r}(G, M)=\operatorname{Im}\left(d^{r-1}\right)$.\\
Let $\varphi: G^{r} \rightarrow M$ be a continuous $r$-cochain. Then $\varphi\left(G^{r}\right)$ is compact (because $G^{r}$ is compact) and discrete (because $M$ is discrete), and so it is finite, and hence it is contained in $M^{H_{0}}$ for some open normal subgroup $H_{0}$ of $G$. The inverse image $\varphi^{-1}(m)$ of each point $m$ of $\varphi\left(G^{r}\right)$ is open, and so contains a translate of $H(m)^{r}$ for some open normal subgroup $H(m)$ of $G$. Let $H_{1}=\bigcap_{m \in \varphi\left(G^{r}\right)} H(m)$. ${ }^{6}$ This is again an open subgroup of $G$, and $\varphi$ factors through $\left(G / H_{1}\right)^{r}$. Let $H=H_{0} \cap H_{1}$. Then $\varphi$ arises by inflation from an $r$-cocycle on $G / H$ with values in $M^{H}$. In other words, the map

$$
\xrightarrow{\lim _{\longrightarrow}} C^{r}\left(G / H, M^{H}\right) \rightarrow C_{\mathrm{cts}}^{r}(G, M)
$$

is surjective. It is not difficult to see that it is also injective. Because passage to the direct limit commutes with the formation of kernels and cokernels, and hence with the formation of cohomology, we obtain the following proposition:

Proposition 4.2 The maps $\operatorname{Inf}: H^{r}\left(G / H, M^{H}\right) \rightarrow H_{\text {cts }}^{r}(G, M)$ realize $H_{\text {cts }}^{r}(G, M)$ as the direct limit of the groups $H^{r}\left(G / H, M^{H}\right)$ as $H$ runs through the open normal subgroups $H$ of $G$ :

$$
\lim _{\longrightarrow} H^{r}\left(G / H, M^{H}\right)=H_{c t s}^{r}(G, M) .
$$

Explicitly, the statement means that each element of $H_{\text {cts }}^{r}(G, M)$ arises by inflation from some group $H^{r}\left(G / H, M^{H}\right)$ and if $a \in H^{r}\left(G / H, M^{H}\right)$ and $a^{\prime} \in H^{r}\left(G / H^{\prime}, M^{H^{\prime}}\right)$ map to the same element in $H_{\text {cts }}^{r}(G, M)$, then they map to the same element in $H^{r}\left(G / H^{\prime \prime}, M^{H^{\prime \prime}}\right)$ for some open subgroup $H^{\prime \prime} \subset H \cap H^{\prime}$.

\footnotetext{${ }^{6}$ This group may not be small enough for $\varphi$ to factor through $\left(G / H_{1}\right)^{r}$. Instead, for any point of $G^{r}$, take a translate of an open normal subgroup containing the point and fully contained in a fibre of $\varphi$. In this way, we get an open covering of $G^{r}$, which may be replaced by a finite subcovering because $G^{r}$ is compact. Now let $H_{1}$ be the intesection of all the subgroups whose translates occur in this subcovering.
}Corollary 4.3 For every profinite group $G$ and discrete $G$-module $M, H^{r}(G, M)$ is a torsion group for all $r>0$.

Proof. Each of the groups $H^{r}\left(G / H, M^{H}\right), r>0$, is torsion (see1.31). \(\square\)

Most of the theory concerning the cohomology groups $H^{r}(G, M)$ for $r \geq 0$ continues to hold for the groups defined by continuous cochains. For example, if $H$ is a closed subgroup of $G$, there are maps Inf, Res, and Cor (the last requires $H$ to be of finite index), and there are cup-product maps.

In future, all cohomology groups will be defined using continuous cochains (and the subscript cts will be dropped). In practice, this will mean that either $G$ is an infinite profinite group, and it matters that we take continuous cochains, or $G$ is finite, in which case it doesn't (and the groups are defined for all $r \in \mathbb{Z}$ ).

Proposition 4.4 Let $G$ be a profinite group, and let $M$ be a discrete $G$-module. If $M=$ $\xrightarrow{\lim } M_{i}$, where $M_{i} \subset M$, then $H^{r}(G, M)=\underset{\longrightarrow}{\lim } H^{r}\left(G, M_{i}\right)$.

Proof. Because $G$ is compact and $M$ is discrete, the image of any $r$-cochain $f: G^{r} \rightarrow M$ is finite. Since the $M_{i}$ form a directed system of submodules of $M$ (i.e., given $M_{i}$ and $M_{j}$, there is an $M_{k}$ containing both of them) and $M=\bigcup M_{i}$, every finite subset of $M$ is contained in an $M_{i}$. It follows that $C^{r}(G, M)=\underset{\longrightarrow}{\lim } C^{r}\left(G, M_{i}\right)$, and so Lemma 4.1 shows that

$$
H^{r}\left(C^{\bullet}(G, M)\right)=\underset{\longrightarrow}{\lim } H^{r}\left(C^{\bullet}\left(G, M_{i}\right)\right) .
$$ \(\square\)

Aside 4.5 (For the experts.) Let $G$ be a profinite group, let $\mathrm{M}(G)$ be the category of all $G$-modules, and $\mathrm{C}(G)$ the category of discrete $G$-modules. Then $\mathrm{C}(G)$ is a full subcategory of $\mathrm{M}(G)$. Moreover, there is a functor

$$
M \mapsto M^{*} \stackrel{\text { def }}{=} \bigcup_{H \text { open in } G} M^{H}: \mathrm{M}(G) \rightarrow \mathrm{C}(G) .
$$

Clearly,

$$
\operatorname{Hom}_{G}\left(M, N^{*}\right)=\operatorname{Hom}_{G}(M, N)
$$

for $M$ a discrete $G$-module. The inclusion functor $i: \mathrm{C}(G) \rightarrow \mathrm{M}(G)$ is exact, but doesn't preserve injectives-hence $H^{r}(G, M) \neq H^{r}(G, i M)$ in general. On the other hand, $M \mapsto M^{*}$ preserves injectives, but is only left exact-hence $\mathrm{C}(M)$ has enough injectives. Again $H^{r}(G, M) \neq$ $H^{r}\left(G, M^{*}\right)$ (however, there is a spectral sequence ...).

Notes In the mid-1930s, Hurewicz showed that the homology groups of an "aspherical space" $X$ depend only on the fundamental group $\pi$ of the space. Thus one could think of the homology groups $H_{r}(X, \mathbb{Z})$ of the space as being the homology groups $H_{r}(\pi, \mathbb{Z})$ of the group $\pi$. It was only in the mid-1940s that Hopf, Eckmann, Eilenberg, MacLane, Freudenthal and others gave purely algebraic definitions of the homology and cohomology groups of a group $G$. It was then found that $H^{1}$ coincided with the group of crossed homomorphisms modulo principal crossed homomorphisms, and $H^{2}$ with the group of equivalence classes of "factor sets", which had been introduced much earlier (e.g., I. Schur, Über die Darstellung der endlichen. . . , 1904; O. Schreier, Über die Erweiterungen von Gruppen, 1926; R. Brauer, Über Zusammenhänge... , 1926). For more on the history, see Mac Lane 1978.

Our proof of Tate's theorem follows Tate's original proof (Tate 1952). At that time, there was no published account of the Tate groups, and so Tate proved his theorem only for $r \geq 0$, but ended with an enigmatic promise to extend the result to negative $r$ and to recover the Artin map. The construction of the Tate cohomology groups was first published in Cartan and Eilenberg 1956 following Tate's ideas.

A good source for the material in the first three sections is Serre 1962, Part 3. There is a somewhat abbreviated version of the same material in Atiyah and Wall 1967. See also Iyanaga 1975, Chapter I, and Weiss 1969. For the cohomology of profinite groups, see Serre 1964 and Shatz 1972.

For a discussion of the origins of cohomology in class field theory, see sx857593.

\section*{A Appendix: Some Homological Algebra}
\begin{displayquote}
Ça me semble extrêmement plaisant de ficher comme ça beaucoup de choses, pas drôles quand on les prend séparément, sous le grand chapeau des foncteurs dérivés.\\
Grothendieck, Letter to Serre, 18.2.1955. ${ }^{7}$
\end{displayquote}

Some exact sequences ${ }^{8}$\\
Lemma A. 1 (The Extended Snake Lemma) The exact commutative diagram in blue gives rise to the exact sequence in red:\\
\includegraphics[max width=\textwidth, alt={}, center]{e7cfc28e-5737-4710-b4d2-10a47512e540-33_556_1376_1246_364}\\
Proof. Except for the first and last terms, this is standard. A small diagram chase shows that $\operatorname{Ker} f \subset \operatorname{Ker}(a)$, from which exactness at $\operatorname{Ker} a$ follows. The proof of exactness at Coker $c$ is similarly straightforward. \(\square\)

Lemma A. 2 (Kernel-Cokernel Lemma) Every pair of homomorphisms

$$
A \xrightarrow{f} B \xrightarrow{g} C
$$

of abelian groups gives rise to an exact sequence

$$
0 \rightarrow \operatorname{Ker} f \rightarrow \operatorname{Ker} g \circ f \xrightarrow{f} \operatorname{Ker} g \rightarrow \operatorname{Coker} f \rightarrow \operatorname{Coker} g \circ f \rightarrow \operatorname{Coker} g \rightarrow 0 .
$$

\footnotetext{${ }^{7}$ To me it seems extremely pleasant to stick all sorts of things, which are not much fun when considered separately, under the big hat of derived functors\\
${ }^{8}$ Based on F. Lemmermeyer, The Snake Lemma.
}Proof. Apply the extended snake lemma to\\
\includegraphics[max width=\textwidth, alt={}]{e7cfc28e-5737-4710-b4d2-10a47512e540-34_193_768_424_639} \(\square\)

The language of category theory\\
A category C consists of a nonempty class ob(C) of objects, a set $\operatorname{Hom}(A, B)$ for each pair of objects $A, B$ (called the set of morphisms from $A$ to $B$ ), and a map

$$
(\alpha, \beta) \mapsto \beta \circ \alpha: \operatorname{Hom}(A, B) \times \operatorname{Hom}(B, C) \rightarrow \operatorname{Hom}(A, C)
$$

for each triple of objects $A, B, C$, satisfying the following conditions:

\begin{itemize}
  \item[(a)] composition of morphisms is associative;
  \item[(b)] for each object $A, \operatorname{Hom}(A, A)$ has an element $\mathrm{id}_{A}$ that is a left and right identity for composition.
\end{itemize}

It is to be understood that the sets $\operatorname{Hom}(A, B)$ are disjoint, so that a morphism determines its source and target.

A covariant functor $F: \mathrm{C} \rightarrow \mathrm{D}$ is a "map" that attaches to each object $A$ of C an object $F(A)$ of D and to each morphism $\alpha: A \rightarrow B$ a morphism $F(\alpha): F(A) \rightarrow F(B)$ such that $F(\alpha \circ \beta)=F(\alpha) \circ F(\beta)$ and $F\left(\operatorname{id}_{A}\right)=\operatorname{id}_{F(A)}$.

A functor $F: \mathrm{C} \rightarrow \mathrm{D}$ is left adjoint to the functor $G: \mathrm{D} \rightarrow \mathrm{C}$ if there is a natural isomorphism

$$
\operatorname{Hom}_{D}(F(A), B) \simeq \operatorname{Hom}_{C}(A, G(B)) .
$$

If the sets $\operatorname{Hom}(A, B)$ are endowed with the structures of abelian groups in such a way that the composition maps are bi-additive, and every finite collection of objects in C has a direct sum, then C (together with the structures) is called an additive category. To say that $A$ and $B$ admit a direct sum means that there is an object $A \oplus B$ in C and maps $i_{A}: A \rightarrow A \oplus B$, $i_{B}: B \rightarrow A \oplus B, p_{A}: A \oplus B \rightarrow A, p_{B}: A \oplus B \rightarrow B$ such that:

$$
\begin{array}{cc}
p_{A} \circ i_{A}=\operatorname{id}_{A} & p_{B} \circ i_{B}=\operatorname{id}_{B} \\
p_{A} \circ i_{B}=0 & p_{B} \circ i_{A}=0 \\
& i_{A} p_{A}+i_{B} p_{B}=1_{A \oplus B} .
\end{array}
$$

Let C be an additive category. A sequence

$$
0 \rightarrow A \rightarrow B \xrightarrow{\alpha} C
$$

is exact if the sequence of abelian groups

$$
0 \rightarrow \operatorname{Hom}(T, A) \rightarrow \operatorname{Hom}(T, B) \rightarrow \operatorname{Hom}(T, C)
$$

is exact for all objects $T$. A sequence

$$
A \xrightarrow{\beta} B \rightarrow C \rightarrow 0
$$

is exact if the sequence of abelian groups

$$
0 \rightarrow \operatorname{Hom}(C, T) \rightarrow \operatorname{Hom}(B, T) \rightarrow \operatorname{Hom}(A, T)
$$

is exact for all objects $T$. When the first sequence is exact, $A$ is called the kernel of $\alpha$, and when the second is exact, $C$ is called the cokernel of $\beta$.

A morphism $A \rightarrow B$ is said to be injective (or a monomorphism) if $0 \rightarrow A \rightarrow B$ is exact. Similarly, it is surjective (or an epimorphism) if $A \rightarrow B \rightarrow 0$ is exact.

Let C be an additive category in which every morphism has both a kernel and a cokernel. Let $\alpha: A \rightarrow B$ be morphism. The kernel of the cokernel $B \rightarrow C$ of $\alpha$ is called the image of $\alpha$, and the cokernel of the kernel of $\alpha$ is called the coimage of $\alpha$. There is a canonical map from the coimage of $\alpha$ to the image of $\alpha$, and if this is always an isomorphism, then C is called an abelian category. ${ }^{9}$

Functors between additive categories will be assumed to be additive, i.e., such that the maps $\operatorname{Hom}(A, B) \rightarrow \operatorname{Hom}(F(A), F(B))$ are homomorphisms of abelian groups. Such a functor is said to be exact if it maps exact sequences to exact sequences.

For example, for every ring $R$, the category of $R$-modules is an abelian category, and, for every topological space $X$, the category of sheaves of abelian groups on $X$ is an abelian category.

In the remainder of this section, C will be an abelian category. The reader will lose little (so far as this course is concerned) by taking C to be the category of modules over a ring, for example, the category of modules over a group ring $\mathbb{Z}[G]$.

\section*{Injective objects}
Let C be an abelian category. An object $I$ of C is injective if $\operatorname{Hom}(\cdot, I)$ is an exact functor, i.e., if

$$
0 \rightarrow A \rightarrow B \rightarrow C \rightarrow 0
$$

exact in C implies that

$$
0 \rightarrow \operatorname{Hom}(C, I) \rightarrow \operatorname{Hom}(B, I) \rightarrow \operatorname{Hom}(A, I) \rightarrow 0
$$

is exact. The last sequence is automatically exact except at\\
\includegraphics[max width=\textwidth, alt={}, center]{e7cfc28e-5737-4710-b4d2-10a47512e540-35_206_185_1819_1451}\\
$\operatorname{Hom}(A, I)$, and so to say that $I$ is injective means that, whenever $A$ is a subobject of $B$, every homomorphism $A \rightarrow I$ extends to $B$.

An abelian category C is said to have enough injectives if every object admits an injective homomorphism into an injective object.

Proposition A. 3 A module $M$ over a principal ideal domain $R$ is injective if and only if it is divisible, i.e., the map $x \mapsto r x$ is surjective for all $r \in R, r \neq 0$.

Proof. injective ⟹ divisible: Let $m \in M$ and let $r$ be a nonzero element of $R$. The map $x \mapsto r x: R \rightarrow R$ is injective, and every extension of $x \mapsto x m: R \rightarrow M$ to $R$ will send 1 to an element $m^{\prime}$ such that $r m^{\prime}=m$.\\
divisible ⇒ injective: Suppose that $M$ is divisible, and consider a homomorphism $\alpha: A \rightarrow M$, where $A$ is a submodule of $B$. On applying Zorn's lemma to the set of pairs

\footnotetext{${ }^{9}$ If you are mystified by this paragraph, you may find sx2263197 helpful.
}
( $A^{\prime}, \alpha^{\prime}$ ), where $A^{\prime}$ is a submodule of $B$ containing $A$ and $\alpha^{\prime}$ extends $\alpha$ to $A^{\prime}$, we obtain a maximal such pair $\left(A_{1}, \alpha_{1}\right)$. If $A_{1} \neq B$, then there exists a $b \in B \backslash A_{1}$, and we define $I=\left\{r \in R \mid r b \in A_{1}\right\}$. Because $M$ is divisible, the map $r \mapsto \alpha_{1}(r b): I \rightarrow M$ extends ${ }^{10}$ to $R$, but this implies that $\alpha_{1}$ extends to $A_{1}+R b$, which contradicts the maximality of $\left(A_{1}, \alpha_{1}\right)$. \(\square\)

Therefore, when $R$ is a principal ideal domain, every quotient of an injective module is injective, for example, $K / R$, where $K$ is the field of fractions of $R$. In particular, $\mathbb{Q} / \mathbb{Z}$ is an injective abelian group.

Proposition A. 4 The category of modules over a principal ideal domain $R$ has enough injectives.

Proof. Write $M$ as a quotient $M \approx F / N$ of a free $R$-module $F$, and consider the exact commutative diagram,\\
\includegraphics[max width=\textwidth, alt={}, center]{e7cfc28e-5737-4710-b4d2-10a47512e540-36_173_951_1017_556}\\
Because $F \rightarrow F \otimes_{R} K$ is an injective map and $F \otimes_{R} K$ is an injective object, the map $M \rightarrow\left(F \otimes_{R} K\right) / N$ is an injective homomorphism of $M$ into an injective $R$-module. \(\square\)

For abelian groups, the proposition can also be proved as follows: for an abelian group $M$, let $M^{\vee}=\operatorname{Hom}(M, \mathbb{Q} / \mathbb{Z})$; choose a free abelian group $F$ mapping onto $M^{\vee}, F \rightarrow$ $M^{\vee}$; then $M \hookrightarrow M^{\vee \vee} \hookrightarrow F^{\vee}$, and $F^{\vee}$ is injective.

Proposition A. 5 A functor that admits an exact left adjoint preserves injectives.\\
Proof. Let $F^{\prime}$ be an exact left adjoint to the functor $F: \mathrm{C} \rightarrow \mathrm{D}$. For every injective object $I$ in C, the functor $\operatorname{Hom}_{\mathrm{D}}(\cdot, F(I))$ is isomorphic to the functor $\operatorname{Hom}_{\mathrm{C}}\left(F^{\prime}(\cdot), I\right)$, which is exact because it is the composite of two exact functors, namely, $F^{\prime}$ and $\operatorname{Hom}_{\mathrm{C}}(\cdot, I)$. \(\square\)

\section*{Right derived functors}
Let C be an abelian category with enough injectives, and let $F: \mathrm{C} \rightarrow \mathrm{D}$ be a left exact functor from C to a second abelian category. Thus, a short exact sequence

$$
0 \rightarrow M^{\prime} \rightarrow M \rightarrow M^{\prime \prime} \rightarrow 0
$$

in C gives rise to an exact sequence

$$
0 \rightarrow F\left(M^{\prime}\right) \rightarrow F(M) \rightarrow F\left(M^{\prime \prime}\right)
$$

in D. The theory of derived functors provides a natural extension of this last sequence to a long exact sequence.

Let $M$ be an object of C. A resolution of $M$ is a long exact sequence

$$
0 \rightarrow M \rightarrow I^{0} \xrightarrow{d^{0}} I^{1} \rightarrow \cdots \rightarrow I^{r} \xrightarrow{d^{r}} I^{r+1} \rightarrow \cdots .
$$

If the $I^{r}$ 's are injective objects of C, then it is called an injective resolution. We sometimes denote this complex by $M \rightarrow I^{\bullet}$.

\footnotetext{${ }^{10}$ For the moment, this is left as an exercise for the reader.
}Lemma A. 6 An injective resolution $M \rightarrow I^{\bullet}$ of $M$ exists, and if $M \rightarrow J^{\bullet}$ is a second injective resolution, then there exists a homomorphism from $M \rightarrow I^{\bullet}$ to $M \rightarrow J^{\bullet}$, i.e., there exists a commutative diagram,\\
\includegraphics[max width=\textwidth, alt={}, center]{e7cfc28e-5737-4710-b4d2-10a47512e540-37_180_704_514_681}

Proof. By assumption, there exists an injective morphism

$$
0 \rightarrow M \rightarrow I^{0}
$$

with $I^{0}$ injective. Let $B^{1}$ be the cokernel of the map. Again, there exists an injective morphism

$$
0 \rightarrow B^{1} \rightarrow I^{1}
$$

with $I^{1}$ injective. Now

$$
0 \rightarrow M \rightarrow I^{0} \rightarrow I^{1}
$$

is exact. Let $B^{2}=\operatorname{Coker}\left(B^{1} \rightarrow I^{1}\right)$, and continue in this fashion.\\
Similarly, a morphism of resolutions can be constructed step by step, using the definition of an injective object. \(\square\)

On applying a left exact functor $F$ to an injective resolution $M \rightarrow I^{\bullet}$ of a object $M$, we obtain a complex

$$
F\left(I^{\bullet}\right): F\left(I^{0}\right) \rightarrow F\left(I^{1}\right) \rightarrow \cdots \rightarrow F\left(I^{r}\right) \xrightarrow{F\left(d^{r}\right)} F\left(I^{r+1}\right) \rightarrow \cdots
$$

which may no longer be exact, and so the cohomology groups

$$
H^{r}\left(F\left(I^{\bullet}\right)\right) \stackrel{\text { def }}{=} \frac{\operatorname{Ker}\left(F\left(d^{r}\right)\right)}{\operatorname{Im}\left(F\left(d^{r-1}\right)\right)}
$$

may be nonzero.\\
Remark A. 7 Because $F$ is left exact, the sequence

$$
0 \rightarrow F(M) \rightarrow F\left(I^{0}\right) \xrightarrow{d^{0}} F\left(I^{1}\right)
$$

is exact. Therefore,

$$
H^{0}\left(F\left(I^{\bullet}\right)\right) \stackrel{\text { def }}{=} \operatorname{Ker}\left(d^{0}\right)=F(M) .
$$

Proposition A. 8 Let $M \rightarrow I^{\bullet}$ and $N \rightarrow J^{\bullet}$ be injective resolutions of objects $M$ and $N$ of C. Every morphism $\alpha: M \rightarrow N$ extends to a map of complexes\\
\includegraphics[max width=\textwidth, alt={}, center]{e7cfc28e-5737-4710-b4d2-10a47512e540-37_170_241_2205_915}\\
For every left exact functor $F$, the morphism

$$
H^{r}\left(F\left(\alpha^{\bullet}\right)\right): H^{r}\left(F\left(I^{\bullet}\right)\right) \rightarrow H^{r}\left(F\left(J^{\bullet}\right)\right)
$$

is independent of the choice of $\alpha^{\bullet}$.

We discuss the proof below.\\
The proposition, applied to the identity morphism $M \rightarrow M$, implies that the objects $H^{r}\left(F\left(I^{\bullet}\right)\right)$ are well-defined up to a well-defined isomorphism: given two injective resolutions $M \rightarrow I^{\bullet}$ and $M \rightarrow J^{\bullet}$ of $M$, there exists a morphism $\alpha^{\bullet}: I^{\bullet} \rightarrow J^{\bullet}$ extending the identity map on $M$, and the maps $H^{r}\left(F\left(I^{\bullet}\right)\right) \rightarrow H^{r}\left(F\left(J^{\bullet}\right)\right)$ it defines are isomorphisms independent of the choice of $\alpha^{\bullet}$. Now, for each object $M$ of C, we choose an injective resolution $M \rightarrow I^{\bullet}$, and we define

$$
\left(R^{r} F\right)(M)=H^{r}\left(F\left(I^{\bullet}\right)\right) .
$$

A morphism $\alpha: M \rightarrow N$ gives rise to a well-defined morphism $\left(R^{r} F\right)(M) \rightarrow\left(R^{r} F\right)(N)$, and these maps make $R^{r} F$ into a functor. They are called the right derived functors of $F$.

The next two lemmas prove something a little more precise than the proposition.\\
Lemma A. 9 Let $M \rightarrow I^{\bullet}$ and $N \rightarrow J^{\bullet}$ be resolutions of objects $M$ and $N$ of C. If $N \rightarrow J^{\bullet}$ is an injective resolution, then every morphism $\alpha: M \rightarrow N$ extends to a morphism\\
\includegraphics[max width=\textwidth, alt={}, center]{e7cfc28e-5737-4710-b4d2-10a47512e540-38_169_228_1051_918}\\
of complexes.\\
Proof. Bucur and Deleanu 1968, 7.5. \(\square\)

Two morphisms $\alpha^{\bullet}, \beta^{\bullet}: I^{\bullet} \rightarrow J^{\bullet}$ of complexes are said to be homotopic if there exists a family of morphisms $k^{r}: I^{r} \rightarrow J^{r-1}$ (a homotopy) such that

$$
\alpha^{r}-\beta^{r}=d^{r-1} \circ k^{r}+k^{r+1} \circ d^{r}
$$

for all $r$.\\
Note that, for every $x \in Z^{r}\left(I^{\bullet}\right) \stackrel{\text { def }}{=} \operatorname{Ker}\left(d^{r}\right)$,

$$
\alpha^{r}(x)-\beta^{r}(x)=d^{r-1}\left(k^{r}(x)\right) \in \operatorname{Im}\left(d^{r-1}\right) \stackrel{\text { def }}{=} B^{r}\left(J^{\bullet}\right) .
$$

Therefore $\alpha^{r}(x)$ and $\beta^{r}(x)$ have the same image in $H^{r}\left(J^{\bullet}\right)$, and so homotopic morphisms define the same morphism on cohomology.

Lemma A. 10 Let $M \rightarrow I^{\bullet}$ be a resolution of $M$, and let $N \rightarrow J^{\bullet}$ be an injective resolution $N$. Any two extensions $\alpha^{\bullet}$ and $\beta^{\bullet}$ of morphisms $M \rightarrow N$ to $I^{\bullet} \rightarrow J^{\bullet}$ are homotopic.

Proof. Bucur and Deleanu 1968, 7.5. \(\square\)

This implies the second statement of Proposition A.8, because the family $\left(F\left(k^{r}\right)\right)$ is a homotopy from $F\left(\alpha^{\bullet}\right)$ to $F\left(\beta^{\bullet}\right)$.

Proposition A. 11 A short exact sequence

$$
0 \rightarrow A \rightarrow B \rightarrow C \rightarrow 0
$$

in C gives rise to a long exact sequence

$$
\begin{aligned}
& 0 \rightarrow F(A) \rightarrow F(B) \rightarrow F(C) \rightarrow R^{1} F(A) \rightarrow \cdots \\
& \quad \cdots \rightarrow R^{r} F(A) \rightarrow R^{r} F(B) \rightarrow R^{r} F(C) \rightarrow \cdots
\end{aligned}
$$

and the association of the long exact sequence to the short exact sequence is functorial.

The second condition means that a commutative diagram\\
\includegraphics[max width=\textwidth, alt={}, center]{e7cfc28e-5737-4710-b4d2-10a47512e540-39_170_671_423_697}\\
gives rise to a commutative diagram\\
\includegraphics[max width=\textwidth, alt={}, center]{e7cfc28e-5737-4710-b4d2-10a47512e540-39_194_1333_708_365}\\
For the proof of the proposition, see Bucur and Deleanu 1968, 7.6.\\
Remark A. 12 The right derived functors of $F$ are uniquely characterized (up to a unique isomorphism of functors) by the following three properties:

\begin{itemize}
  \item[(a)] $R^{0} F=F$;
  \item[(b)] $R^{r} F(I)=0$ for $r>0$ when $I$ is injective;
  \item[(c)] the property in (A.11).
\end{itemize}

Variants\\
By reversing the directions of some of the arrows, one obtains variants of some of the above definitions, for example, projective objects, left derived functors, etc.

The Ext groups\\
Let C be an abelian category.\\
Let $A \in \mathrm{C}$. If C has enough injectives, then we can define the right derived functors of the left exact functor $\operatorname{Hom}(A, \cdot)$. Denote the $r$ th right derived functor by $\operatorname{Ext}^{r}(A, \cdot)$. To compute $\operatorname{Ext}^{r}(A, B)$, we choose an injective resolution $B \rightarrow I^{\bullet}$ of $B$, and set

$$
\operatorname{Ext}^{r}(A, B)=H^{r}\left(\operatorname{Hom}\left(A, I^{\bullet}\right)\right) .
$$

Let $B \in \mathrm{C}$. If C has enough projectives, then we can define the right derived functors of the left exact contravariant functor $\operatorname{Hom}(\cdot, B)$. Denote the $r$ th right derived functor by $\operatorname{Ext}^{r}(\cdot, B)$. To compute $\operatorname{Ext}^{r}(A, B)$, we choose a projective resolution $P_{\bullet} \rightarrow A$ of $A$, and we set

$$
\operatorname{Ext}^{r}(A, B)=H^{r}\left(\operatorname{Hom}\left(P_{\bullet}, B\right)\right) .
$$

Proposition A. 13 If C has enough injectives and enough projectives, then the two definitions of $\operatorname{Ext}^{r}(A, B)$ coincide.

Proof. We define the $\operatorname{Ext}^{r}$ using projectives, and prove that they have the properties characterizing the right derived functors of $\operatorname{Hom}(A, \cdot)$.

First, certainly $\operatorname{Ext}^{0}(A, B)=\operatorname{Hom}(A, B)$.\\
To say that $I$ is injective means that $\operatorname{Hom}(\cdot, I)$ is exact. Therefore $\operatorname{Hom}\left(P_{\bullet}, I\right)$ is exact, and so

$$
\operatorname{Ext}^{r}(A, I) \stackrel{\text { def }}{=} H^{r}\left(\operatorname{Hom}\left(P_{\bullet}, I\right)\right)=0 .
$$

Finally, if

$$
0 \rightarrow B^{\prime} \rightarrow B \rightarrow B^{\prime \prime} \rightarrow 0
$$

is exact, then because $P_{\bullet} \rightarrow A$ is a resolution of $A$ by projectives, the sequence of complexes

$$
0 \rightarrow \operatorname{Hom}\left(P_{\bullet}, B^{\prime}\right) \rightarrow \operatorname{Hom}\left(P_{\bullet}, B\right) \rightarrow \operatorname{Hom}\left(P_{\bullet}, B^{\prime \prime}\right) \rightarrow 0
$$

is exact. By a standard procedure, we get out of this a long exact sequence

$$
\cdots \rightarrow H^{r}\left(\operatorname{Hom}\left(P_{\bullet}, B^{\prime}\right)\right) \rightarrow H^{r}\left(\operatorname{Hom}\left(P_{\bullet}, B\right)\right) \rightarrow H^{r}\left(\operatorname{Hom}\left(P_{\bullet}, B^{\prime \prime}\right)\right) \rightarrow \cdots .
$$ \(\square\)

Example A. 14 Let $G$ be a group. Then $\operatorname{Mod}_{G}$ has both enough injectives and enough projectives. For every $G$-module $M, \operatorname{Hom}_{G}(\mathbb{Z}, M)=M^{G}$, and so the functors $\operatorname{Hom}(\mathbb{Z}, \cdot)$ and $H^{0}(G, \cdot)$ agree. Hence, so also do their right derived functors:

$$
\operatorname{Ext}_{G}^{r}(\mathbb{Z}, M) \simeq H^{r}(G, M) .
$$

The last proposition allows us to compute these groups by choosing a projective resolution $P_{\bullet} \rightarrow \mathbb{Z}$ of $\mathbb{Z}$ and setting


\begin{equation*}
H^{r}(G, M)=H^{r}\left(\operatorname{Hom}_{G}\left(P_{\bullet}, M\right)\right) . \tag{25}
\end{equation*}


Remark A. 15 It would shorten the exposition in this chapter a little by adopting the formula (25) as the definition of $H^{r}(G, M)$. This is the approach taken by Atiyah and Wall (1967), but it is not the natural definition.

References\\
For the general notion of derived functors, see Chapter 7 of Bucur and Deleanu 1968.

\section*{Chapter III}
\section*{Local Class Field Theory: Cohomology}
In this chapter, we develop the cohomological approach to local class field theory. We first prove that there exists a local Artin map (Theorem 3.4). Together with the theorems of Chapter I, §§1-3, this suffices to complete the proofs of the main theorems of local class field theory. For those who skipped the Lubin-Tate theory, we then give a proof of the main theorems that is independent of Chapter I (except for §1).

Throughout this chapter, "local field" means "nonarchimedean local field". As before, $K^{\mathrm{al}}$ denotes an algebraic closure of $K$ (or separable algebraic closure in the case that $K$ has characteristic $p \neq 0$ ), and "extension of $K$ " means "subfield of $K^{\text {al }}$ containing $K$ ". All cohomology groups will computed using continuous cochains (see II, §4). For a Galois extension of fields $L / K$ (possibly infinite), set

$$
H^{2}(L / K)=H^{2}\left(\operatorname{Gal}(L / K), L^{\times}\right) .
$$

For the moment, $H^{2}(L / K)$ is just a cohomology group, but in the next chapter we shall see that it has an explicit interpretation as the relative Brauer group of $L / K$.

\section*{1 The Cohomology of Unramified Extensions}
The cohomology of the units\\
Let $K$ be a local field.\\
Proposition 1.1 Let $L / K$ be a finite unramified extension with Galois group $G$, and let $U_{L}$ be the group of units in $L$. Then

$$
H_{T}^{r}\left(G, U_{L}\right)=0, \quad \text { all } r .
$$

Proof. If $\pi$ is a prime element of $L$, then every element of $L^{\times}$can be written uniquely in the form $\alpha=u \pi^{m}, u \in U_{L}, m \in \mathbb{Z}$. Thus


\begin{equation*}
L^{\times}=U_{L} \cdot \pi^{\mathbb{Z}} \simeq U_{L} \times \mathbb{Z} . \tag{26}
\end{equation*}


Since $L$ is unramified over $K$, we can choose $\pi \in K$. Then $\tau \alpha=\tau\left(u \pi^{m}\right)=(\tau u) \pi^{m}$ for $\tau \in \operatorname{Gal}(L / K)$, and so (26) becomes a decomposition of $G$-modules when we let\\
$G$ act trivially on $\pi^{\mathbb{Z}} \simeq \mathbb{Z}$. Therefore, $H^{r}\left(G, U_{L}\right)$ is a direct summand of $H^{r}\left(G, L^{\times}\right)$ (see II, 1.25). Since $H^{1}\left(G, L^{\times}\right)=0$ by Hilbert's theorem 90 (II, 1.22), this shows that $H^{1}\left(G, U_{L}\right)=0$. Because $G$ is cyclic, to complete the proof of the theorem, it suffices (by II, 3.4) to show that $H_{T}^{0}\left(G, U_{L}\right)=0$. This is accomplished by the next proposition. \(\square\)

Proposition 1.2 Let $L / K$ be a finite unramified extension. Then the norm map $\operatorname{Nm}_{L / K}: U_{L} \rightarrow$ $U_{K}$ is surjective.

We first need some lemmas. Let $l$ and $k$ be the residue fields of $L$ and $K$. Because $L / K$ is unramified, the action of $G$ on $\mathcal{O}_{L}$ defines an isomorphism $G \simeq \operatorname{Gal}(l / k)$.

Lemma 1.3 For $m>0$, let $U_{L}^{(m)}=1+\mathfrak{m}_{L}^{m}$. Then

$$
\begin{aligned}
U_{L} / U_{L}^{(1)} & \xrightarrow{\simeq} l^{\times} \\
U_{L}^{(m)} / U_{L}^{(m+1)} & \approx l
\end{aligned}
$$

as $G$-modules.\\
Proof. Let $\pi$ be a prime element of $K$. It remains prime in $L$, and

$$
U_{L}^{(m)}=\left\{1+a \pi^{m} \mid a \in \mathcal{O}_{L}\right\} .
$$

The maps

$$
\begin{aligned}
u \mapsto u & \bmod \mathfrak{m}_{L}: U_{L} \rightarrow l^{\times} \\
1+a \pi^{m} \mapsto a & \bmod \mathfrak{m}_{L}: U_{L}^{(m)} \rightarrow l
\end{aligned}
$$

induce the required isomorphisms. \(\square\)

Lemma 1.4 For all $r, H_{T}^{r}\left(G, l^{\times}\right)=0$. In particular, the norm map $l^{\times} \rightarrow k^{\times}$is surjective.\\
Proof. By Hilbert's Theorem 90 (II, 1.22), $H^{1}\left(G, l^{\times}\right)=0$, and because $l^{\times}$is finite, its Herbrand quotient $h\left(l^{\times}\right)=1$ (see II 3.8). Therefore $H^{2}\left(G, l^{\times}\right)=0$, and this implies that all the groups are zero (by II, 3.4). \(\square\)

Lemma 1.5 The group $H_{T}^{r}(G, l)=0$ for all $r$. In particular, the trace map $l \rightarrow k$ is surjective.

Proof. In (II, 1.24) this is proved for $r>0$, which implies it for all $r$ (by II, 3.4). \(\square\)

Proof (of Proposition 1.2) There are commutative diagrams $(m>0)$\\
\includegraphics[max width=\textwidth, alt={}, center]{e7cfc28e-5737-4710-b4d2-10a47512e540-42_215_621_2038_726}\\
Consider $u \in U_{K}$. Because the norm map $l^{\times} \rightarrow k^{\times}$is surjective, there exists a $v_{0} \in$ $U_{L}$ such that $\operatorname{Nm}\left(v_{0}\right)$ and $u$ have the same image in $k^{\times}$, i.e., such that $u / \operatorname{Nm}\left(v_{0}\right) \in$ $U_{K}^{(1)}$. Because the trace map $l \rightarrow k$ is surjective, there exists a $v_{1} \in U_{L}^{(1)}$ such that $\operatorname{Nm}\left(v_{1}\right) \equiv u / \operatorname{Nm}\left(v_{0}\right) \bmod U_{K}^{(2)}$. Continuing in this fashion, we obtain a sequence of elements $v_{0}, v_{1}, v_{2}, v_{3}, \ldots, v_{i} \in U_{L}^{(i)}$, such that $u / \operatorname{Nm}\left(v_{0} \cdots v_{i}\right) \in U_{K}^{(i+1)}$. Let $v=$ $\lim _{m \rightarrow \infty} \prod_{j=0}^{m} v_{j}$. Then $u / \operatorname{Nm}(v) \in \bigcap U_{K}^{(i)}=\{1\}$. \(\square\)

Corollary 1.6 Let $L / K$ be an infinite unramified extension with Galois group $G$. Then $H^{r}\left(G, U_{L}\right)=0$ for $r>0$ (continuous cochains).

Proof. The field $L$ is a union of finite extensions $K^{\prime}$ of $K$, and so (see II, 4.2),

$$
H^{r}\left(\operatorname{Gal}(L / K), U_{L}\right) \simeq \lim _{\longrightarrow K^{\prime}} H^{r}\left(\operatorname{Gal}\left(K^{\prime} / K\right), U_{K^{\prime}}\right)=0 .
$$ \(\square\)

The invariant map\\
Let $L$ be an unramified extension of $K$ (possibly infinite), and let $G=\operatorname{Gal}(L / K)$.\\
As $H^{2}\left(G, U_{L}\right)=0=H^{3}\left(G, U_{L}\right)$, the cohomology sequence of the short exact sequence

$$
0 \rightarrow U_{L} \rightarrow L^{\times} \xrightarrow{\operatorname{ord}_{L}} \mathbb{Z} \rightarrow 0,
$$

gives an isomorphism

$$
H^{2}\left(G, L^{\times}\right) \xrightarrow[\simeq]{H^{2}\left(\operatorname{ord}_{L}\right)} H^{2}(G, \mathbb{Z}) .
$$

The groups $H^{r}(G, \mathbb{Q})$ are torsion for $r>0$ (see II, 4.3) and uniquely divisible (because $\mathbb{Q}$ is), and hence zero. Therefore the cohomology sequence of the short exact sequence

$$
0 \rightarrow \mathbb{Z} \rightarrow \mathbb{Q} \rightarrow \mathbb{Q} / \mathbb{Z} \rightarrow 0
$$

(trivial $G$-actions) gives an isomorphism

$$
H^{1}(G, \mathbb{Q} / \mathbb{Z}) \xrightarrow{\delta} H^{2}(G, \mathbb{Z}) .
$$

Recall that

$$
H^{1}(G, \mathbb{Q} / \mathbb{Z}) \simeq \operatorname{Hom}_{\text {cts }}(G, \mathbb{Q} / \mathbb{Z})
$$

and that $G$ has a canonical topological generator, namely, the Frobenius element $\sigma=$ $\operatorname{Frob}_{L / K}$. The composite of

$$
H^{2}(L / K) \xrightarrow[\simeq]{\text { ord }_{L}} H^{2}(G, \mathbb{Z}) \stackrel{\delta}{\simeq} H^{1}(G, \mathbb{Q} / \mathbb{Z}) \simeq \operatorname{Hom}_{\text {cts }}(G, \mathbb{Q} / \mathbb{Z}) \xrightarrow{f \mapsto f(\sigma)} \mathbb{Q} / \mathbb{Z}
$$

is called the invariant map

$$
\operatorname{inv}_{L / K}: H^{2}(L / K) \rightarrow \mathbb{Q} / \mathbb{Z} .
$$

Consider a tower of field extensions

$$
E \supset L \supset K
$$

with both $E$ and $L$ unramified (hence Galois) over $K$. Then\\
\includegraphics[max width=\textwidth, alt={}, center]{e7cfc28e-5737-4710-b4d2-10a47512e540-43_204_374_2307_848}\\
commutes, because all the maps in the definition of inv are compatible with Inf.

Theorem 1.7 There exists a unique isomorphism

$$
\operatorname{inv}_{K}: H^{2}\left(K^{\text {un }} / K\right) \rightarrow \mathbb{Q} / \mathbb{Z},
$$

with the property that, for every $L \subset K^{\text {un }}$ of finite degree $n$ over $K, \operatorname{inv}_{K}$ induces the isomorphism

$$
\operatorname{inv}_{L / K}: H^{2}(L / K) \rightarrow \frac{1}{[L: K]} \mathbb{Z} / \mathbb{Z} .
$$

Proof. This is an immediate consequence of the above discussion. \(\square\)

Proposition 1.8 Let $L$ be a finite extension of $K$ of degree $n$, and let $K^{\text {un }}$ and $L^{\text {un }}$ be the largest unramified extensions of $K$ and $L$. Then the following diagram commutes:\\
\includegraphics[max width=\textwidth, alt={}, center]{e7cfc28e-5737-4710-b4d2-10a47512e540-44_190_524_864_770}

(27)\\
Proof. The largest unramified extension of a local field is obtained by adjoining all $m$ th roots of 1 for $m$ not divisible by the residue characteristic. Therefore, $L^{\text {un }}=L \cdot K^{\text {un }}$, and so the map

$$
\tau \mapsto \tau \mid K^{\text {un }}: \operatorname{Gal}\left(L^{\text {un }} / L\right) \rightarrow \operatorname{Gal}\left(K^{\text {un }} / K\right)
$$

is injective. The map denoted Res in (27) is that defined by the compatible homomorphisms

$$
\begin{array}{ccc}
\operatorname{Gal}\left(K^{\text {un }} / K\right) & \leftarrow & \operatorname{Gal}\left(L^{\text {un }} / L\right) \\
K^{\text {un× }} & \rightarrow & L^{\text {un× }}
\end{array}
$$

(II, 1.27b).\\
Let $\Gamma_{K}=\operatorname{Gal}\left(K^{\text {un }} / K\right)$ and $\Gamma_{L}=\operatorname{Gal}\left(L^{\text {un }} / L\right)$, and consider the diagram,\\
\includegraphics[max width=\textwidth, alt={}, center]{e7cfc28e-5737-4710-b4d2-10a47512e540-44_242_1206_1636_432}

Here $e$ is the ramification index for $L / K$ and $f$ is the residue class degree. The first square is obtained from the commutative square\\
\includegraphics[max width=\textwidth, alt={}, center]{e7cfc28e-5737-4710-b4d2-10a47512e540-44_180_262_2029_905}\\
The second square expresses the fact that the restriction map commutes with the boundary map. Apart from the factor " $e$ ", the third square is\\
\includegraphics[max width=\textwidth, alt={}, center]{e7cfc28e-5737-4710-b4d2-10a47512e540-44_231_512_2365_776}

The Frobenius elements $\sigma_{K}$ and $\sigma_{L}$ are determined by the fact that they induce $x \mapsto x^{q}$ and $x \mapsto x^{q^{f}}$ respectively on the residue fields, where $q=|k|$ and $q^{f}=|l|$, and so $\sigma_{L} \mid K^{\text {un }}=\sigma_{K}^{f}$. It is now clear that the square commutes, and since $n=e f$, this proves the proposition. \(\square\)

\section*{The local Artin map}
Let $L$ be a finite unramified extension of $K$ with Galois group $G$, and let $n=[L: K]$. The local fundamental class $u_{L / K}$ is the element of $H^{2}(L / K)$ mapped to the generator $\frac{1}{[L: K]}$ of $\frac{1}{[L: K]} \mathbb{Z} / \mathbb{Z}$ by the invariant map $\operatorname{inv}_{L / K}: H^{2}(L / K) \xrightarrow{\simeq} \frac{1}{[L: K]} \mathbb{Z} / \mathbb{Z}$. The pair $\left(G, L^{\times}\right)$satisfies the hypotheses of Tate's theorem (II, 3.11), ${ }^{1}$ and so cup-product with the fundamental class $u_{L / K}$ defines an isomorphism

$$
H_{T}^{r}(G, \mathbb{Z}) \rightarrow H_{T}^{r+2}\left(G, L^{\times}\right)
$$

for all $r \in \mathbb{Z}$. For $r=-2$, this becomes\\
\includegraphics[max width=\textwidth, alt={}, center]{e7cfc28e-5737-4710-b4d2-10a47512e540-45_154_538_1137_764}\\
We now compute this map explicitly.\\
A prime element $\pi$ of $K$ is also a prime element of $L$, and defines a decomposition

$$
L^{\times}=U_{L} \cdot \pi^{\mathbb{Z}} \simeq U_{L} \times \mathbb{Z}
$$

of $G$-modules. Thus

$$
H^{r}\left(G, L^{\times}\right) \simeq H^{r}\left(G, U_{L}\right) \oplus H^{r}\left(G, \pi^{\mathbb{Z}}\right) .
$$

Choose a generator $\sigma$ of $G$ (e.g., the Frobenius generator), and let

$$
f \in H^{1}(G, \mathbb{Q} / \mathbb{Z}) \simeq \operatorname{Hom}(G, \mathbb{Q} / \mathbb{Z})
$$

be the element such that $f\left(\sigma^{i}\right)=\frac{i}{n} \bmod \mathbb{Z}$ for all $i$. It generates $H^{1}(G, \mathbb{Q} / \mathbb{Z})$.\\
From the exact sequence

$$
0 \rightarrow \mathbb{Z} \rightarrow \mathbb{Q} \rightarrow \mathbb{Q} / \mathbb{Z} \rightarrow 0
$$

and the fact that $H^{r}(G, \mathbb{Q})=0$ for all $r$, we obtain an isomorphism

$$
\delta: H^{1}(G, \mathbb{Q} / \mathbb{Z}) \rightarrow H^{2}(G, \mathbb{Z}) .
$$

According to the description of $\delta$ in (II, 1.21), to construct $\delta f$, we first choose a lifting of $f$ to 1 -cochain $\tilde{f}: G \rightarrow \mathbb{Q}$. We take $\tilde{f}$ to be the map $\sigma^{i} \mapsto \frac{i}{n}$, where $0 \leq i<n-1$. Then

$$
d \tilde{f}\left(\sigma^{i}, \sigma^{j}\right)=\sigma^{i} \tilde{f}\left(\sigma^{j}\right)-\tilde{f}\left(\sigma^{i+j}\right)+\tilde{f}\left(\sigma^{i}\right)=\left\{\begin{array}{lll}
0 & \text { if } & i+j \leq n-1 \\
1 & \text { if } & i+j>n-1 .
\end{array}\right.
$$

\footnotetext{${ }^{1}$ For the condition on $H^{1}$, use Hilbert's Theorem 90, and for the condition on $H^{2}$, use the invariant map.
}When we identify $\mathbb{Z}$ with the subgroup $\pi^{\mathbb{Z}}$ of $L^{\times}$, we find that the fundamental class $u_{L / K} \in H^{2}\left(G, L^{\times}\right)$is represented by the cocycle:

$$
\varphi\left(\sigma^{i}, \sigma^{j}\right)=\left\{\begin{array}{lll}
1 & \text { if } & i+j \leq n-1 \\
\pi & \text { if } & i+j>n-1
\end{array}\right.
$$

From the exact sequences

$$
\begin{gathered}
0 \rightarrow I \rightarrow \mathbb{Z}[G] \rightarrow \mathbb{Z} \rightarrow 0 \\
0 \rightarrow L^{\times} \rightarrow L^{\times}(\varphi) \rightarrow I \rightarrow 0
\end{gathered}
$$

(see the proof of II, 3.11) we obtain boundary maps

$$
\begin{aligned}
& H^{-2}(G, \mathbb{Z}) \rightarrow H^{-1}(G, I) \\
& H^{-1}(G, I) \rightarrow H^{0}\left(G, L^{\times}\right),
\end{aligned}
$$

which are isomorphisms because $\mathbb{Z}[G]$ and $L^{\times}(\varphi)$ have trivial cohomology. Here $L^{\times}(\varphi)$ is the splitting module $L^{\times} \oplus \bigoplus_{\sigma \in G, \sigma \neq 1} \mathbb{Z} x_{\sigma}$ of $\varphi$.

Finally, $H^{-2}(G, \mathbb{Z}) \stackrel{\text { def }}{=} H_{1}(G, \mathbb{Z}) \simeq G$ (see II, 2.7).

\begin{figure}[h]
\begin{center}
\captionsetup{labelformat=empty}
\caption{Proposition 1.9 Under the composite}
  \includegraphics[alt={},max width=\textwidth]{e7cfc28e-5737-4710-b4d2-10a47512e540-46_154_537_1360_764}
\end{center}
\end{figure}

of the above maps, the Frobenius element $\sigma \in G$ maps to the class of $\pi$ in $K^{\times} / \operatorname{Nm}\left(L^{\times}\right)$.\\
Note that, because all units in $K$ are norms from $L^{\times}$(see 1.2), the class of $\pi \bmod$ $\operatorname{Nm}\left(L^{\times}\right)$is independent of the prime element $\pi$. On the other hand, the $G$-module $L^{\times}(\varphi)$ and the map depend on the choice of the generator $\sigma$ for $G$.

Proof. From the construction of the isomorphism $H^{-2}(G, \mathbb{Z}) \simeq G$, we see that the image of $\sigma$ under the boundary map $H^{-2}(G, \mathbb{Z}) \rightarrow H^{-1}\left(G, I_{G}\right) \subset I_{G} / I_{G}^{2}$ is represented by $\sigma-1$.

The boundary map $H^{-1}\left(G, I_{G}\right) \rightarrow H^{0}\left(G, L^{\times}\right)$is that given by the snake lemma from the diagram (we write $I$ for $I_{G}$ ):\\
\includegraphics[max width=\textwidth, alt={}, center]{e7cfc28e-5737-4710-b4d2-10a47512e540-46_464_1042_2090_511}

The vertical maps connecting the rows are $\mathrm{Nm}_{G}=\sum_{i=0}^{n-1} \sigma^{i}$. The element $(\sigma-1)+I^{2}$ is the image of $x_{\sigma}+I \cdot L^{\times}(\varphi)$ in $L^{\times}(\varphi)_{G}$, and $\operatorname{Nm}_{G}\left(x_{\sigma}+I \cdot L^{\times}(\varphi)\right)$ is the sum of the elements:

$$
\begin{aligned}
x_{\sigma} & =x_{\sigma} \\
\sigma x_{\sigma} & =x_{\sigma^{2}}-x_{\sigma}+\varphi(\sigma, \sigma) \\
\sigma^{2} x_{\sigma} & =x_{\sigma^{3}}-x_{\sigma^{2}}+\varphi\left(\sigma, \sigma^{2}\right) \\
\cdots & \cdots \\
\sigma^{n-1} x_{\sigma} & ={ }^{\prime} x_{1},-x_{\sigma^{n-1}}+\varphi\left(\sigma, \sigma^{n-1}\right)
\end{aligned}
$$

On summing these, remembering that ' $x_{1}$ ' $=\varphi(1,1)=1$ and that + on the factor $L^{\times}$of $L(\varphi)$ is actually •, we find that

$$
\operatorname{Nm}_{G}\left(x_{\sigma}\right)=\prod_{i=1}^{n-1} \varphi\left(\sigma, \sigma^{i}\right)=\pi .
$$

This completes the proof. \(\square\)

Remark 1.10 The above proof of Proposition 1.9, using Tate's original definition of the isomorphism $H^{r}(G, \mathbb{Z}) \rightarrow H^{r+2}(G, C)$, is simpler than that found in other references, which uses the description of the map in terms of cup products.

\section*{2 The Cohomology of Ramified Extensions}
Because of Hilbert's Theorem 90 (II, 1.22), there is an exact sequence

$$
0 \rightarrow H^{2}(L / K) \xrightarrow{\operatorname{Inf}} H^{2}(E / K) \xrightarrow{\operatorname{Res}} H^{2}(E / L)
$$

for any tower of Galois extensions $E \supset L \supset K$ (see II, 1.36).\\
The next theorem extends Theorem 1.7 to ramified extensions.\\
Theorem 2.1 For every local field $K$, there exists a canonical isomorphism

$$
\operatorname{inv}_{K}: H^{2}\left(K^{\mathrm{al}} / K\right) \rightarrow \mathbb{Q} / \mathbb{Z} .
$$

Let $L$ be a Galois extension of $K$ of degree $n<\infty$. Then the diagram\\
\includegraphics[max width=\textwidth, alt={}, center]{e7cfc28e-5737-4710-b4d2-10a47512e540-47_204_945_1961_557}

(28)\\
commutes, and therefore defines an isomorphism

$$
\operatorname{inv}_{L / K}: H^{2}(L / K) \rightarrow \frac{1}{n} \mathbb{Z} / \mathbb{Z} .
$$

The proof will occupy most of the rest of the section.\\
Lemma 2.2 If $L / K$ is Galois of finite degree $n$, then $H^{2}(L / K)$ contains a subgroup canonically isomorphic to $\frac{1}{n} \mathbb{Z} / \mathbb{Z}$.

Proof. Consider the diagram\\
\includegraphics[max width=\textwidth, alt={}, center]{e7cfc28e-5737-4710-b4d2-10a47512e540-48_198_953_427_556}\\
Since the two inflation maps are injective, so also is the first vertical map, but 1.8 shows that the kernel of the restriction map on the top row is canonically isomorphic to $\frac{1}{n} \mathbb{Z} / \mathbb{Z}$. □

To complete the proof of Theorem 2.1, it suffices to prove that the map $\frac{1}{n} \mathbb{Z} / \mathbb{Z} \hookrightarrow$ $H^{2}(L / K)$ is an isomorphism. There are two different approaches to proving this. In the next chapter on Brauer groups, we shall show that $H^{2}\left(K^{\text {un }} / K\right) \simeq H^{2}\left(K^{\text {al }} / K\right)$ (every central simple algebra over $K$ is split by an unramified extension; see IV §4), and so the vertical maps in the above diagram are isomorphisms. The second proof, which we now present, shows that $\left|H^{2}(L / K)\right| \leq n$.

Lemma 2.3 Let $L$ be a finite Galois extension of $K$ with Galois group $G$. Then there exists an open subgroup $V$ of $\mathcal{O}_{L}$, stable under $G$, such that $H^{r}(G, V)=0$ for all $r>0$.

Proof. Let $\left\{x_{\tau} \mid \tau \in G\right\}$ be a normal basis for $L$ over $K$ (see FT, 5.18). The $x_{\tau}$ have a common denominator $d$ in $\mathcal{O}_{K}$ (see ANT, 2.6). After replacing each $x_{\tau}$ with $d \cdot x_{\tau}$, we may suppose that they lie in $\mathcal{O}_{L}$. Take $V=\sum \mathcal{O}_{K} x_{\tau}$. It is stable under $G$ because the normal basis is. Let $\pi$ be a prime element of $\mathcal{O}_{L}$. Then $V \supset \pi^{m} \mathcal{O}_{L}$ for some $m>0$, which shows that $V$ is open (it is union of cosets of $\pi^{m} \mathcal{O}_{L}$ ). Finally,

$$
V \simeq \mathcal{O}_{K}[G] \simeq \operatorname{Ind}^{G} \mathcal{O}_{K}
$$

as $G$-modules, and so $H^{r}(G, V)=0$ for all $r>0$ (II, 1.12). \(\square\)

Lemma 2.4 Let $L, K$, and $G$, be as in the last lemma. Then there exists an open subgroup $V$ of $U_{L}$ stable under $G$ such that $H^{r}(G, V)=0$ for all $r>0$.

Proof. I prove this only for $K$ of characteristic zero. ${ }^{2}$ The power series

$$
e^{x}=1+x+\cdots+x^{n} / n!+\cdots
$$

converges for $\operatorname{ord}(x)>\operatorname{ord}(p) /(p-1)$ (ANT, 7.29). It defines an isomorphism of an open neighbourhood of 0 in $L$ onto an open neighbourhood of 1 in $L^{\times}$, with inverse mapping

$$
\log (x)=(x-1)-(x-1)^{2} / 2+(x-1)^{3} / 3-\cdots .
$$

Both maps commute with the actions of $G$ (because $G$ acts continuously). If $V^{\prime}$ is an open neighbourhood of 0 as in (2.3), then $\pi^{M} V^{\prime}$ will have the same properties, and we can take $V=\exp \left(\pi^{M} V^{\prime}\right)$ with $M$ chosen to be sufficiently large that the exponential function is defined and an isomorphism on $\pi^{M} V^{\prime}$. \(\square\)

Lemma 2.5 Let $L / K$ be a cyclic extension of degree $n$; then $h\left(U_{L}\right)=1$ and $h\left(L^{\times}\right)=n$.

\footnotetext{${ }^{2}$ For a proof of the characteristic $p$ case, see p. 134 of Serre 1967b.
}Proof. Let $V$ be an open subgroup of $U_{L}$ with $H^{r}(G, V)=0$ for all $r$. Because $U_{L}$ is compact, the quotient $U_{L} / V$ is finite, and so $h\left(U_{L}\right)=h(V)=1$ by II, 3.9. Now $h\left(L^{\times}\right)=h\left(U_{L}\right) \cdot h(\mathbb{Z})=h(\mathbb{Z})$, and

$$
h(\mathbb{Z})=\left|H_{T}^{0}(G, \mathbb{Z})\right| /\left|H^{1}(G, \mathbb{Z})\right|=\left|H_{T}^{0}(G, \mathbb{Z})\right|=|\mathbb{Z} / n \mathbb{Z}|=n .
$$ \(\square\)

Lemma 2.6 Let $L$ be a finite Galois extension of $K$ of order $n$, then $H^{2}(L / K)$ has order $n$.

Proof. We know that the order of $H^{2}(L / K)$ is divisible by $n$, and that it equals $n$ when $L / K$ is cyclic. We prove the lemma by induction on $[L: K]$. Because the group $\operatorname{Gal}(L / K)$ is solvable (ANT, 7.59), there exists a Galois extension $K^{\prime} / K$ with $L \supsetneqq K^{\prime} \supsetneqq K$. From the exact sequence

$$
0 \rightarrow H^{2}\left(K^{\prime} / K\right) \rightarrow H^{2}(L / K) \rightarrow H^{2}\left(L / K^{\prime}\right)
$$

we see that

$$
\left|H^{2}(L / K)\right| \leq\left|H^{2}\left(K^{\prime} / K\right)\right| \cdot\left|H^{2}\left(L / K^{\prime}\right)\right|=n .
$$ \(\square\)

Proof (of Theorem 2.1) From the diagram in the proof of (2.2) we see that, for every finite Galois extension $L$ of $K$, the subgroup $H^{2}(L / K)$ of $H^{2}\left(K^{\mathrm{al}} / K\right)$ is contained in $H^{2}\left(K^{\text {un }} / K\right)$. Since $H^{2}\left(K^{\text {al }} / K\right)=\bigcup H^{2}(L / K)$, this proves that the inflation map $H^{2}\left(K^{\text {un }} / K\right) \rightarrow H^{2}\left(K^{\text {al }} / K\right)$ is an isomorphism. Compose the inverse of this with the invariant map $\operatorname{inv}_{K}: H^{2}\left(K^{\text {un }} / K\right) \rightarrow \mathbb{Q} / \mathbb{Z}$ of Theorem 1.7. \(\square\)

In fact, the diagram in Theorem 2.1 commutes even when $L / K$ is not Galois:


\begin{equation*}
\operatorname{inv}_{L} \circ \operatorname{Res}=[L: K] \operatorname{inv}_{K} . \tag{29}
\end{equation*}


(In this setting, we define $H^{2}(L / K)$ to be the kernel of the restriction map.)\\
I claim that


\begin{equation*}
\operatorname{inv}_{K} \circ \operatorname{Cor}=\operatorname{inv}_{L} . \tag{30}
\end{equation*}


Because Res is surjective (see diagram (28)), to prove this it suffices to show that

$$
\operatorname{inv}_{K} \circ \text { Cor } \circ \operatorname{Res}=\operatorname{inv}_{L} \circ \operatorname{Res} .
$$

But Cor $\circ$ Res $=[L: K]$ by (II, 1.30), and so this is the preceding equality (29).

The fundamental class\\
Let $L$ be a finite Galois extension of $K$ with Galois group $G$. As in the unramified case, we define the fundamental class $u_{L / K}$ of $L / K$ to be the element $u_{L / K}$ of $H^{2}(L / K)$ such that

$$
\operatorname{inv}_{L / K}\left(u_{L / K}\right)=\frac{1}{[L: K]} \bmod \mathbb{Z},
$$

or, equivalently, such that

$$
\operatorname{inv}_{K}\left(u_{L / K}\right)=\frac{1}{[L: K]} \bmod \mathbb{Z} .
$$

Lemma 2.7 Let $L \supset E \supset K$ with $L / K$ finite and Galois. Then


\begin{align*}
& \operatorname{Res}\left(u_{L / K}\right)=u_{L / E},  \tag{31}\\
& \operatorname{Cor}\left(u_{L / E}\right)=[E: K] u_{L / K} . \tag{32}
\end{align*}


Moreover, if $E / K$ is also Galois, then


\begin{equation*}
\operatorname{Inf}\left(u_{E / K}\right)=[L: E] u_{L / K} . \tag{33}
\end{equation*}


Proof. Consider\\
\includegraphics[max width=\textwidth, alt={}, center]{e7cfc28e-5737-4710-b4d2-10a47512e540-50_227_916_817_575}\\
On applying the kernel-cokernel lemma (II, A.2) to the rows (and using that the lemma is functorial), we obtain a commutative diagram,\\
\includegraphics[max width=\textwidth, alt={}, center]{e7cfc28e-5737-4710-b4d2-10a47512e540-50_245_1057_1217_507}\\
Here $H^{2}(E / K)$ is defined to be the kernel of $\operatorname{Res}: H^{2}\left(K^{\mathrm{al}} / K\right) \rightarrow H^{2}\left(K^{\mathrm{al}} / E\right)$. The equality (31) expresses the fact that the second square commutes. When $E / K$ is Galois, the equality (33) expresses the fact that the first square commutes.

Equality (32) follows from comparing

$$
\begin{aligned}
\operatorname{inv}_{K}\left(\operatorname{Cor}_{E / K}\left(u_{L / E}\right)\right) & \stackrel{(30)}{=} \operatorname{inv}_{E}\left(u_{L / E}\right)=\frac{1}{[L: E]} \bmod \mathbb{Z} \\
& \operatorname{inv}_{K}\left(u_{L / K}\right)=\frac{1}{[L: K]} \bmod \mathbb{Z} .
\end{aligned}
$$ \(\square\)

REMARK 2.8 There are groups $K^{\times} / \operatorname{Nm}\left(L^{\times}\right)$and $H^{2}(L / K)$ attached to every finite extension $L / K$ of local fields. When $L / K$ is cyclic, they are isomorphic, but not otherwise. The first group is always isomorphic to $\operatorname{Gal}(M / K)$, where $M$ is the largest abelian subextension of $L / K$ (see 3.4, 3.5), and the second is always cyclic of order $[L: K]$. Thus, when $L / K$ is abelian but not cyclic, the two groups have the same order but are not isomorphic, and when $L / K$ is not abelian, they have different orders.

\section*{3 The Local Artin Map}
The pair $\left(G, L^{\times}\right)$satisfies the hypotheses of Tate's theorem (II, 3.11),and so we have proved the following result.

\begin{figure}[h]
\begin{center}
\captionsetup{labelformat=empty}
\caption{Theorem 3.1 For every finite Galois extension of local fields $L / K$ and $r \in \mathbb{Z}$, the homomorphism}
$$
H_{T}^{r}(\operatorname{Gal}(L / K), \mathbb{Z}) \rightarrow H_{T}^{r+2}\left(\operatorname{Gal}(L / K), L^{\times}\right)
$$\end{center}
\end{figure}

defined by $x \mapsto x \cup u_{L / K}$ is an isomorphism. When $r=-2$, this becomes an isomorphism

$$
G^{\mathrm{ab}} \simeq K^{\times} / \mathrm{Nm}_{L / K}\left(L^{\times}\right) .
$$

If $L \supset E \supset K$ with $L / K$ Galois, then the homomorphism in the theorem commutes with both Res and Cor. For example, to prove the statement for Res we must show that

$$
\operatorname{Res}\left(x \cup u_{L / K}\right)=\operatorname{Res}(x) \cup u_{L / E}
$$

for all $x \in H_{T}^{r}\left(\operatorname{Gal}(L / K), L^{\times}\right)$. But a standard property of cup-products (II, 1.39c) is that

$$
\operatorname{Res}\left(x \cup u_{L / K}\right)=\operatorname{Res}(x) \cup \operatorname{Res}\left(u_{L / K}\right),
$$

and so this follows from (31). The proof for Cor is similar.\\
If $L \supset E \supset K$ with $L / K$ and $E / K$ Galois and $x \in H^{r}\left(\operatorname{Gal}(E / K), E^{\times}\right)$with $r \geq 1$, then

$$
\operatorname{Inf}\left(x \cup u_{E / K}\right)=[L: E] \operatorname{Inf}(x) \cup u_{L / K} .
$$

Again, this follows from a standard property of cup-products (1.39e) and the formula (33).\\
For $r=-2$, the map in the theorem becomes

$$
\operatorname{Gal}(L / K)^{\mathrm{ab}} \xrightarrow{\simeq} K^{\times} / \mathrm{Nm}_{L / K} L^{\times} .
$$

We denote the inverse map by

$$
\phi_{L / K}: K^{\times} / \mathrm{Nm}_{L / K} L^{\times} \rightarrow \operatorname{Gal}(L / K)^{\mathrm{ab}}
$$

and call it the local Artin map, or by $\operatorname{rec}_{L / K}$ and call it the local reciprocity map. (The second name is more common, but it is hard to see any reciprocity in the map.)

Lemma 3.2 Let $L \supset E \supset K$ be local fields with $L / K$ Galois. Then the following diagrams commute:\\
\includegraphics[max width=\textwidth, alt={}, center]{e7cfc28e-5737-4710-b4d2-10a47512e540-51_242_1144_1763_461}\\
The unmarked vertical arrows are induced by the inclusions $\operatorname{Gal}(L / E) \subset \operatorname{Gal}(L / K)$ and $K \subset E$.

Proof. To be added (the second is not really needed). \(\square\)\\
3.3 Let $L \supset E \supset K$ be local fields with both $L$ and $E$ Galois over $K$. Then the following diagram commutes:\\
\includegraphics[max width=\textwidth, alt={}, center]{e7cfc28e-5737-4710-b4d2-10a47512e540-51_207_391_2324_838}\\
The unmarked vertical arrow is induced by the surjection $\sigma \mapsto \sigma \mid E: \operatorname{Gal}(L / K) \rightarrow \operatorname{Gal}(E / K)$.

Proof. This follows directly from the definition of the local Artin map, using (33). ${ }^{3}$ \(\square\)

In particular, if $L \supset E \supset K$ is a tower of finite abelian extensions of $K$, then $\phi_{L / K}(a) \mid E=\phi_{E / K}(a)$ for all $a \in K^{\times}$, and so we can define $\varphi_{K}: K^{\times} \rightarrow \operatorname{Gal}\left(K^{\mathrm{ab}} / K\right)$ to be the homomorphism such that, for every finite abelian extension $L / K, \phi_{K}(a) \mid L=$ $\varphi_{L / K}(a)$.

Theorem 3.4 For every local field $K$, there exists a homomorphism (local Artin map)

$$
\phi_{K}: K^{\times} \rightarrow \operatorname{Gal}\left(K^{\mathrm{ab}} / K\right)
$$

with the following properties:

\begin{itemize}
  \item[(a)] for every prime element $\pi$ of $K, \phi_{K}(\pi) \mid K^{\mathrm{un}}=\operatorname{Frob}_{K}$;
  \item[(b)] for every finite abelian extension $L$ of $K, \operatorname{Nm}_{L / K}\left(L^{\times}\right)$is contained in the kernel of $a \mapsto \phi_{K}(a) \mid L$, and $\phi_{K}$ induces an isomorphism
$$
\phi_{L / K}: K^{\times} / \operatorname{Nm}_{L / K}\left(L^{\times}\right) \rightarrow \operatorname{Gal}(L / K) .
$$
\end{itemize}

Proof. Everything is obvious from the above, except (a), which follows from the fact that, for an unramified extension $L$ of $K, \phi_{L / K}$ agrees with that defined in §1, and so we can apply Proposition 1.9. \(\square\)

For those who have read §3 of Chapter I, this completes the proof of the main theorems of local class field theory - they can now skip to Chapter V if they wish. Others will need to continue to the end of the chapter.

Theorem 3.5 (Norm Limitation Theorem) Let $L$ be a finite extension of $K$, and let $E$ be the largest abelian extension of $K$ contained in $L$; then

$$
\operatorname{Nm}_{L / K}\left(L^{\times}\right)=\operatorname{Nm}_{E / K}\left(E^{\times}\right) .
$$

Proof. Because of the transitivity of the norm map, $\mathrm{Nm}_{L / K}\left(L^{\times}\right)$is a subgroup of $\mathrm{Nm}_{E / K}\left(E^{\times}\right)$. If $L / K$ is Galois, then $\operatorname{Gal}(E / K)=\operatorname{Gal}(L / K)^{\text {ab }}$, and so Theorem 3.1 shows that both norm groups have index $[E: K]$ in $K^{\times}$. This implies that they are equal.

In the general case, we let $L^{\prime}$ be a finite Galois extension of $K$ containing $L$. Let $G=\operatorname{Gal}\left(L^{\prime} / K\right)$ and $H=\operatorname{Gal}\left(L^{\prime} / L\right)$. Thus:\\
\includegraphics[max width=\textwidth, alt={}, center]{e7cfc28e-5737-4710-b4d2-10a47512e540-52_377_146_2015_959}

\footnotetext{${ }^{3}$ Hervé Jacquet has pointed out that this doesn't appear to be correct, but that the statement can be proved using Proposition 3.6. I don't know whether it can be proved without using 3.6.
}Then $E$ is the largest subfield of $L^{\prime}$ that is abelian over $K$ and contained in $L$. Therefore, the subgroup of $G$ fixing it is $G^{\prime} \cdot H$, where $G^{\prime}$ is the derived group of $G$. Let $a \in \operatorname{Nm}\left(E^{\times}\right)$. We have to show that $a \in \operatorname{Nm}\left(L^{\times}\right)$. Consider the commutative diagram\\
\includegraphics[max width=\textwidth, alt={}, center]{e7cfc28e-5737-4710-b4d2-10a47512e540-53_337_322_521_871}\\
The element $\phi_{L^{\prime} / K}(a)$ of $G / G^{\prime}$ maps to 1 in $G / G^{\prime} H$. As $\phi_{L^{\prime} / L}$ is surjective, we see from the diagram that there exists a $b \in L^{\times}$such $\phi_{L^{\prime} / K}(a)=\phi_{L^{\prime} / K}(\operatorname{Nm}(b))$, and hence $a / \operatorname{Nm}(b) \in \operatorname{Nm}\left(L^{\prime \times}\right)$, say, $a / \operatorname{Nm}(b)=\operatorname{Nm}(c)$. Now

$$
a=\operatorname{Nm}_{L / K}\left(b \cdot \operatorname{Nm}_{L^{\prime} / L}(c)\right) \in \operatorname{Nm}_{L / K}\left(L^{\times}\right) .
$$ \(\square\)

Theorem 3.5 shows that there is no hope of classifying nonabelian extensions of a local field in terms of the norm groups.

Alternative description of the local Artin map\\
Let $L / K$ be a finite abelian extension with Galois group $G$, and let $u_{L / K} \in H^{2}\left(G, L^{\times}\right)$be the fundamental class. The local Artin map $\phi_{L / K}$ is the inverse to the isomorphism

$$
x \mapsto x \cup u_{L / K}: H_{T}^{-2}(G, \mathbb{Z}) \rightarrow H_{T}^{0}\left(G, L^{\times}\right) .
$$

This definition is difficult to work with because cup-products involving both homology and cohomology groups have no very convenient description. Instead, we re-interprete the map purely in terms of cohomology groups. Consider the cup-product pairing

$$
H^{0}\left(G, L^{\times}\right) \times H^{2}(G, \mathbb{Z}) \longrightarrow H^{2}\left(G, L^{\times}\right) \xrightarrow{\operatorname{inv}_{L / K}} \mathbb{Q} / \mathbb{Z} .
$$

Given an element $a \in H^{0}\left(G, L^{\times}\right)=K^{\times}$and a class $c \in H^{2}(G, \mathbb{Z})$ represented by a cocycle $f: G \times G \rightarrow \mathbb{Z}$, the cup-product class $a \cup c$ is represented by the cocycle $(\sigma, \tau) \mapsto$ $a^{f(\sigma, \tau)}$. Recall also that we have an isomorphism

$$
\operatorname{Hom}(G, \mathbb{Q} / \mathbb{Z})=H^{1}(G, \mathbb{Q} / \mathbb{Z}) \xrightarrow{\delta} H^{2}(G, \mathbb{Z}) .
$$

Proposition 3.6 For every $\chi \in \operatorname{Hom}_{c t s}(G, \mathbb{Q} / \mathbb{Z})$ and $a \in K^{\times}$,

$$
\chi\left(\phi_{L / K}(a)\right)=\operatorname{inv}_{K}(a \cup \delta \chi) .
$$

Proof. See Serre 1962, "Annexe" to Chapter XI, and Serre 1967a, p. 140. \(\square\)

Using this, we can get another proof of Proposition 3.4.\\
Lemma 3.7 If $L / K$ is unramified, $\phi_{L / K}$ sends $a \in K^{\times}$to Frob $^{\text {ord }_{K}(a)}$.

Proof. Recall that $\operatorname{inv}_{K}$ is defined to be the composite

$$
H^{2}\left(G, L^{\times}\right) \xrightarrow{\text { ord }} H^{2}(G, \mathbb{Z}) \stackrel{\delta}{\leftarrow} \operatorname{Hom}(G, \mathbb{Q} / \mathbb{Z}) \xrightarrow{\chi \mapsto \chi(\sigma)} \mathbb{Q} / \mathbb{Z} .
$$

Because of the functoriality of cup-products

$$
\operatorname{ord}(a \cup \delta \chi)=\operatorname{ord}(a) \cup \delta \chi, \quad a \in K^{\times}, \quad \chi \in \operatorname{Hom}(G, \mathbb{Q} / \mathbb{Z})
$$

where, on the left ord denotes the map on $H^{2}$ induced by $\operatorname{ord}_{L}: L^{\times} \rightarrow \mathbb{Z}$, and on the right it is the map itself. Let $a \in H^{0}\left(G, L^{\times}\right)=K^{\times}$, and let $m=\operatorname{ord}_{L}(a)$. For every $\chi \in \operatorname{Hom}(G, \mathbb{Q} / \mathbb{Z})$, in the above diagram,

$$
a \cup \delta \chi \mapsto \operatorname{ord}(a) \cup \delta \chi \mapsto m \chi \mapsto \chi\left(\sigma^{m}\right), \quad \sigma=\text { Frob },
$$

i.e., $\operatorname{inv}_{K}(a \cup \delta \chi)=\chi\left(\sigma^{m}\right)$. On combining this with the formula in (3.6) we find that

$$
\chi(\phi(\alpha))=\chi\left(\sigma^{\operatorname{ord}(\alpha)}\right)
$$

for all $\chi \in \operatorname{Hom}(G, \mathbb{Q} / \mathbb{Z})$, and so $\phi(\alpha)=\sigma^{\text {ord }(\alpha)}$. \(\square\)

For every character $\chi$ of $G$, we get a character $a \mapsto \operatorname{inv}_{K}(a \cup \delta \chi)$ of $K^{\times}$. By duality we get a map $K^{\times} \rightarrow G$. Proposition 3.6 shows that this map is $\phi_{L / K}: K^{\times} \rightarrow \operatorname{Gal}(L / K)$.

\section*{4 The Hilbert symbol}
Since the subgroups of finite index in $K^{\times}$intersect in \{1\}, the existence theorem implies that the norm groups intersect in 1. The following key result is the starting point in our proof of the existence theorem.

Proposition 4.1 Let $K$ be a local field containing a primitive $n$th root of 1 . Any element of $K^{\times}$that is a norm from every cyclic extension of $K$ of degree dividing $n$ is an $n$th power.

This will be a consequence of the theory of the Hilbert symbol.

A very special case\\
We explain in this subsection how to obtain the proposition in the special case $K=\mathbb{Q}_{p}$, $n=2$, by using only elementary results from Serre 1970, Chapter III.

For $a, b \in \mathbb{Q}_{p}^{\times}$, define

$$
(a, b)_{p}=\left\{\begin{aligned}
1 & \text { if } z^{2}=a x^{2}+b y^{2} \text { has a nontrivial solution in } \mathbb{Q}_{p} \\
-1 & \text { otherwise }
\end{aligned}\right.
$$

Clearly, $(a, b)_{p}$ depends only $a$ and $b$ modulo squares, and so this is a pairing

$$
a, b \mapsto(a, b)_{p}: \mathbb{Q}_{p}^{\times} / \mathbb{Q}_{p}^{\times 2} \times \mathbb{Q}_{p}^{\times} / \mathbb{Q}_{p}^{\times 2} \rightarrow\{ \pm 1\} .
$$

Serre shows (ibid. 1.1, 1.2) that this pairing is bi-multiplicative

$$
\left(a a^{\prime}, b\right)_{p}=(a, b)_{p}\left(a^{\prime}, b\right)_{p}, \quad\left(a, b b^{\prime}\right)_{p}=(a, b)_{p}\left(a, b^{\prime}\right)_{p},
$$

symmetric (which in this case is the same as skew-symmetric)

$$
(b, a)_{p}=(a, b)_{p}^{-1}=(a, b)_{p},
$$

and nondegenerate (i.e., the left and right kernels are trivial).\\
Let $a$ be an element of $\mathbb{Q}_{p}^{\times}$that is not a square. Then

$$
\begin{aligned}
b \text { is a norm from } \mathbb{Q}_{p}[\sqrt{a}] & \Longleftrightarrow b=z^{2}-a x^{2} \text { has a solution in } \mathbb{Q}_{p} \\
& \Longleftrightarrow z^{2}=a x^{2}+b y^{2} \text { has a nontrivial solution in } \mathbb{Q}_{p} \\
& \Longleftrightarrow(a, b)_{p}=1
\end{aligned}
$$

Thus,

$$
b \text { is a norm from } \mathbb{Q}_{p}[\sqrt{a}] \text { for all } a \Longrightarrow(a, b)_{p}=1 \text { for all } a \Longrightarrow b \in \mathbb{Q}_{p}^{\times 2},
$$

by the nondegeneracy of the pairing. This proves the proposition in this case.\\
Aside 4.2 The map $a, b \mapsto(a, b)_{p}$ is called the Hilbert symbol. It is defined also for $\mathbb{R}$, and Serre proves (Hilbert product formula)


\begin{equation*}
\prod_{p \leq \infty}(a, b)_{p}=1 . \tag{34}
\end{equation*}


\section*{The general case}
In this subsection, $\operatorname{Hom}(G, \mathbb{Z} / n \mathbb{Z})$ denotes the group of continuous homomorphisms $G \rightarrow$ $\mathbb{Z} / n \mathbb{Z}$.

Step 1. Review.\\
Recall that for every cyclic group $\bar{G}$ and $\bar{G}$-module $M$, there are isomorphisms

$$
H_{T}^{r}(\bar{G}, M) \rightarrow H_{T}^{r+2}(\bar{G}, M),
$$

which become canonical once one chooses a generator $\sigma$ of $\bar{G}$. More precisely, let $n=(\bar{G}$ : 1) and let $\chi$ be the isomorphism $\bar{G} \rightarrow \mathbb{Z} / n \mathbb{Z}$ sending $\sigma$ to $1 \bmod n$. From the cohomology sequence of

$$
0 \rightarrow \mathbb{Z} \xrightarrow{n} \mathbb{Z} \rightarrow \mathbb{Z} / n \mathbb{Z} \rightarrow 0,
$$

we obtain a boundary map $\delta: H^{1}(\bar{G}, \mathbb{Z} / n \mathbb{Z}) \rightarrow H^{2}(\bar{G}, \mathbb{Z})$. The isomorphism

$$
H_{T}^{r}(\bar{G}, M) \rightarrow H_{T}^{r+2}(\bar{G}, M)
$$

is $b \mapsto \delta \chi \cup b$.\\
Now let $K$ be a field (not necessarily local), and let $L / K$ cyclic extension with Galois group $\bar{G}$ (of order $n$ ). When $r=0$ and $M=L^{\times}$, the isomorphism becomes

$$
b \mapsto \delta \chi \cup b: K^{\times} / \operatorname{Nm} L^{\times} \rightarrow H^{2}\left(\bar{G}, L^{\times}\right) .
$$

For $\chi \in \operatorname{Hom}(\bar{G}, \mathbb{Z} / n \mathbb{Z})$ and $b \in K^{\times}$, define

$$
(\chi, b)^{\prime}=\delta \chi \cup b \in H^{2}\left(\bar{G}, L^{\times}\right) .
$$

The above discussion shows that, if $\chi$ is an isomorphism $\bar{G} \rightarrow \mathbb{Z} / n \mathbb{Z}$, then

$$
(\chi, b)^{\prime}=0 \Longleftrightarrow b \text { is a norm from } L^{\times} .
$$

Step 2.\\
Let $K$ be a field (not necessarily local). If $K$ has characteristic $p \neq 0$, then we require $n$ to be relatively prime to $p$. Let $K^{\text {al }}$ be a separable algebraic closure of $K$ (thus, simply an algebraic closure if $K$ is perfect $)$, and let $G=\operatorname{Gal}\left(K^{\mathrm{al}} / K\right)$. Let $\mu_{n}$ be the group of $n$th roots of 1 in $K^{\text {al }}$, regarded as a $G$-module. Under our assumption on $n, \mu_{n}$ is a cyclic group of order $n$.

From the cohomology sequence of

$$
0 \longrightarrow \mu_{n} \longrightarrow K^{\mathrm{al} \times} \xrightarrow{x \mapsto x^{n}} K^{\mathrm{al} \times} \longrightarrow 0
$$

and Hilbert's theorem 90 (II, 1.22) we find that

$$
\delta: K^{\times} / K^{\times n} \xrightarrow{\simeq} H^{1}\left(G, \mu_{n}\right), \quad H^{2}\left(G, \mu_{n}\right) \xrightarrow{\simeq} H^{2}\left(G, K^{\text {al× }}\right)_{n} .
$$

Here $H^{2}\left(G, K^{\text {al× }}\right)_{n}$ denotes the group of elements in $H^{2}\left(G, K^{\text {al× }}\right)$ killed by $n$.\\
Consider the cup-product pairing

$$
\begin{array}{ccccc}
H^{1}(G, \mathbb{Z} / n \mathbb{Z}) & \times & H^{1}\left(G, \mu_{n}\right) & \rightarrow & H^{2}\left(G, \mu_{n}\right) \\
\text { । } & & \text { ।ʔ } & & \text { ।ʔ } \\
\operatorname{Hom}(G, \mathbb{Z} / n \mathbb{Z}) & & K^{\times} / K^{\times n} & & H^{2}\left(G, K^{\text {al× }}\right)_{n} .
\end{array}
$$

For $\chi \in \operatorname{Hom}(G, \mathbb{Z} / n \mathbb{Z})$ and $b \in K^{\times} / K^{\times n}$, we write $(\chi, b)$ for the image of the pair under the above pairing: $(\chi, b)=\chi \cup \delta b$ (as an element of $H^{2}\left(G, K^{\mathrm{al} \times}\right)$ ).

Let $\chi \in \operatorname{Hom}(G, \mathbb{Z} / n \mathbb{Z})$ and (for simplicity) assume that $\chi$ has order $n$. Let $L_{\chi}$ be the subfield of $K^{\mathrm{al}}$ fixed by $\operatorname{Ker}(\chi)$. Thus $L_{\chi}$ is a cyclic extension of $K$ of degree $n$, and $\chi$ induces an isomorphism of its Galois group $\bar{G}$ with $\mathbb{Z} / n \mathbb{Z}$. Let $b \in K^{\times}$, and let

$$
(\chi, b)^{\prime}=\delta \chi \cup b \in H^{2}\left(\bar{G}, L^{\times}\right),
$$

as in Step 1. Then

$$
\operatorname{Inf}\left((\chi, b)^{\prime}\right)=(\chi, b) .
$$

Since the inflation map $H^{2}\left(\bar{G}, L_{\chi}\right) \rightarrow H^{2}\left(G, K^{\text {al× }}\right)$ is injective (see II, 1.36), we deduce that

$$
(\chi, b)=0 \Longleftrightarrow b \text { is a norm from } L_{\chi}^{\times} .
$$

Step 3.\\
Now assume that $K$ is a local field, so that we have a canonical isomorphism

$$
\operatorname{inv}_{K}: H^{2}\left(G, K^{\mathrm{alx}}\right) \rightarrow \mathbb{Q} / \mathbb{Z} .
$$

For $\chi \in \operatorname{Hom}(G, \mathbb{Z} / n \mathbb{Z}) \simeq H^{1}(G, \mathbb{Z} / n \mathbb{Z})$ and $b \in K^{\times} / K^{\times n} \simeq H^{1}\left(G, \mu_{n}\right)$, we now define

$$
(\chi, b)=\operatorname{inv}_{K}(\chi \cup \delta b) \in \frac{1}{n} \mathbb{Z} / \mathbb{Z} .
$$

Thus, we have a canonical pairing

$$
\chi, b \mapsto(\chi, b): H^{1}(G, \mathbb{Z} / n \mathbb{Z}) \times H^{1}\left(G, \mu_{n}\right) \rightarrow H^{2}\left(G, \mu_{n}\right) \simeq \frac{1}{n} \mathbb{Z} / \mathbb{Z} .
$$

Proposition 4.3 The left kernel of this pairing is zero.

Proof. Let $\chi \in \operatorname{Hom}(G, \mathbb{Z} / n \mathbb{Z})$; we have to show that if $(\chi, b)=0$ for all $b \in K^{\times} / K^{\times n}$, then $\chi=0$.

Recall that the composite of the local Artin map $\phi_{K}: K^{\times} \rightarrow \operatorname{Gal}\left(K^{\mathrm{ab}} / K\right)$ with

$$
\tau \mapsto \tau \mid L: \operatorname{Gal}\left(K^{\mathrm{ab}} / K\right) \rightarrow \operatorname{Gal}(L / K)
$$

is surjective for every finite extension $L$ of $K$ contained in $K^{\mathrm{ab}}$. This implies that the image of $\phi_{K}$ is dense in $\operatorname{Gal}\left(K^{\mathrm{ab}} / K\right)$. Recall (3.6) that, for $b \in K^{\times}$,

$$
\left(\chi, b \bmod K^{\times n}\right)=\chi\left(\phi_{K}(b)\right) .
$$

Because $\chi$ is continuous, if it is zero on $\phi_{K}\left(K^{\times}\right)$, then it is zero on the whole of $\operatorname{Gal}\left(K^{\mathrm{ab}} / K\right) . \square$\\
Step 4.\\
Now assume that $K$ contains a primitive $n$th root of 1 . Thus, $\mu_{n} \approx \mathbb{Z} / n \mathbb{Z}$ as a $G$ -module (the isomorphism depends on the choice of a primitive $n$th root of 1 ). In particular, $\mu_{n} \otimes \mu_{n} \approx \mu_{n}$, and so $H^{2}\left(G, \mu_{n} \otimes \mu_{n}\right) \approx H^{2}\left(G, \mu_{n}\right)$ (noncanonically). There is a canonical isomorphism

$$
H^{2}\left(G, \mu_{n}\right) \otimes \mu_{n} \xrightarrow{\simeq} H^{2}\left(G, \mu_{n} \otimes \mu_{n}\right),
$$

which can be defined as $x, y \mapsto x \cup y$. On combining this with the isomorphism $H^{2}\left(G, \mu_{n}\right) \rightarrow$ $\frac{1}{n} \mathbb{Z} / \mathbb{Z}$ given by the invariant map, we obtain an isomorphism

$$
H^{2}\left(G, \mu_{n} \otimes \mu_{n}\right) \simeq \mu_{n} .
$$

Thus, we have a cup-product pairing

$$
\begin{array}{cccc}
H^{1}\left(G, \mu_{n}\right) & \times & H^{1}\left(G, \mu_{n}\right) & \rightarrow \\
\text { ।2 } & \text { ।2 } & H^{2}\left(G, \mu_{n} \otimes \mu_{n}\right) \\
K^{\times} / K^{\times n} & K^{\times} / K^{\times n} & \text { ।2 } \\
\mu_{n .}
\end{array}
$$

For $a, b \in K^{\times}$, their image in $\mu_{n}$ under this pairing is denoted $(a, b)$. The pairing is called the Hilbert symbol.

Theorem 4.4 The Hilbert symbol has the following properties.

\begin{itemize}
  \item[(a)] It is bi-multiplicative, i.e.,
$$
\left(a a^{\prime}, b\right)=(a, b)\left(a^{\prime}, b\right), \quad\left(a, b b^{\prime}\right)=(a, b)\left(a, b^{\prime}\right) .
$$
  \item[(b)] It is skew-symmetric, i.e.,
$$
(b, a)=(a, b)^{-1} .
$$
  \item[(c)] It is nondegenerate, i.e.,
$$
\begin{aligned}
& (a, b)=1 \text { for all } b \in K^{\times} / K^{\times n} \Longrightarrow a \in K^{\times n}, \\
& (a, b)=1 \text { for all } a \in K^{\times} / K^{\times n} \Longrightarrow b \in K^{\times n} .
\end{aligned}
$$
  \item[(d)] $(a, b)=1$ if and only if $b$ is a norm from $K[\sqrt[n]{a}]$.
\end{itemize}

Proof. Statements (a) and (b) are immediate from the definition of the Hilbert symbol in terms of cup-products.

We next prove (c). Choose an $n$th root of 1 , and use it to define isomorphisms $\mathbb{Z} / n \mathbb{Z} \rightarrow$ $\mu_{n}$ and $\mu_{n} \rightarrow \mu_{n} \otimes \mu_{n}$. The diagram defined by these maps,

$$
\begin{array}{ccccc}
H^{1}\left(G, \mu_{n}\right) & \times & H^{1}\left(G, \mu_{n}\right) & \rightarrow & H^{2}\left(G, \mu_{n} \otimes \mu_{n}\right) \\
\uparrow & & \| & & \uparrow \\
H^{1}(G, \mathbb{Z} / n \mathbb{Z}) & \times & H^{1}\left(G, \mu_{n}\right) & \rightarrow & H^{2}\left(G, \mu_{n}\right),
\end{array}
$$

commutes. By Proposition 4.3, the left kernel of the lower pairing is zero. Hence, the left kernel of the upper pairing is also zero, which, by skew-symmetry, implies that the right kernel is zero.

Finally, we prove (d). Assume (for simplicity), that $K[\sqrt[n]{a}]$ has degree $n$ over $K$. Choose a primitive $n$th root $\zeta$ of 1. Then $K[\sqrt[n]{a}]$ is cyclic of degree $n$ over $K$, and its Galois group is generated by the map $\sigma$ sending $\sqrt[n]{a}$ to $\zeta \sqrt[n]{a}$. With this choice of $\zeta$, $a \in K^{\times} / K^{\times n}$ corresponds to the element of $\chi_{a} \in \operatorname{Hom}(G, \mathbb{Z} / n \mathbb{Z})$ sending $\sigma$ to $1 \bmod$ $n \mathbb{Z}$. As $K[\sqrt[n]{a}]=L_{\chi_{a}}$, (d) follows from the boxed formula above. \(\square\)

Proof (of Proposition 4.1) We can deduce Proposition 4.1 from Theorem 4.4 by the same argument as in the case $K=\mathbb{Q}_{p}, n=2$, namely, if $b$ is a norm from $K[\sqrt[n]{a}]$ for all $a$, then $(a, b)=1$ for all $a$ (by 4.4d), and hence $b \in K^{\times n}$ (by 4.4c). \(\square\)

REMARK 4.5 The Hilbert symbol is related to the local Artin map by the formula

$$
\phi_{K}(b)\left(a^{\frac{1}{n}}\right)=(a, b) a^{\frac{1}{n}} .
$$

Note that Galois theory tells us that, for every $\tau \in \operatorname{Gal}\left(K\left[a^{\frac{1}{n}}\right] / K\right), \tau a^{\frac{1}{n}}=\zeta a^{\frac{1}{n}}$ for some $n$th root of one $\zeta$ (remember, we are assuming that $K$ contains the $n$th roots of 1 ), and so the point of the formula is that roots of 1 are the same. The proof of the formula is an exercise in cup-products, starting from Proposition 3.6.(a)

REMARK 4.6 The existence theorem will show that $\phi_{K}$ defines an isomorphism $K^{\times} / K^{\times n} \rightarrow$ $\operatorname{Gal}(L / K)$, where $L / K$ is the largest abelian extension of $K$ of exponent $n$ (without the existence theorem we know only that the map is surjective). Clearly, $\operatorname{Hom}(G, \mathbb{Z} / n \mathbb{Z})=$ $\operatorname{Hom}(\operatorname{Gal}(L / K), \mathbb{Z} / n \mathbb{Z})$, which is the dual of $\operatorname{Gal}(L / K)$. Thus, $K^{\times} / K^{\times n}$ and $\operatorname{Hom}(G, \mathbb{Z} / n \mathbb{Z})$ have the same order. This shows that the pairing in Step 2,

$$
H^{1}(G, \mathbb{Z} / n \mathbb{Z}) \times H^{1}\left(G, \mu_{n}\right) \rightarrow H^{2}\left(G, \mu_{n}\right) \simeq \frac{1}{n} \mathbb{Z} / \mathbb{Z}
$$

is nondegenerate. More generally, for every finite $G$-module $M$, the cup-product pairing

$$
H^{r}(G, M) \times H^{2-r}(G, \check{M}) \rightarrow H^{2}\left(G, \mu_{n}\right) \simeq \frac{1}{n} \mathbb{Z} / \mathbb{Z}
$$

is nondegenerate - here $\check{M}=\operatorname{Hom}\left(M, \mu_{n}\right)$. This duality theorem was proved by Poitou and Tate in the early 1950s, and was a starting point for the theorems in my book "Arithmetic Duality Theorems". Note that it implies that $H^{2}\left(G, \mu_{n} \otimes \mu_{n}\right) \simeq\left(\mu_{n}\right)_{G}$.

Remark 4.7 The Hilbert symbol has an interpretation in terms of central simple algebras (see the next chapter). For $a, b \in K^{\times}$and $\zeta$ a primitive $n$th root of 1 , define $A(a, b ; \zeta)$ to be the $K$-algebra with generators $i$ and $j$ satisfying the relations $i^{n}=a, j^{n}=b$, and $i j=\zeta j i$. Then, $(a, b)_{v}$ is the class in $\operatorname{Br}(K)_{n}$ represented by $A(a, b ; \zeta)$. Note that for $n=2$, we recover the quaternion algebra $H(a, b)$.

Remark 4.8 The Hilbert symbol

$$
(,): K^{\times} / K^{\times n} \times K^{\times} / K^{\times n} \rightarrow \mu_{n}
$$

is defined also for $K=\mathbb{R}$ and $n=2$. Moreover, the product formula (34) holds for all number fields and all $n$ (see V, 5.4 below).

\section*{5 The Existence Theorem}
Let $K$ be a local field. Recall that a subgroup $N$ of $K^{\times}$is a norm group if there is a finite abelian extension $L / K$ such that $\operatorname{Nm}_{L / K}\left(L^{\times}\right)=N$. Because $K^{\times} / N \xrightarrow{\simeq} \operatorname{Gal}(L / K)$, such a group $N$ is of finite index in $K^{\times}$and hence open (I, 1.3).

Theorem 5.1 (Existence Theorem) Every open subgroup of finite index in $K^{\times}$is a norm group.

In this draft, I prove this only for $K$ of characteristic zero. Then subgroups of finite index are open, and we can use the theory of the Hilbert symbol.

Remark 5.2 Before starting the proof of Theorem 5.1, we should consider the abstract situation: Let $Z$ be an infinite abelian group, and let $\mathcal{N}$ be a family of subgroups of finite index; what do we need to know in order to prove that $\mathcal{N}$ contains all subgroups of $Z$ of finite index? Clearly, a start is to know,

\begin{itemize}
  \item[(a)] if $I$ contains an element of $\mathcal{N}$, then $I$ is in $\mathcal{N}$;
  \item[(b)] if $N_{1}$ and $N_{2}$ are in $\mathcal{N}$, then $N_{1} \cap N_{2}$ contains an element of $\mathcal{N}$ (and hence is in $\mathcal{N}$ ).
\end{itemize}

We know that the family of norm subgroups of $K^{\times}$has these properties (Corollary 1.2), but so also, for example, does the set of subgroups of $\mathbb{Z}$ containing 5Z. Clearly, we need to know more. Let $D$ be the intersection of all the groups in $\mathcal{N}$. If every subgroup $I$ of finite index in $Z$ contains an element of $\mathcal{N}$, then it contains $D$. Conversely, if $D$ is divisible, i.e., $D=n D$ for all integers $n$, then $D$ is contained in every subgroup of finite index (a subgroup of $Z$ of index $n$ contains $n Z$, which contains $n D=D$ ). Therefore, we need that $D$ be divisible, but even this isn't sufficient. For example, let $Z=\mathbb{Z}$ and $\mathcal{N}=\left\{p^{n} \mathbb{Z} \mid n \in \mathbb{N}\right\}$. Then $D=\bigcap p^{n} \mathbb{Z}=0$ is divisible, but not every subgroup of finite index in $\mathbb{Z}$ contains an element of $\mathcal{N}$. [Probably I should include a complete list of "axioms" that a family $\mathcal{N}$ should satisfy in order to contain all open subgroups of finite index - see Serre 1962, XI.5, or Artin and Tate 1961, Chapter 14. This may be useful when we come to the global existence theorem.]

Proof of Theorem 5.1\\
Step 1. For all finite extensions $L / K$, the norm map $L^{\times} \rightarrow K^{\times}$has closed image and compact kernel.

Proof. Since the image has finite index, it is open (see I, 1.3), and hence closed. The kernel is closed, and the equality

$$
\operatorname{ord}_{L}(\operatorname{Nm}(a))=[L: K] \operatorname{ord}_{L}(a)=f \cdot \operatorname{ord}_{K}(a),
$$

shows that it contained in $U_{K}$, which is compact. \(\square\)

Let $D_{K}=\bigcap \mathrm{Nm}_{L / K}\left(L^{\times}\right)$, where $L$ runs over the finite extensions of $K$ (or only the finite abelian extensions - after Theorem 3.5, it is the same).

Step 2. For any finite extension $K^{\prime} / K, \operatorname{Nm}_{K^{\prime} / K} D_{K^{\prime}}=D_{K}$.\\
Proof. Let $a \in D_{K}$, and consider the sets

$$
\operatorname{Nm}_{L / K^{\prime}} L^{\times} \cap \operatorname{Nm}_{K^{\prime} / K}^{-1}(a)
$$

for $L / K^{\prime}$ finite. They are compact and nonempty, and any two contain a third. Therefore their intersection is nonempty (pointset topology). An element in the intersection lies in $D_{K^{\prime}}$ and has norm $a$. \(\square\)

Step 3. The group $D_{K}$ is divisible.\\
Proof. Let $n>1$ be an integer - we have to show that $n D_{K}=D_{K}$. Let $a \in D_{K}$. For each finite extension $L$ of $K$ containing a primitive $n$th root of 1, consider the set

$$
E(L)=\left\{b \in K^{\times} \mid b^{n}=a, \quad b \in \operatorname{Nm}_{L / K} L^{\times}\right\} .
$$

It is nonempty: $a=\operatorname{Nm}_{L / K} a^{\prime}$ for some $a^{\prime} \in D_{L}$; according to Proposition 4.1, $a^{\prime}=c^{n}$ for some $c \in L$, and

$$
\operatorname{Nm}_{L / K}(c)^{n}=\operatorname{Nm}_{L / K}\left(a^{\prime}\right)=a ;
$$

therefore $b \stackrel{\text { def }}{=} \operatorname{Nm}_{L / K}(c) \in E(L)$. Moreover, $E\left(L \cdot L^{\prime}\right) \subset E(L) \cap E\left(L^{\prime}\right)$. Since each set is finite, their intersection is nonempty. An element in their intersection lies in $D_{K}$ and has $n$th power $a$. \(\square\)

Step 4. $D_{K}=\{1\}$. [Actually, we don't use this here.]\\
Proof. Choose a prime element of $K$. Let $V_{m, n}=U^{(m)} \times \pi^{n \mathbb{Z}}$. Then $V_{m, n}$ is an (open) subgroup of finite index in $K^{\times}$, and therefore contains $D_{K}$. Now $D_{K} \subset \bigcap_{m, n} V_{m, n}=$ $\{1\}$. \(\square\)

Step 5. Every subgroup $I$ of finite index in $K^{\times}$that contains $U_{K}$ is a norm subgroup.\\
Proof. Since $\operatorname{ord}_{K}: K^{\times} \rightarrow \mathbb{Z}$ is surjective with kernel $U_{K}$, the subgroups $I$ in question are the subgroups of $K^{\times}$of the form $\operatorname{ord}_{K}^{-1}(n \mathbb{Z}), n \geq 1$. Let $K_{n}$ be the unramified extension of $K$ of degree $n$. Then $\operatorname{Nm}_{K_{n} / K}\left(K_{n}^{\times}\right)$is a subgroup of $K^{\times}$containing $U_{K}$ (see 1.2) with image $n \mathbb{Z}$ in $\mathbb{Z}$. The statement follows. \(\square\)

We now prove the theorem. Let $\mathcal{N}$ be the set of norm groups in $K^{\times}$, so that $D_{K}=$ $\bigcap_{N \in \mathcal{N}} N$. Let $I$ be a subgroup of $K^{\times}$of finite index. Because $D_{K}$ is divisible, $I \supset D_{K}$, and so $I \supset \bigcap_{N \in \mathcal{N}} N \supset \bigcap_{N \in \mathcal{N}}\left(N \cap U_{K}\right)$. Therefore the sets $\left(N \cap U_{K}\right) \backslash I$ have empty intersection. As they are compact some finite subfamily has empty intersection, i.e., for\\
some finite subset $\mathcal{S}$ of $\mathcal{N}, I \supset \bigcap_{N \in \mathcal{S}}\left(N \cap U_{K}\right)$. As any two of the sets $N \cap U_{K}$ contains a third, this implies that $I \supset N \cap U_{K}$ for some $N$.

Fix a norm group $N$ such that $N \cap U_{K} \subset I$. I claim that $I$ contains

$$
N \cap\left(U_{K} \cdot(N \cap I)\right) .
$$

Indeed, every element of the intersection is of the form $a b, a \in U_{K}, b \in N \cap I, a b \in N$. The last two statements imply that $a \in N$. Hence $a \in N \cap U_{K} \subset I$, and so $a b \in I$.

Now $N \cap I$ is of finite index in $K^{\times}$because both $N$ and $I$ are and $K^{\times} / N \cap I$ injects into $\left(K^{\times} / N\right) \times\left(K^{\times} / I\right)$. Hence $U_{K} \cdot(N \cap I)$ is a subgroup of finite index in $K^{\times}$containing $U_{K}$, and so is a norm group (Step 5). Now $N \cap\left(U_{K} \cdot(N \cap I)\right)$, being the intersection of two norm groups, contains a norm group. Therefore, $I$ contains a norm group, which implies that it is a norm group.

Notes It follows from Krasner's lemma (ANT, 7.60) that every finite abelian extension of local fields arises by completing a finite abelian extension of global fields. In the 1930s Hasse and F.K. Schmidt ${ }^{4}$ were able to deduce the main theorems of local class field theory from those of global class field theory.

From the modern perspective, this seems a strange way to do things. In the 1940s, in his algebraic approach to class field theory Chevalley developed local class field theory directly, and used it in the construction of global class field theory. F. K. Schmidt also showed that local class field theory can be constructed independently of global class field theory.

At that time, there was no good description of the local Artin map, and nor was there an explicit way of constructing the maximal abelian extension of a local field (except for $\mathbb{Q}_{p}$ of course).

In 1958 Dwork gave an explicit description of the local Artin map, which is reproduced in Serre 1962, but it was not very pleasant.

In 1965 Lubin and Tate introduced the Lubin-Tate formal group laws, and gave an explicit construction of $K^{\mathrm{ab}}$ and an explicit description of the local Artin map. However, they made use of the existence of the local Artin map (our Theorem I, 1.1) in their proofs.

In the early 1980s Lubin, Gold, and Rosen independently gave "elementary" proofs that $K^{\mathrm{ab}}=$ $K_{\pi} \cdot K^{\text {un }}$. In his book (Iwasawa 1986), Iwasawa develops the whole of local class field theory from the Lubin-Tate perspective, and also gives explicit formulas (due to de Shalit and Wiles) for the Hilbert symbols etc..

Other noncohomological approaches can be found in Hazewinkel 1975, Neukirch 1986, and Fesenko and Vostokov 1993. The disadvantage of the noncohomological approaches is, naturally, that they provide no information about the cohomology groups of local fields, which have important applications to other topics, for example elliptic curves.

For an explicit description of the local Artin map in the case of tame extensions, see Newton, Rachel, Explicit local reciprocity for tame extensions. Math. Proc. Cambridge Philos. Soc. 152 (2012), no. 3, 425-454.

In this chapter, I have largely followed Serre 1962 and Serre 1967a.

\footnotetext{${ }^{4}$ H. Hasse, Die Normenresttheorie relative-Abelscher Zahlkörper als Klassenkörper im Kleinen, J. für Mathematik (Crelle) 162 (1930), 145-154.\\
F. K. Schmidt, Zur Klassenkörpertheorie im Kleinen, ibid. 155-168.
}
\end{document}